% Découverte d’un environnement d’IA — Informatique 3ème — Cours signé OnBuch+
% Programme officiel MINESEC. Compiler avec : tectonic N01-decouverte-environnement-ia.tex
\newcommand{\DOCMATIERE}{Informatique}
\newcommand{\DOCNIVEAU}{3ème}
\newcommand{\DOCMODULE}{Informatique – classe de 3ème}
\newcommand{\DOCLECON}{N01}
\newcommand{\DOCTITRE}{Découverte d’un environnement d’IA}
% OnBuch+ — préambule « cours signés » ENRICHI (compilé avec tectonic / XeLaTeX)
% Système visuel « Pop » d'OnBuch+ : fond crème (jamais blanc), bordures encre
% épaisses, ombres « dures » décalées, titres Archivo Black en capitales.
% Version MATHÉMATIQUES : graphes (pgfplots/tikz), boîtes pédagogiques
% (définition, propriété, méthode, attention, le savais-tu, exercice résolu,
% exercice, TP, bilan), page de garde et signature OnBuch+ ancrée en bas.
%
% Macros attendues dans main.tex AVANT \input{preamble} :
%   \DOCMATIERE  \DOCNIVEAU  \DOCMODULE  \DOCLECON (numéro)  \DOCTITRE
\documentclass[11pt]{article}

\usepackage{fontspec}
\usepackage[french]{babel}
\usepackage[a4paper,margin=2cm,top=3.4cm,bottom=2.8cm,headheight=1.4cm,footskip=1.3cm]{geometry}
\usepackage{amsmath,amssymb,amsfonts}
\usepackage{mathtools} % \xmapsto, \coloneqq... souvent utilisés par les modèles
\usepackage{cancel} % simplifications barrées (bilans de forces)
% \abs{z} (module, valeur absolue) et \norm{u} (norme), utilisés par les modèles
\providecommand{\abs}{}\let\abs\relax\DeclarePairedDelimiter{\abs}{\lvert}{\rvert}
\providecommand{\norm}{}\let\norm\relax\DeclarePairedDelimiter{\norm}{\lVert}{\rVert}
\usepackage[table]{xcolor} % [table] : fournit aussi \rowcolors (alternance de lignes)
\usepackage{pagecolor}
\usepackage{tikz}
\usetikzlibrary{arrows.meta,positioning,calc,shapes.geometric,shapes.misc,shapes.arrows,shapes.symbols,decorations.pathmorphing,decorations.pathreplacing,decorations.markings,patterns,shadows,mindmap,trees,fit,backgrounds,matrix,intersections,angles,quotes}
\usepackage{pgfplots}
\usepackage[version=4]{mhchem} % équations nucléaires \ce{^{235}_{92}U}
\usepackage[european,siunitx]{circuitikz} % schémas électriques
\pgfplotsset{compat=1.18}
\usepgfplotslibrary{fillbetween,groupplots} % aires entre courbes (calcul intégral)
\usepackage{enumitem}
\usepackage{fancyhdr}
% [most] charge toutes les bibliothèques tcolorbox (listings, minted, raster,
% documentation...) et gonfle inutilement la mémoire TeX sur un document long
% (~300 boîtes) : on ne charge que ce qui est réellement utilisé.
\usepackage[skins,breakable]{tcolorbox}
\usepackage{graphicx}
\usepackage{colortbl}
\usepackage{booktabs}
\usepackage{tabularx}
\usepackage{diagbox} % case d'en-tête diagonale des tableaux à double entrée
% Colonnes extensibles centrée / gauche / droite (C, L, R), souvent utilisées par les modèles
\newcolumntype{C}{>{\centering\arraybackslash}X}
\newcolumntype{L}{>{\raggedright\arraybackslash}X}
\newcolumntype{R}{>{\raggedleft\arraybackslash}X}
\usepackage{array}
\usepackage{multirow}
\usepackage{multicol}
\usepackage{siunitx}
% parse-numbers=false : accepte aussi \SI{2,5 \times 10^{-3}}{...}
\sisetup{locale=FR,output-decimal-marker={,},group-minimum-digits=4,parse-numbers=false}
\DeclareSIUnit\ohm{\ensuremath{\symup{\Omega}}} % U+2126 absent de la police
\DeclareSIUnit\division{div} % réglages d'oscilloscope (ms/div, V/div)
\DeclareSIUnit\tr{tr} % tours (vitesse de rotation, ex. \SI{1500}{\tr\per\minute})
\DeclareSIUnit\bit{bit}
\DeclareSIUnit\byte{o} % octet
\DeclareSIUnit\inch{po} % pouce
\DeclareSIUnit\ppp{ppp} % points par pouce (résolution d'image)
\usepackage{qrcode}
% Blocs de code source (PHP, C, HTML, JavaScript, SQL) — spécifique à la
% matière Informatique : listings + habillage aux couleurs OnBuch+ (voir le
% style \lstset ci-dessous, à côté des autres boîtes pédagogiques).
\usepackage{listings}
% \degree : commande fréquemment utilisée par les modèles pour un angle ou une
% température en dehors de \SI{}{} (non fournie par les paquets chargés).
\providecommand{\degree}{^{\circ}}
% \qed (fin de démonstration), utilisé par les modèles sans amsthm
\providecommand{\qed}{\hfill$\square$}
% \barre : double barre des valeurs interdites dans un tableau de variations (tkz-tab)
\providecommand{\barre}{\ensuremath{\|}}
% environnement proof (amsthm non chargé) utilisé par les modèles
\ifdefined\proof\else\newenvironment{proof}[1][Démonstration]{\par\noindent\textit{#1.}\ }{\qed\par}\fi
\usepackage{lastpage}
\usepackage{hyperref}
\hypersetup{hidelinks}

% ============================================================
% Palette « Pop » OnBuch+ — mêmes hex que la classe P (plus_ui.dart)
% ============================================================
\definecolor{poporange}{HTML}{F59321}
\definecolor{popdark}{HTML}{E07A0C}
\definecolor{popcream}{HTML}{FFF6E8}
\definecolor{popink}{HTML}{1C1714}
\definecolor{poppurple}{HTML}{7A5AE0}
\definecolor{popgreen}{HTML}{2E9E6A}
\definecolor{poppink}{HTML}{E0457B}
\definecolor{popblue}{HTML}{2F7DE1}
\definecolor{popgold}{HTML}{C98A12}
\definecolor{popmuted}{HTML}{8A7F76}
\definecolor{popline}{HTML}{EADFD2}
\definecolor{popgray}{HTML}{8A7F76}
\colorlet{popgrey}{popgray}
% Alias tolérants (le modèle nomme parfois les couleurs différemment)
\colorlet{popwhite}{white}
\colorlet{popblack}{popink}
\colorlet{popred}{poppink}
\colorlet{popyellow}{popgold}
% Teintes claires pour fonds de tableaux / zones
\colorlet{popblueL}{popblue!10!white}
\colorlet{poporangeL}{poporange!14!white}
\colorlet{popgreenL}{popgreen!12!white}
\colorlet{poppurpleL}{poppurple!12!white}
\colorlet{poppinkL}{poppink!12!white}
\colorlet{popgoldL}{popgold!14!white}

\pagecolor{popcream}

% ============================================================
% Typographie
% ============================================================
\defaultfontfeatures{Ligatures=TeX}
\setmainfont{PlusJakartaSans-Medium.ttf}[Path=../../../fonts/,BoldFont=PlusJakartaSans-Bold.ttf]
% Maths en sans-serif (Fira Math) avec chiffres et lettres droites de Plus
% Jakarta, pour que les formules se fondent dans le texte.
\usepackage[math-style=ISO,bold-style=ISO]{unicode-math}
\setmathfont{FiraMath-Regular.otf}[Path=../../../fonts/]
\setmathfont{PlusJakartaSans-Medium.ttf}[Path=../../../fonts/,range={up/num,up/{Latin,latin},\mathup}]
\usepackage{icomma} % virgule décimale sans espace en mode math
% Caractères mathématiques Unicode absents des polices de texte (Plus Jakarta,
% Archivo Black) : rendus par leur équivalent mathématique, en texte comme en
% formule — sinon ils s'affichent en carré vide (ex. « Arithmétique dans ℤ »).
\usepackage{newunicodechar}
\newunicodechar{ℝ}{\ensuremath{\mathbb{R}}}
\newunicodechar{ℤ}{\ensuremath{\mathbb{Z}}}
\newunicodechar{ℂ}{\ensuremath{\mathbb{C}}}
\newunicodechar{ℕ}{\ensuremath{\mathbb{N}}}
\newunicodechar{ℚ}{\ensuremath{\mathbb{Q}}}
\newunicodechar{𝔻}{\ensuremath{\mathbb{D}}}
\newunicodechar{α}{\ensuremath{\alpha}}
\newunicodechar{β}{\ensuremath{\beta}}
\newunicodechar{γ}{\ensuremath{\gamma}}
\newunicodechar{δ}{\ensuremath{\delta}}
\newunicodechar{ε}{\ensuremath{\varepsilon}}
\newunicodechar{θ}{\ensuremath{\theta}}
\newunicodechar{λ}{\ensuremath{\lambda}}
\newunicodechar{μ}{\ensuremath{\mu}}
\newunicodechar{σ}{\ensuremath{\sigma}}
\newunicodechar{τ}{\ensuremath{\tau}}
\newunicodechar{φ}{\ensuremath{\varphi}}
\newunicodechar{ψ}{\ensuremath{\psi}}
\newunicodechar{ω}{\ensuremath{\omega}}
\newunicodechar{Σ}{\ensuremath{\Sigma}}
% majuscules produites par \MakeUppercase dans les titres (α-aminés -> Α)
\newunicodechar{Α}{\ensuremath{\alpha}}
\newunicodechar{Β}{\ensuremath{\beta}}
\newunicodechar{Γ}{\ensuremath{\Gamma}}
\newunicodechar{Δ}{\ensuremath{\Delta}}
\newunicodechar{Ε}{\ensuremath{\varepsilon}}
\newunicodechar{Θ}{\ensuremath{\Theta}}
\newunicodechar{Λ}{\ensuremath{\Lambda}}
\newunicodechar{Μ}{\ensuremath{\mu}}
\newunicodechar{Τ}{\ensuremath{\tau}}
\newunicodechar{Φ}{\ensuremath{\Phi}}
\newunicodechar{Ψ}{\ensuremath{\Psi}}
\newunicodechar{Ω}{\ensuremath{\Omega}}
\newunicodechar{↦}{\ensuremath{\mapsto}}
\newunicodechar{∈}{\ensuremath{\in}}
\newunicodechar{∉}{\ensuremath{\notin}}
\newunicodechar{∧}{\ensuremath{\wedge}}
\newunicodechar{⇔}{\ensuremath{\Leftrightarrow}}
\newunicodechar{⟺}{\ensuremath{\Longleftrightarrow}}
\newunicodechar{⇒}{\ensuremath{\Rightarrow}}
\newunicodechar{⟹}{\ensuremath{\Longrightarrow}}
\newunicodechar{∘}{\ensuremath{\circ}}
\newunicodechar{∪}{\ensuremath{\cup}}
\newunicodechar{∩}{\ensuremath{\cap}}
\newunicodechar{≡}{\ensuremath{\equiv}}
\newunicodechar{⊂}{\ensuremath{\subset}}
\newunicodechar{⊥}{\ensuremath{\perp}}
\newunicodechar{∃}{\ensuremath{\exists}}
\newunicodechar{∀}{\ensuremath{\forall}}
\newunicodechar{∛}{\ensuremath{\sqrt[3]{\,}}}
\newunicodechar{∜}{\ensuremath{\sqrt[4]{\,}}}
\newunicodechar{⃗}{\ensuremath{{}^{\to}}}
\newunicodechar{ⁿ}{\ifmmode^{n}\else\textsuperscript{\ensuremath{n}}\fi}
\newunicodechar{ˣ}{\ifmmode^{x}\else\textsuperscript{\ensuremath{x}}\fi}
\newunicodechar{ᵇ}{\ifmmode^{b}\else\textsuperscript{\ensuremath{b}}\fi}
\newunicodechar{ʳ}{\ifmmode^{r}\else\textsuperscript{\ensuremath{r}}\fi}
\newunicodechar{ᵘ}{\ifmmode^{u}\else\textsuperscript{\ensuremath{u}}\fi}
\newunicodechar{ʸ}{\ifmmode^{y}\else\textsuperscript{\ensuremath{y}}\fi}
\newunicodechar{ᵃ}{\ifmmode^{a}\else\textsuperscript{\ensuremath{a}}\fi}
\newunicodechar{ᵉ}{\ifmmode^{e}\else\textsuperscript{\ensuremath{e}}\fi}
\newunicodechar{ᵏ}{\ifmmode^{k}\else\textsuperscript{\ensuremath{k}}\fi}
\newunicodechar{ᵖ}{\ifmmode^{p}\else\textsuperscript{\ensuremath{p}}\fi}
\newunicodechar{ᴺ}{\ifmmode^{N}\else\textsuperscript{\ensuremath{N}}\fi}
\newunicodechar{ᶜ}{\ifmmode^{c}\else\textsuperscript{\ensuremath{c}}\fi}
\newunicodechar{ʰ}{\ifmmode^{h}\else\textsuperscript{\ensuremath{h}}\fi}
\newunicodechar{ᵠ}{\ifmmode^{q}\else\textsuperscript{\ensuremath{q}}\fi}
\newunicodechar{ₙ}{\ifmmode_{n}\else\textsubscript{\ensuremath{n}}\fi}
\newunicodechar{ₐ}{\ifmmode_{a}\else\textsubscript{\ensuremath{a}}\fi}
\newunicodechar{ᵢ}{\ifmmode_{i}\else\textsubscript{\ensuremath{i}}\fi}
\newunicodechar{ᵣ}{\ifmmode_{r}\else\textsubscript{\ensuremath{r}}\fi}
\newunicodechar{ₖ}{\ifmmode_{k}\else\textsubscript{\ensuremath{k}}\fi}
\newunicodechar{ᵦ}{\ifmmode_{\beta}\else\textsubscript{\ensuremath{\beta}}\fi}
\newunicodechar{ₓ}{\ifmmode_{x}\else\textsubscript{\ensuremath{x}}\fi}
\newfontfamily\popdisplay{ArchivoBlack-Regular.ttf}[Path=../../../fonts/]
\newfontfamily\poplogo{SpaceGrotesk-Bold.ttf}[Path=../../../fonts/]
\newfontfamily\popsemi{PlusJakartaSans-SemiBold.ttf}[Path=../../../fonts/]
\newfontfamily\popxbold{PlusJakartaSans-ExtraBold.ttf}[Path=../../../fonts/]
\newfontfamily\popxboldls{PlusJakartaSans-ExtraBold.ttf}[Path=../../../fonts/,LetterSpace=3.0]

\setlist[enumerate]{leftmargin=1.7em,itemsep=5pt,topsep=4pt}
\setlist[itemize]{leftmargin=1.7em,itemsep=4pt,topsep=4pt,label={\color{poporange}\textbullet}}
\setlength{\parindent}{0pt}
\setlength{\parskip}{5pt}
\renewcommand{\arraystretch}{1.25}

% pgfplots : style Pop par défaut
\pgfplotsset{
  every axis/.append style={axis line style={popink,line width=1pt},tick style={popink},
    grid style={popline,line width=0.5pt},label style={font=\small},tick label style={font=\footnotesize}},
  popcurve/.style={poporange,line width=1.8pt,no markers,smooth},
}
% Même style disponible dans un \draw TikZ simple (hors pgfplots)
\tikzset{popcurve/.style={poporange,line width=1.8pt}}
\tikzset{popgrid/.style={popline,very thin}}
\colorlet{popcurve}{poporange} % le nom du style est parfois utilisé comme couleur

% ============================================================
% Pill / squiggle
% ============================================================
\newcommand{\poppill}[3]{%
  \tikz[baseline=-0.55ex]\node[draw=popink,line width=1.2pt,rounded corners=2.6pt,%
    fill=#2,inner xsep=6.5pt,inner ysep=3pt]{%
    \popxboldls\fontsize{7}{7}\selectfont\color{#3}\MakeUppercase{#1}};%
}
\newcommand{\popsquiggle}[1]{%
  \tikz[baseline=-0.4ex]{
    \draw[#1,line width=2.6pt,line cap=round]
      plot [smooth] coordinates {(0,0.09) (0.34,0.01) (0.68,0.21) (1.02,0.01) (1.36,0.10)};
  }%
}
% Mot-clé surligné (marqueur orange sous le texte)
\newcommand{\cle}[1]{{\popxbold\color{popdark}#1}}

% ============================================================
% En-tête / pied
% ============================================================
\pagestyle{fancy}
\fancyhf{}
\renewcommand{\headrulewidth}{0pt}
\renewcommand{\footrulewidth}{0pt}
\lhead{\raisebox{-0.15ex}{\poplogo\fontsize{13}{13}\selectfont\color{popink}OnBuch\color{poporange}+}}
\rhead{\poppill{\DOCMATIERE\ \textbullet\ \DOCNIVEAU\ \textbullet\ Leçon \DOCLECON}{white}{popink}}
\lfoot{\popxbold\fontsize{7.6}{7.6}\selectfont\color{popmuted}\MakeUppercase{Cours signé OnBuch+}\ \textbullet\ {\popsemi\DOCTITRE}}
\rfoot{\tikz[baseline=-0.6ex]\node[rounded corners=8.5pt,fill=popink,minimum height=17pt,minimum width=17pt,inner xsep=5pt,inner sep=0]{\popxbold\fontsize{8}{8}\selectfont\color{popcream}\thepage\,/\,\pageref*{LastPage}};}

% ============================================================
% Titres
% ============================================================
\newcommand{\chaptitre}[2]{%
  \begin{tcolorbox}[enhanced,colback=poporange,colframe=popink,boxrule=2.6pt,arc=11pt,
      drop shadow=popink,left=15pt,right=15pt,top=12pt,bottom=11pt,before skip=3pt,after skip=18pt]
    \poppill{#2}{popink}{popcream}\par\vspace{9pt}
    {\raggedright\hyphenpenalty=10000\exhyphenpenalty=10000\popdisplay\fontsize{22}{24}\selectfont\color{popcream}\MakeUppercase{#1}\par}
  \end{tcolorbox}
}
% Section numérotée : pastille orange avec le numéro + titre Archivo
\newcounter{popsec}
\newcommand{\coursec}[1]{%
  \stepcounter{popsec}\setcounter{popsub}{0}%
  \par\vspace{16pt}\noindent
  \begin{minipage}{\linewidth}
  \tikz[baseline=-0.7ex]\node[draw=popink,line width=1.6pt,fill=poporange,rounded corners=4pt,minimum size=22pt,inner sep=0]{\popdisplay\fontsize{12}{12}\selectfont\color{popcream}\thepopsec};%
  \hspace{8pt}{\popdisplay\fontsize{15}{17}\selectfont\color{popink}\MakeUppercase{#1}}\par
  \vspace{4pt}\noindent{\color{poporange}\rule{36pt}{3.6pt}}
  \end{minipage}\par\nopagebreak\vspace{8pt}
}
\newcounter{popsub}
\newcommand{\courssub}[1]{%
  \stepcounter{popsub}%
  \par\vspace{9pt}\noindent{\popxbold\fontsize{11.5}{13}\selectfont\color{popink}{\color{poporange}\thepopsec.\thepopsub}\hspace{6pt}#1}\par\nopagebreak\vspace{4pt}
}

% ============================================================
% Boîtes pédagogiques — même recette : fond blanc, bordure épaisse colorée,
% ombre dure encre, badge plein. Titre optionnel [#1].
% ============================================================
\tcbset{popbase/.style={breakable,enhanced,colback=white,boxrule=2.3pt,arc=9pt,
  shadow={3pt}{-3pt}{0pt}{popink},left=14pt,right=14pt,top=11pt,bottom=11pt,before skip=11pt,after skip=15pt,
  fontupper=\popsemi\fontsize{10.6}{14.5}\selectfont\color{popink},
  fontlower=\popsemi\fontsize{10.6}{14.5}\selectfont\color{popink}}}
% Coche dessinée (le glyphe ✓ n'existe pas dans les polices du document)
\newcommand{\popcheck}[1]{\tikz[baseline=-0.5ex]\draw[#1,line width=1.6pt,line cap=round,line join=round](-0.13,0.0)--(-0.04,-0.09)--(0.13,0.1);}
\providecommand{\coche}{\popcheck{popgreen}}
\providecommand{\resultat}[1]{\par\smallskip\noindent\fcolorbox{popgreen}{popgreenL}{\parbox{\dimexpr\linewidth-2\fboxsep-2\fboxrule\relax}{\textbf{Résultat.} #1}}\par\smallskip}
\newcommand{\popboxhead}[3]{\poppill{#1}{#2}{white}\ifx\relax#3\relax\else\hspace{6pt}{\popxbold\fontsize{10.6}{12}\selectfont\color{popink}#3}\fi\par\vspace{7pt}}

\newtcolorbox{aretenir}[1][]{popbase,colframe=popgreen,before upper={\popboxhead{À retenir}{popgreen}{#1}}}
\newtcolorbox{exemplebox}[1][]{popbase,colframe=poporange,before upper={\popboxhead{Exemple}{poporange}{#1}}}
\newtcolorbox{definition}[1][]{popbase,colframe=popblue,before upper={\popboxhead{Définition}{popblue}{#1}}}
\newtcolorbox{propriete}[1][]{popbase,colframe=poppurple,before upper={\popboxhead{Propriété}{poppurple}{#1}}}
\newtcolorbox{methode}[1][]{popbase,colframe=popgold,before upper={\popboxhead{Méthode}{popgold}{#1}}}
\newtcolorbox{attention}[1][]{popbase,colframe=poppink,before upper={\popboxhead{Attention}{poppink}{#1}}}
\newtcolorbox{experience}[1][]{popbase,colframe=popblue,colback=popblueL,before upper={\popboxhead{Expérience}{popblue}{#1}}}
% « Le savais-tu ? » : carte violette pleine (contexte, culture, Cameroun)
\newtcolorbox{savaistu}[1][]{popbase,colback=poppurple,colframe=popink,
  fontupper=\popsemi\fontsize{10.4}{14.2}\selectfont\color{white},
  before upper={\poppill{Le savais-tu\,?}{white}{poppurple}\ifx\relax#1\relax\else\hspace{6pt}{\popxbold\color{white}#1}\fi\par\vspace{7pt}}}
% Exercice résolu : énoncé en haut, \tcblower puis la solution sur fond crème
\newcounter{popexo}\newcounter{popexr}
\newtcolorbox{exoresolu}[1][]{popbase,colframe=poporange,colbacklower=popcream,
  segmentation style={popink,line width=1.2pt,dashed},
  before upper={\stepcounter{popexr}\popboxhead{Exercice résolu \thepopexr}{poporange}{#1}},
  before lower={\poppill{Solution}{popink}{popcream}\par\vspace{6pt}}}
% Exercice (non résolu en place) — la correction est dans la partie Corrigés
\newtcolorbox{exercice}[1][]{popbase,colframe=popink,
  before upper={\stepcounter{popexo}\popboxhead{Exercice \thepopexo}{popink}{#1}}}
% Corrigé
\newtcolorbox{corrige}[1][]{popbase,colframe=popgreen,colback=popgreenL,
  before upper={\popboxhead{Corrigé}{popgreen}{#1}}}
% Objectifs / prérequis / bilan
\newtcolorbox{objectifs}{popbase,colframe=popink,colback=poporangeL,
  before upper={\popboxhead{Objectifs de la leçon}{popink}{}}}
\newtcolorbox{prerequis}{popbase,colframe=popink,colback=popblueL,
  before upper={\popboxhead{Prérequis}{popblue}{}}}
\newtcolorbox{bilan}{popbase,colframe=popink,colback=white,boxrule=2.8pt,
  before upper={\popboxhead{Fiche bilan}{poporange}{L'essentiel à connaître pour l'examen}}}
% Figure encadrée avec légende
\newtcolorbox{popfigure}[1][]{enhanced,breakable=false,colback=white,colframe=popline,boxrule=1.4pt,arc=8pt,
  left=10pt,right=10pt,top=10pt,bottom=8pt,before skip=10pt,after skip=12pt,halign=center,
  lower separated=false,halign lower=center,
  fontlower=\popsemi\fontsize{9}{11.5}\selectfont\color{popmuted},
  #1}
\newcommand{\legende}[1]{\par\vspace{6pt}{\popsemi\fontsize{9}{11.5}\selectfont\color{popmuted}#1}}

% ============================================================
% Bloc de code source — \begin{lstlisting}[language=Php]...\end{lstlisting}
% (ou language=C, HTML, JavaScript, SQL ; option omise = pseudo-code brut).
% Même recette visuelle que les figures (\popfigure) : cadre encre épais,
% fond clair (popcream), légende possible juste après via \legende{...}
% comme pour une figure (listings ne sait pas arrondir ses coins : on garde
% un cadre rectangulaire, cohérent avec les autres tableaux du document).
% CONTRAT : les modèles ne doivent utiliser QUE \begin{lstlisting}[...] tel
% quel (pas de \begin{tcblisting}, pas de \verb pour un bloc multi-lignes).
% ============================================================
\lstdefinestyle{poplst}{
  backgroundcolor=\color{popcream},
  basicstyle=\popsemi\fontsize{8.6}{11}\selectfont\color{popink},
  keywordstyle=\bfseries\color{popblue},
  commentstyle=\itshape\color{popmuted},
  stringstyle=\color{popgreen},
  numberstyle=\tiny\color{popmuted},
  identifierstyle=\color{popink},
  frame=l,
  rulecolor=\color{popink},
  framerule=1.6pt,
  framesep=7pt,
  framexleftmargin=7pt,
  framexrightmargin=7pt,
  xleftmargin=2pt,
  xrightmargin=2pt,
  breaklines=true,
  breakatwhitespace=false,
  showstringspaces=false,
  showspaces=false,
  showtabs=false,
  tabsize=2,
  columns=flexible,
  keepspaces=true,
  numbers=none,
  aboveskip=10pt,
  belowskip=10pt,
  captionpos=b,
}
\lstset{style=poplst, language=PHP}
% La distribution TeX embarquée ne fournit pas le fichier de langue
% JavaScript de listings (lstlang*.dat) : on la redéfinit nous-mêmes
% pour que \begin{lstlisting}[language=JavaScript] fonctionne.
\lstdefinelanguage{JavaScript}{
  keywords={break,case,catch,class,const,continue,debugger,default,delete,do,
    else,export,extends,finally,for,function,if,import,in,instanceof,let,new,
    return,static,super,switch,this,throw,try,typeof,var,void,while,with,yield,
    async,await,of,get,set},
  keywordstyle=\bfseries\color{popblue},
  ndkeywords={true,false,null,undefined,NaN,Infinity},
  ndkeywordstyle=\bfseries\color{poporange},
  identifierstyle=\color{popink},
  sensitive=true,
  comment=[l]{//},
  morecomment=[s]{/*}{*/},
  string=[b]",
  morestring=[b]',
  morestring=[b]`,
}
% Même limitation pour CSS et les fichiers de configuration (apacheconf,
% nginx, ini...) : ils ne sont pas fournis. CSS et un style générique de
% type "config" (commentaires # ou ;, pas de mots-clés) couvrent les besoins.
\lstdefinelanguage{CSS}{
  keywords={color,background,background-color,margin,padding,border,
    border-radius,font,font-size,font-weight,font-family,width,height,
    display,position,top,left,right,bottom,float,clear,text-align,
    line-height,list-style,cursor,overflow,box-shadow,transition,
    flex,grid,align-items,justify-content,z-index},
  keywordstyle=\bfseries\color{popblue},
  identifierstyle=\color{popink},
  sensitive=false,
  comment=[s]{/*}{*/},
  string=[b]",
  morestring=[b]',
}
\lstdefinelanguage{apacheconf}{
  sensitive=false,
  morecomment=[l]{\#},
}
\lstdefinelanguage{Apache}{
  sensitive=false,
  morecomment=[l]{\#},
}
\lstdefinelanguage{text}{sensitive=false}
\lstdefinelanguage{powershell}{
  keywords={New-Object,New-SmbShare,Get-Acl,Set-Acl,Get-ChildItem,Set-Content,
    Get-Content,Write-Host,ForEach-Object,Where-Object,Select-Object,
    Test-Path,Remove-Item,Copy-Item,Move-Item,Start-Process,Stop-Process,
    If,Else,ElseIf,ForEach,While,Function,Return,Try,Catch},
  keywordstyle=\bfseries\color{popblue},
  identifierstyle=\color{popink},
  sensitive=false,
  comment=[l]{\#},
  string=[b]",
  morestring=[b]',
}
\lstdefinelanguage{ini}{
  sensitive=false,
  morecomment=[l]{;},
  morecomment=[l]{\#},
}

% ============================================================
% Signature OnBuch+
% ============================================================
\newcommand{\poplogoinline}[1]{{\poplogo\fontsize{#1}{#1}\selectfont\color{popink}OnBuch\color{poporange}+}}
\newcommand{\signatureonbuch}{%
  \begin{tcolorbox}[enhanced,colback=popink,colframe=popink,boxrule=2.2pt,arc=10pt,
      drop shadow=poporange,left=14pt,right=14pt,top=10pt,bottom=10pt,before skip=0pt,after skip=0pt]
    \tikz[baseline=-0.7ex]\node[circle,fill=popgreen,draw=popcream,line width=1.3pt,minimum size=22pt,inner sep=0]{\popcheck{white}};%
    \hspace{9pt}\begin{minipage}[c]{0.62\linewidth}
      {\popxbold\fontsize{12}{13}\selectfont\color{popcream}Signé\ \textbullet\ {\poplogo OnBuch\color{poporange}+}}\par\vspace{2pt}
      {\raggedright\popsemi\fontsize{8}{10.5}\selectfont\color{popline}Ludovic A.\\\MakeUppercase{Conforme au programme officiel MINESEC}\par}
    \end{minipage}\hfill
    {\poplogo\fontsize{22}{22}\selectfont\color{popcream}OnBuch\color{poporange}+}
  \end{tcolorbox}%
}

% ============================================================
% Page de garde
% #1 titre · #2 matière · #3 niveau · #4 module · #5 n° leçon · #6 durée ·
% #7 description / objectifs (texte ou itemize)
% ============================================================
\newcommand{\pagegarde}[7]{%
  \thispagestyle{empty}
  \newgeometry{a4paper,margin=1.7cm,top=1.6cm,bottom=1.4cm}%
  \begin{tikzpicture}[remember picture,overlay]
    \fill[poporange!18] ([xshift=-3.2cm,yshift=-3.6cm]current page.north east) circle (4.6cm);
    \fill[poppurple!14] ([xshift=2.4cm,yshift=7.6cm]current page.south west) circle (3.8cm);
  \end{tikzpicture}%
  {\poplogo\fontsize{30}{30}\selectfont\color{popink}OnBuch\color{poporange}+}\hfill
  \raisebox{0.6ex}{\poppill{Cours signé \textbullet\ Premium}{popink}{popcream}}\par
  \vspace{1pt}\popsquiggle{poppurple}\par
  \vspace{8pt}
  \begin{tcolorbox}[enhanced,colback=poporange,colframe=popink,boxrule=2.8pt,arc=13pt,
      drop shadow=popink,left=18pt,right=18pt,top=12pt,bottom=13pt,after skip=10pt]
    \poppill{#2 \textbullet\ #3}{popink}{popcream}\hspace{5pt}\poppill{#4}{white}{popink}\par\vspace{8pt}
    {\popdisplay\fontsize{14}{14}\selectfont\color{popink}LEÇON #5}\par\vspace{4pt}
    {\raggedright\hyphenpenalty=10000\exhyphenpenalty=10000\popdisplay\fontsize{25}{27}\selectfont\color{popcream}\MakeUppercase{#1}\par}
    \vspace{8pt}
    {\popxbold\fontsize{9.5}{9.5}\selectfont\color{popink}\MakeUppercase{Durée conseillée : #6}}
  \end{tcolorbox}
  \begin{tcolorbox}[enhanced,colback=white,colframe=popink,boxrule=2.2pt,arc=10pt,
      drop shadow=popink,left=16pt,right=16pt,top=10pt,bottom=10pt,after skip=8pt]
    \poppill{Dans cette leçon}{poppurple}{white}\par\vspace{5pt}
    \popsemi\fontsize{9.3}{12.6}\selectfont\color{popink}#7
  \end{tcolorbox}
  \noindent\begin{minipage}[t]{0.487\linewidth}
    \begin{tcolorbox}[enhanced,colback=popgreen,colframe=popink,boxrule=2.2pt,arc=10pt,
        drop shadow=popink,left=11pt,right=11pt,top=10pt,bottom=10pt,height=3.1cm,valign=center]
      \tcbox[colback=white,colframe=popink,boxrule=1.3pt,arc=5pt,boxsep=0pt,nobeforeafter,
        left=3pt,right=3pt,top=3pt,bottom=3pt,valign=center]{\qrcode[height=1.9cm]{https://play.google.com/store/apps/details?id=com.luvvix.onbuch&hl=fr}}%
      \hspace{8pt}\begin{minipage}[c]{0.5\linewidth}
        {\popdisplay\fontsize{11}{12.5}\selectfont\color{white}\MakeUppercase{Télécharge OnBuch}}\par\vspace{4pt}
        {\raggedright\popsemi\fontsize{8.4}{11}\selectfont\color{white}Gratuit \textbullet\ Google Play\\Tuteur IA, annales, concours, résultats\par}
      \end{minipage}
    \end{tcolorbox}
  \end{minipage}\hfill
  \begin{minipage}[t]{0.487\linewidth}
    \begin{tcolorbox}[enhanced,colback=poppurple,colframe=popink,boxrule=2.2pt,arc=10pt,
        drop shadow=popink,left=11pt,right=11pt,top=10pt,bottom=10pt,height=3.1cm,valign=center]
      \tikz[baseline=-0.7ex]\node[circle,fill=white,draw=popink,line width=1.3pt,minimum size=1.6cm,inner sep=0]{\popdisplay\fontsize{24}{24}\selectfont\color{poppurple}+};%
      \hspace{8pt}\begin{minipage}[c]{0.6\linewidth}
        {\popdisplay\fontsize{11}{12.5}\selectfont\color{white}\MakeUppercase{Passe à OnBuch+}}\par\vspace{4pt}
        {\raggedright\popsemi\fontsize{8.4}{11}\selectfont\color{white}Tous les cours signés, Tuteur IA illimité, fiches PDF — onglet OnBuch+ de l'app\par}
      \end{minipage}
    \end{tcolorbox}
  \end{minipage}
  \vspace{8pt}
  \popcontenu
  \vfill
  \signatureonbuch
  \restoregeometry
  \newpage
}

% Rangée « Ce cours contient » (page de garde)
\newcommand{\poptile}[3]{% #1 couleur #2 gros texte #3 légende
  \begin{tcolorbox}[enhanced,colback=#1,colframe=popink,boxrule=2pt,arc=9pt,shadow={2.5pt}{-2.5pt}{0pt}{popink},
      left=6pt,right=6pt,top=9pt,bottom=9pt,halign=center,valign=center,height=1.75cm,width=\linewidth,nobeforeafter]
    {\popdisplay\fontsize{12.5}{13}\selectfont\color{white}\MakeUppercase{#2}}\par\vspace{3pt}
    {\centering\popsemi\fontsize{7.8}{9.5}\selectfont\color{white}#3\par}
  \end{tcolorbox}}
\newcommand{\popcontenu}{%
  \par\noindent{\popxboldls\fontsize{7.5}{7.5}\selectfont\color{popmuted}CE COURS SIGNÉ CONTIENT}\par\vspace{7pt}
  \noindent\begin{minipage}[t]{0.235\linewidth}\poptile{popblue}{Cours}{complet, illustré, pas à pas}\end{minipage}\hfill
  \begin{minipage}[t]{0.235\linewidth}\poptile{poporange}{Méthodes}{et exercices résolus}\end{minipage}\hfill
  \begin{minipage}[t]{0.235\linewidth}\poptile{poppink}{Exercices}{type examen + corrigés}\end{minipage}\hfill
  \begin{minipage}[t]{0.235\linewidth}\poptile{popgreen}{Fiche bilan}{l'essentiel à retenir}\end{minipage}\par}

% Sommaire
\newtcolorbox{sommaire}{enhanced,colback=white,colframe=popink,boxrule=2.2pt,arc=10pt,
  drop shadow=popink,left=18pt,right=18pt,top=13pt,bottom=13pt,before skip=6pt,after skip=20pt,
  before upper={\poppill{Sommaire}{poppurple}{white}\par\vspace{9pt}\popsemi\fontsize{10.6}{14.5}\selectfont\color{popink}}}

% Signature de fin de cours — poussée tout en bas de la dernière page
\newcommand{\findecours}{%
  \par\vfill\nopagebreak
  \begin{center}\popsquiggle{poporange}\end{center}\vspace{4pt}
  \signatureonbuch
}
\AtBeginDocument{\def\llbracket{\mathopen{[\mkern-3.5mu[}}\def\rrbracket{\mathclose{]\mkern-3.5mu]}}} % intervalles d'entiers (glyphe ⟦ absent de la police)
% Glyphes absents de Fira Math : redéfinis à partir de glyphes disponibles
\AtBeginDocument{%
  \def\setminus{\mathbin{\backslash}}\let\smallsetminus\setminus
  \def\vdots{\mathord{\vbox{\baselineskip3.2pt\lineskiplimit0pt\kern4pt\hbox{.}\hbox{.}\hbox{.}}}}%
  \def\ddots{\mathinner{\mkern1mu\raise6.5pt\hbox{.}\mkern2mu\raise3.6pt\hbox{.}\mkern2mu\raise0.7pt\hbox{.}\mkern1mu}}%
  \def\triangle{\mathord{\scalebox{0.9}{$\symup{\Delta}$}}}%
  \def\nabla{\mathord{\raisebox{\depth}{\scalebox{1}[-1]{$\symup{\Delta}$}}}}%
  \def\checkmark{\popcheck{popdark}}%
  \def\star{\mathbin{\ast}}\def\bigstar{\mathord{\ast}}%
  \def\diamond{\mathbin{\scalebox{0.6}{$\Diamond$}}}%
  \def\bigcup{\mathop{\mathchoice{\raisebox{-0.3em}{\scalebox{1.7}{$\cup$}}}{\scalebox{1.25}{$\cup$}}{\cup}{\cup}}\displaylimits}%
  \def\bigcap{\mathop{\mathchoice{\raisebox{-0.3em}{\scalebox{1.7}{$\cap$}}}{\scalebox{1.25}{$\cap$}}{\cap}{\cap}}\displaylimits}%
  \def\lesseqqgtr{\mathrel{\lessgtr}}\def\bigoplus{\mathop{\oplus}\displaylimits}%
}
\providecommand{\ssi}{si et seulement si} % abréviation fréquente
\AtBeginDocument{\providecommand{\wideparen}{\overparen}} % arc de cercle

%% FIGFIX-BEGIN v4 — correctifs globaux des figures (docs/onbuch-generation/tools/figfix)
\usepackage{adjustbox}\usepackage{etoolbox}
% toute figure TikZ/pgfplots plus large que la ligne est réduite à la largeur disponible
\BeforeBeginEnvironment{tikzpicture}{\begin{adjustbox}{max width=\linewidth}}
\AfterEndEnvironment{tikzpicture}{\end{adjustbox}}
% légendes pgfplots lisibles par-dessus les courbes
\pgfplotsset{every axis/.append style={legend style={fill=white,fill opacity=0.92,text opacity=1,draw=black!15,inner sep=3pt}}}
% étiquettes d'arêtes (syntaxe edge["..."]) avec fond blanc
\tikzset{every edge quotes/.append style={fill=white,inner sep=1.2pt}}
%% FIGFIX-END
\begin{document}

\pagegarde{Découverte d’un environnement d’IA}{Informatique}{3ème}{Informatique – classe de 3ème}{N01}{4 séances (fiche de progression harmonisée)}{Cette leçon initie les élèves de 3ème à l'Intelligence Artificielle : définition, usages concrets dans la vie courante camerounaise, avantages et limites. Elle développe aussi la capacité à dialoguer efficacement avec une IA et à adopter une attitude critique et responsable.
\vspace{4pt}
\begin{multicols}{2}\small
\begin{itemize}
\item À la fin de cette leçon, je sais définir l'Intelligence Artificielle et la distinguer d'un programme informatique classique
\item À la fin de cette leçon, je sais citer au moins cinq domaines d'application de l'IA dans la vie courante
\item À la fin de cette leçon, je sais identifier des exemples d'IA utilisés au Cameroun (mobile money, agriculture, santé, réseaux sociaux)
\item À la fin de cette leçon, je sais énumérer les avantages et les inconvénients de l'IA
\item À la fin de cette leçon, je sais expliquer ce que signifie une utilisation responsable de l'IA
\item À la fin de cette leçon, je sais formuler une consigne claire (prompt) pour échanger efficacement avec une IA
\end{itemize}\end{multicols}}

\begin{sommaire}
\begin{multicols}{2}
\begin{enumerate}
\item Généralités sur l'Intelligence Artificielle
\item Importance et usages de l'IA dans la vie courante
\item Avantages et inconvénients de l'IA
\item Utilisation responsable de l'IA
\item Apprendre à échanger avec une IA
\item Activité d'intégration
\item Méthodes et exercices résolus
\item Exercices
\item Corrigés des exercices
\item Fiche bilan
\end{enumerate}
\end{multicols}
\end{sommaire}

\begin{objectifs}
\begin{itemize}
\item À la fin de cette leçon, je sais définir l'Intelligence Artificielle et la distinguer d'un programme informatique classique
\item À la fin de cette leçon, je sais citer au moins cinq domaines d'application de l'IA dans la vie courante
\item À la fin de cette leçon, je sais identifier des exemples d'IA utilisés au Cameroun (mobile money, agriculture, santé, réseaux sociaux)
\item À la fin de cette leçon, je sais énumérer les avantages et les inconvénients de l'IA
\item À la fin de cette leçon, je sais expliquer ce que signifie une utilisation responsable de l'IA
\item À la fin de cette leçon, je sais formuler une consigne claire (prompt) pour échanger efficacement avec une IA
\item À la fin de cette leçon, je sais vérifier et critiquer une réponse fournie par une IA
\item À la fin de cette leçon, je sais rédiger et justifier une conclusion sur l'usage de l'IA dans une situation concrète
\end{itemize}
\end{objectifs}

\begin{prerequis}
\begin{itemize}
\item Notions de base sur l'ordinateur et ses périphériques (classe de 4ème)
\item Notion de logiciel et d'application (classe de 4ème)
\item Utilisation élémentaire d'Internet et d'un moteur de recherche
\item Notion de données et d'information
\item Lecture et rédaction d'un texte argumentatif simple
\end{itemize}
\end{prerequis}

\newpage

\coursec{Généralités sur l'Intelligence Artificielle}

Aminatou, élève en classe de 3ème à Yaoundé, rentre de l'école et pose son cartable. Dans le salon, le haut-parleur intelligent de son grand frère est allumé. Elle demande à voix haute : « Quelle heure fait-il à Douala ? » La machine répond aussitôt, en français, avec une voix presque humaine. Quelques minutes plus tard, le même appareil lui suggère une chanson de Locko, « parce qu'elle aime la musique camerounaise ». Aminatou reste songeuse : comment cette machine arrive-t-elle à \emph{comprendre} ce qu'elle dit, à \emph{chercher} une réponse et même à \emph{deviner} ses goûts musicaux ? Est-ce vraiment de l'intelligence, comme celle d'un humain ? Et peut-on faire confiance à tout ce que la machine répond ?

C'est exactement le genre de questions que se posent aujourd'hui des millions de personnes, des élèves comme toi jusqu'aux grands chercheurs. Cette leçon va nous permettre de découvrir ce qu'est l'\cle{Intelligence Artificielle}, d'où elle vient, comment elle fonctionne, et pourquoi il faut apprendre à l'utiliser de manière responsable.

\textbf{Problématique :} qu'est-ce que l'Intelligence Artificielle, comment une machine peut-elle « apprendre », et en quoi est-ce différent de l'intelligence humaine ?

\courssub{Qu'est-ce que l'Intelligence Artificielle ?}

\textbf{Intuition.} Depuis toujours, l'être humain fabrique des machines pour l'aider : la roue pour porter des charges, la calculatrice pour calculer vite, le téléphone pour parler à distance. Mais certaines tâches semblaient réservées au cerveau humain : comprendre une phrase, reconnaître un visage, apprendre de ses erreurs, jouer aux échecs. L'Intelligence Artificielle est née d'une question simple : \emph{une machine peut-elle réaliser, elle aussi, ces tâches qui demandent de l'intelligence ?}

\begin{definition}[Intelligence Artificielle]
L'\cle{Intelligence Artificielle} (IA) est un \textbf{ensemble de techniques informatiques} permettant à une machine de réaliser des tâches qui demandent habituellement de l'\textbf{intelligence humaine} : raisonner, apprendre, comprendre le langage, percevoir le monde (voir, écouter) et prendre des décisions.
\end{definition}

Décomposons cette définition pour bien la comprendre :

\begin{itemize}
\item \textbf{« Ensemble de techniques informatiques »} : l'IA n'est pas un objet magique, c'est de l'informatique. Ce sont des programmes, des mathématiques et des ordinateurs très puissants.
\item \textbf{« Une machine »} : un ordinateur, un téléphone, un robot, un haut-parleur intelligent comme celui d'Aminatou.
\item \textbf{« Des tâches qui demandent habituellement de l'intelligence humaine »} : reconnaître un ami sur une photo, traduire une phrase du français vers l'anglais, comprendre une question posée à voix haute, prédire quelle vidéo va te plaire.
\end{itemize}

\begin{attention}[L'IA ne « pense » pas comme un humain]
Erreur fréquente : croire que l'IA possède une conscience, des émotions ou des intentions. \textbf{C'est faux.} Une IA ne « comprend » pas au sens humain : elle \textbf{calcule des probabilités} à partir de données. Quand le haut-parleur d'Aminatou répond, il n'a pas « réfléchi » comme toi : il a calculé la réponse la plus probable. L'IA n'a ni conscience, ni sentiments, ni bon sens. C'est un outil très performant, mais qui reste une machine.
\end{attention}

\courssub{Une brève histoire de l'IA}

L'IA n'est pas née avec les téléphones portables ! C'est une aventure scientifique qui dure depuis plus de 70 ans. Retenons les grandes étapes.

\textbf{1950 : Alan Turing et son célèbre test.} Le mathématicien britannique \cle{Alan Turing}, l'un des pères de l'informatique, se demande : « Les machines peuvent-elles penser ? » Il propose une expérience, le \cle{test de Turing} : un humain discute \emph{par écrit} (au clavier) avec deux interlocuteurs cachés : une personne et une machine. Si, à la fin, il n'arrive pas à distinguer laquelle est la machine, alors celle-ci peut être dite « intelligente ».

\begin{savaistu}[Le test de Turing]
Le test de Turing propose qu'une machine soit dite « intelligente » si un humain, en discutant par écrit avec elle, ne parvient pas à la distinguer d'une personne réelle. Alan Turing est aussi un héros de l'histoire : pendant la Seconde Guerre mondiale, ses travaux ont permis de casser les codes secrets ennemis. On le considère comme le père de l'informatique moderne.
\end{savaistu}

\textbf{1956 : naissance du mot « Intelligence Artificielle ».} Lors d'une conférence au Dartmouth College (États-Unis), le chercheur John McCarthy invente officiellement l'expression \emph{Artificial Intelligence}. C'est la naissance de l'IA comme domaine scientifique.

\textbf{1997 : Deep Blue bat le champion du monde d'échecs.} L'ordinateur Deep Blue, construit par IBM, bat le champion du monde Garry Kasparov. Pour la première fois, une machine bat le meilleur humain dans un jeu de stratégie réputé très difficile.

\textbf{2012 : l'essor du deep learning.} Grâce aux \textbf{données massives} (des milliards de photos, de textes, de vidéos disponibles sur Internet) et à des ordinateurs de plus en plus puissants, une technique appelée \emph{apprentissage profond} (deep learning) fait faire un bond gigantesque à l'IA : reconnaissance d'images, traduction automatique, assistants vocaux.

\textbf{2022 : les chatbots grand public.} Des agents conversationnels comme ChatGPT deviennent accessibles à tous, y compris au Cameroun, depuis un simple téléphone. L'IA entre dans la vie quotidienne : rédaction, aide aux devoirs, traduction, conseils.

\begin{popfigure}
\begin{center}
\begin{tikzpicture}[font=\small]
\draw[-{Stealth}, thick, popline] (0,0) -- (14.6,0);
\foreach \x/\annee/\texte/\coul in {
  0.8/1950/{Test de\\Turing}/popblue,
  3.6/1956/{Naissance du\\terme « IA »}/popgreen,
  6.4/1997/{Deep Blue bat\\Kasparov (échecs)}/poporange,
  9.6/2012/{Essor du\\deep learning}/poppurple,
  13.2/2022/{Chatbots\\grand public}/poppink}{
  \draw[thick, \coul] (\x,0) -- (\x,0.35);
  \node[above=0.15cm, align=center, text=\coul, font=\bfseries\footnotesize] at (\x,0.35) {\texte};
  \node[below=0.1cm, font=\bfseries] at (\x,0) {\annee};
  \fill[\coul] (\x,0) circle (2.5pt);
}
\end{tikzpicture}
\end{center}
\legende{Frise chronologique : les grandes dates de l'Intelligence Artificielle.}
\end{popfigure}

\courssub{Programme classique ou Intelligence Artificielle ?}

\textbf{Intuition.} Tu as déjà vu (ou tu verras) comment fonctionne un programme classique : un humain écrit des instructions précises, et l'ordinateur les exécute. Par exemple, pour calculer la moyenne de tes notes, le programmeur écrit : « additionne les notes, divise par le nombre de notes ». La machine ne fait \emph{que} cela, toujours de la même façon.

Mais comment écrire des instructions pour reconnaître un visage ? Impossible de décrire par des règles fixes tous les visages du monde ! C'est là que l'IA change tout : au lieu de donner les règles à la machine, on lui donne des \textbf{exemples}, et c'est elle qui \emph{découvre} les règles.

\begin{propriete}[Programme classique]
Un \cle{programme classique} suit des \textbf{instructions fixes} écrites par un humain. On fournit à la machine des \textbf{données} et des \textbf{règles} : elle applique les règles aux données et produit un \textbf{résultat}.
\[ \text{Données} + \text{Règles (écrites par l'humain)} \longrightarrow \text{Résultat} \]
\end{propriete}

\begin{propriete}[Système d'IA]
Un \cle{système d'IA} \textbf{apprend à partir d'exemples}. On fournit à la machine des \textbf{données} et les \textbf{résultats attendus} : elle découvre elle-même les \textbf{règles} qui relient les données aux résultats.
\[ \text{Données} + \text{Résultats (exemples)} \longrightarrow \text{la machine découvre les Règles} \]
\end{propriete}

\begin{popfigure}
\begin{center}
\begin{tikzpicture}[font=\small, node distance=0.45cm,
  boite/.style={draw, thick, rounded corners=2pt, align=center, minimum width=2.6cm, minimum height=0.85cm},
  fleche/.style={-{Stealth}, thick}]
% --- Programme classique (gauche) ---
\node[font=\bfseries, text=popblue] at (2.6,3.4) {Programme classique};
\node[boite, fill=popblueL, draw=popblue] (donn1) at (1.2,2.3) {Données};
\node[boite, fill=popblueL, draw=popblue, align=center] (regl1) at (4.0,2.3){Règles\\(écrites par l'humain)};
\node[boite, fill=popcream, draw=popdark] (prog1) at (2.6,1.1) {Programme};
\node[boite, fill=popgreenL, draw=popgreen] (res1) at (2.6,0) {Résultat};
\draw[fleche, popblue] (donn1.south) -- (prog1.north west);
\draw[fleche, popblue] (regl1.south) -- (prog1.north east);
\draw[fleche, popdark] (prog1) -- (res1);
% --- IA (droite) ---
\node[font=\bfseries, text=poppurple] at (10.2,3.4) {Système d'IA};
\node[boite, fill=poppurpleL, draw=poppurple, align=center] (donn2) at (8.8,2.3){Données\\(exemples)};
\node[boite, fill=poppurpleL, draw=poppurple, align=center] (res2) at (11.6,2.3){Résultats\\attendus};
\node[boite, fill=popcream, draw=popdark, align=center] (ia2) at (10.2,1.1){La machine\\apprend};
\node[boite, fill=poporangeL, draw=poporange, align=center] (regl2) at (10.2,0){Règles découvertes\\par la machine};
\draw[fleche, poppurple] (donn2.south) -- (ia2.north west);
\draw[fleche, poppurple] (res2.south) -- (ia2.north east);
\draw[fleche, popdark] (ia2) -- (regl2);
\end{tikzpicture}
\end{center}
\legende{À gauche, le programme classique applique des règles écrites par l'humain. À droite, le système d'IA découvre les règles à partir d'exemples.}
\end{popfigure}

\begin{exemplebox}[Correcteur orthographique contre chatbot]
Un \textbf{correcteur orthographique classique} applique des \emph{règles fixes} écrites par des humains : « après "les", un nom pluriel prend un s », « "manger" se conjugue ainsi au présent ». Il ne connaît que ce qu'on lui a programmé.

Un \textbf{chatbot d'IA}, lui, a \emph{appris} le langage en analysant des \textbf{milliards de phrases} écrites par des humains. Personne ne lui a dicté les règles de grammaire : il les a découvertes tout seul en observant les exemples. C'est pourquoi il peut répondre à des questions que personne n'avait prévues... et aussi se tromper !
\end{exemplebox}

\courssub{L'apprentissage automatique (machine learning)}

Comment la machine « apprend-elle » à partir d'exemples ? C'est le rôle de l'\cle{apprentissage automatique}.

\begin{definition}[Apprentissage automatique]
L'\cle{apprentissage automatique} (en anglais \emph{machine learning}) est une méthode d'IA où la machine \textbf{apprend à partir de nombreux exemples} au lieu de suivre uniquement des instructions programmées. Plus elle analyse de données, plus elle \textbf{améliore ses réponses}.
\end{definition}

\textbf{L'idée avec un exemple concret.} Imaginons qu'on veuille fabriquer une application capable de reconnaître un fruit : banane ou mangue ?

\begin{enumerate}
\item \textbf{On rassemble des exemples} : on montre à la machine des milliers de photos de bananes et de mangues, en lui disant pour chacune la bonne réponse (« ceci est une banane », « ceci est une mangue »).
\item \textbf{La machine s'entraîne} : elle compare les photos et remarque toute seule des indices : la forme allongée, la couleur jaune ou verte, la taille... Elle construit ses propres « règles ».
\item \textbf{On teste} : on lui montre une \emph{nouvelle} photo, jamais vue. Elle calcule : « à 95 \%, c'est une mangue ».
\item \textbf{On corrige ses erreurs} : si elle se trompe, on ajuste, et elle s'améliore. C'est exactement comme toi quand tu fais des exercices : plus tu t'entraînes, meilleur tu deviens !
\end{enumerate}

\begin{attention}[Des données de mauvaise qualité donnent une mauvaise IA]
Piège à éviter : si les exemples montrés à la machine sont faux, incomplets ou injustes, l'IA apprendra mal. Par exemple, une IA entraînée uniquement avec des photos de bananes jaunes pourra croire qu'une banane verte... n'est pas une banane ! \textbf{La qualité des données est essentielle.} C'est pourquoi une IA peut parfois donner des réponses fausses ou injustes : elle ne fait que reproduire ce qu'elle a vu dans ses données.
\end{attention}

\courssub{IA faible et IA forte}

Toutes les IA ne se valent pas. Les scientifiques distinguent deux grandes catégories.

\begin{definition}[IA faible et IA forte]
\begin{itemize}
\item L'\cle{IA faible} (ou IA étroite) est \textbf{spécialisée dans une seule tâche} : reconnaissance vocale, traduction, recommandation de vidéos, jeu d'échecs... C'est la seule qui existe aujourd'hui. Toutes les IA que tu utilises sont des IA faibles.
\item L'\cle{IA forte} serait une \textbf{intelligence générale}, capable de tout faire comme un humain : apprendre n'importe quoi, comprendre le monde, s'adapter à toutes les situations. \textbf{Elle n'existe pas encore} : elle relève pour l'instant de la recherche... et de la science-fiction !
\end{itemize}
\end{definition}

\begin{aretenir}[L'essentiel sur IA faible et IA forte]
L'IA qui bat le champion du monde d'échecs ne sait même pas faire cuire un œuf ni tenir une conversation sur la pluie et le beau temps : elle est très forte \emph{dans son domaine uniquement}. Ne confonds jamais la performance d'une IA sur une tâche précise avec une intelligence humaine complète.
\end{aretenir}

\courssub{Les capacités de l'IA : voir, écouter, parler, prédire}

Pour bien retenir ce que sait faire l'IA, on peut comparer ses capacités à celles des humains :

\begin{center}
\begin{tabularx}{\linewidth}{|l|X|X|}
\hline
\rowcolor{popblueL} \textbf{Capacité} & \textbf{Ce que fait l'IA} & \textbf{Exemples d'applications} \\
\hline
\textbf{Voir} & Reconnaissance d'images : identifier des visages, des objets, des animaux sur des photos & Déverrouillage du téléphone par le visage, tri automatique des photos \\
\hline
\textbf{Écouter} & Reconnaissance vocale : transformer la voix en texte et comprendre une demande & Assistants vocaux, haut-parleurs intelligents, dictée vocale \\
\hline
\textbf{Parler et écrire} & Produire des phrases en langage humain, discuter, traduire & Chatbots, traduction automatique français--anglais \\
\hline
\textbf{Prédire} & Anticiper un résultat à partir de données passées & Recommandations de vidéos et de musiques, prévision météo \\
\hline
\end{tabularx}
\end{center}

Le haut-parleur intelligent d'Aminatou combine donc plusieurs capacités : il \emph{écoute} sa question (reconnaissance vocale), il \emph{cherche} une réponse, il la \emph{dit} à voix haute, et il \emph{prédit} les musiques qu'elle aimera. Tout cela reste de l'IA faible : chaque capacité est une tâche spécialisée.

\begin{exoresolu}[Classer des applications : programme classique ou IA ?]
\textbf{Énoncé.} Pour chacune des applications suivantes, indique s'il s'agit plutôt d'un programme classique ou d'un système d'IA, et justifie ta réponse :
\begin{enumerate}
\item Une calculatrice qui additionne deux nombres.
\item Une application qui reconnaît les visages sur les photos d'un téléphone.
\item Un tableur qui calcule la moyenne des notes d'une classe.
\item Une application qui traduit une phrase du français vers l'ewondo.
\end{enumerate}
\tcblower
\textbf{Solution détaillée.}
\begin{enumerate}
\item \textbf{Programme classique.} L'addition suit des règles fixes écrites par le programmeur : données + règles $\rightarrow$ résultat. Aucun apprentissage n'est nécessaire.
\item \textbf{Système d'IA.} Impossible d'écrire des règles décrivant tous les visages du monde : l'application a \emph{appris} à partir de milliers d'exemples de photos (capacité « voir »).
\item \textbf{Programme classique.} Le calcul de la moyenne est une formule fixe (somme des notes divisée par leur nombre) : ce sont des instructions précises, appliquées toujours de la même façon.
\item \textbf{Système d'IA.} La traduction ne peut pas reposer sur des règles simples (chaque langue a trop d'exceptions) : l'application a appris en analysant d'énormes quantités de phrases traduites (capacité « parler et écrire »).
\end{enumerate}
\textbf{À retenir :} si la tâche peut être décrite entièrement par des règles fixes, un programme classique suffit ; si elle demande de reconnaître, comprendre ou prédire à partir d'exemples variés, on utilise l'IA.
\end{exoresolu}

\begin{exercice}[À toi de jouer]
\begin{enumerate}
\item Donne, avec tes propres mots, la définition de l'Intelligence Artificielle.
\item En quelle année le terme « Intelligence Artificielle » a-t-il été inventé ? Qui a proposé le test portant son nom, et en quoi consiste ce test ?
\item Explique la différence entre un programme classique et un système d'IA, en t'aidant du schéma « données + règles » contre « données + résultats ».
\item Le haut-parleur intelligent d'Aminatou est-il une IA faible ou une IA forte ? Justifie.
\item Cite deux capacités de l'IA et donne pour chacune un exemple d'application utilisée au quotidien.
\end{enumerate}
\end{exercice}

\begin{corrige}[Exercice : À toi de jouer]
\begin{enumerate}
\item L'Intelligence Artificielle est un ensemble de techniques informatiques permettant à une machine de réaliser des tâches qui demandent habituellement de l'intelligence humaine : raisonner, apprendre, comprendre le langage, percevoir.
\item Le terme a été inventé en \textbf{1956} (conférence de Dartmouth, John McCarthy). \textbf{Alan Turing} a proposé en 1950 le test de Turing : une machine est dite « intelligente » si un humain, en discutant par écrit avec elle, ne parvient pas à la distinguer d'une personne réelle.
\item Un programme classique suit des instructions fixes écrites par l'humain : on lui donne des données et des règles, il produit un résultat. Un système d'IA apprend à partir d'exemples : on lui donne des données et les résultats attendus, et il découvre lui-même les règles.
\item C'est une \textbf{IA faible} : il est spécialisé dans des tâches précises (reconnaissance vocale, recherche de réponses, recommandations musicales) et ne possède pas d'intelligence générale comme un humain.
\item Par exemple : \emph{voir} (reconnaissance d'images : déverrouillage du téléphone par le visage) ; \emph{prédire} (recommandations de vidéos ou de musiques). On pouvait aussi citer \emph{écouter} (assistants vocaux) et \emph{parler/écrire} (chatbots, traduction).
\end{enumerate}
\end{corrige}

\begin{aretenir}[Bilan de la section]
\begin{itemize}
\item L'\cle{IA} est un ensemble de techniques informatiques permettant à une machine d'imiter certaines capacités humaines (raisonner, apprendre, comprendre, percevoir).
\item Née avec le \textbf{test de Turing} (1950) et le terme « Intelligence Artificielle » (1956), elle connaît un essor récent grâce aux \textbf{données massives} et aux ordinateurs puissants.
\item Contrairement au \textbf{programme classique} qui suit des instructions fixes, le système d'IA \textbf{apprend à partir d'exemples} : c'est l'\cle{apprentissage automatique}.
\item Toutes les IA actuelles sont des \textbf{IA faibles}, spécialisées dans une tâche ; l'\textbf{IA forte} (intelligence générale) n'existe pas encore.
\item L'IA sait \textbf{voir, écouter, parler, écrire et prédire}... mais elle ne pense pas comme un humain : elle calcule des probabilités à partir de données.
\end{itemize}
\end{aretenir}

\coursec{Importance et usages de l'IA dans la vie courante}

Maintenant que tu sais ce qu'est l'IA, une question se pose : \emph{à quoi sert-elle concrètement ?} Reprenons la journée d'Aminatou. Le matin, son téléphone se déverrouille en reconnaissant son visage. Sur le chemin de l'école, son père consulte une application qui prédit les embouteillages à Yaoundé. Le soir, elle révise ses leçons avec un chatbot qui répond à ses questions de sciences. Sans même s'en rendre compte, Aminatou utilise l'IA \textbf{plusieurs fois par jour}. Et toi aussi, très probablement !

\textbf{Problématique :} dans quels domaines de la vie courante l'IA est-elle utilisée, et pourquoi est-elle devenue si importante ?

\courssub{Pourquoi l'IA est-elle devenue si importante ?}

L'IA est partout parce qu'elle apporte trois grands atouts aux machines :

\begin{itemize}
\item \textbf{La rapidité} : une IA analyse en quelques secondes des millions de données qu'un humain mettrait des années à étudier.
\item \textbf{La disponibilité} : une machine ne dort jamais, ne se fatigue pas : elle peut répondre 24 heures sur 24, 7 jours sur 7.
\item \textbf{La précision sur les tâches répétitives} : pour reconnaître des images ou trier des données en grand nombre, l'IA fait souvent moins d'erreurs qu'un humain fatigué.
\end{itemize}

C'est pourquoi les entreprises, les hôpitaux, les écoles et même les administrations utilisent de plus en plus l'IA : elle permet de \textbf{gagner du temps} et d'\textbf{aider les humains à prendre de meilleures décisions}.

\courssub{Les usages de l'IA dans la vie courante}

Voici les principaux domaines où tu croises l'IA chaque jour, souvent sans le savoir.

\textbf{1. La communication et les réseaux sociaux.} Les applications comme TikTok, YouTube ou Facebook utilisent l'IA pour te \emph{recommander} des vidéos adaptées à tes goûts. Les filtres qui embellissent tes photos ou ajoutent des oreilles de chien sur ton visage utilisent la reconnaissance d'images.

\textbf{2. L'éducation.} Des chatbots répondent aux questions des élèves, expliquent des leçons, corrigent des exercices. Des applications de traduction aident à comprendre un texte en anglais. Attention cependant : l'IA doit t'aider à \emph{comprendre}, pas faire tes devoirs à ta place !

\textbf{3. La santé.} L'IA aide les médecins à analyser des radios pour détecter des maladies plus tôt, à rechercher de nouveaux médicaments, ou à gérer les rendez-vous d'un hôpital.

\textbf{4. L'agriculture.} Des applications reconnaissent les maladies des plantes à partir d'une simple photo des feuilles : un planteur de cacao ou de manioc peut ainsi soigner ses cultures plus vite. D'autres prédisent la météo pour choisir le bon moment des semis.

\textbf{5. Les transports.} Les applications de navigation calculent le meilleur itinéraire en tenant compte du trafic en temps réel. Les voitures autonomes (qui se conduisent presque seules) utilisent l'IA pour voir la route.

\textbf{6. La sécurité et les banques.} Les banques utilisent l'IA pour détecter les fraudes (par exemple un paiement suspect par mobile money). La reconnaissance faciale sécurise les téléphones et certains bâtiments.

\textbf{7. Le commerce.} Les boutiques en ligne te suggèrent des produits (« les clients qui ont acheté ceci ont aussi acheté cela »). Des chatbots répondent aux clients des entreprises à toute heure.

\begin{savaistu}[L'IA au Cameroun]
Au Cameroun, l'IA commence à transformer la vie quotidienne : paiements par mobile money sécurisés, applications agricoles pour les planteurs, chatbots des banques et des opérateurs téléphoniques, services administratifs en ligne. Le pays forme de jeunes développeurs dans ce domaine : peut-être seras-tu l'un d'eux demain !
\end{savaistu}

\begin{exemplebox}[Une journée avec l'IA, sans le savoir]
Recense les usages de l'IA dans cette journée d'Aminatou :
\begin{itemize}
\item déverrouillage du téléphone par le visage $\rightarrow$ \textbf{reconnaissance d'images} ;
\item application de navigation pour éviter les embouteillages $\rightarrow$ \textbf{prédiction} ;
\item vidéos recommandées sur YouTube $\rightarrow$ \textbf{prédiction des goûts} ;
\item chatbot qui aide à réviser $\rightarrow$ \textbf{compréhension et production du langage} ;
\item filtre photo amusant $\rightarrow$ \textbf{reconnaissance d'images}.
\end{itemize}
Conclusion : l'IA est déjà partout, même dans la poche d'un élève de 3ème !
\end{exemplebox}

\begin{exercice}[À toi de jouer]
\begin{enumerate}
\item Cite trois domaines de la vie courante où l'IA est utilisée, et décris pour chacun un usage précis.
\item Pourquoi dit-on que l'IA est « disponible 24 heures sur 24 » ? Est-ce le cas d'un humain ?
\item Trouve, dans ta propre vie quotidienne, deux moments où tu utilises (sans forcément le savoir) une application reposant sur l'IA.
\end{enumerate}
\end{exercice}

\begin{corrige}[Exercice : À toi de jouer]
\begin{enumerate}
\item Par exemple : la \textbf{santé} (analyse des radios pour détecter des maladies), l'\textbf{agriculture} (reconnaissance des maladies des plantes à partir de photos), les \textbf{transports} (calcul du meilleur itinéraire selon le trafic). On pouvait aussi citer l'éducation, le commerce, la sécurité bancaire, les réseaux sociaux.
\item Une machine ne dort pas et ne se fatigue pas : un chatbot ou un serveur peut répondre à toute heure du jour et de la nuit. Un humain, lui, a besoin de sommeil et de repos.
\item Réponses personnelles possibles : vidéos recommandées sur TikTok/YouTube, déverrouillage facial du téléphone, correcteur ou traducteur automatique, assistant vocal, filtres photo.
\end{enumerate}
\end{corrige}

\begin{aretenir}[Bilan de la section]
\begin{itemize}
\item L'IA est importante car elle est \textbf{rapide}, \textbf{toujours disponible} et \textbf{précise} sur les tâches répétitives.
\item On la retrouve dans de nombreux domaines : \textbf{communication, éducation, santé, agriculture, transports, sécurité, commerce}.
\item Nous utilisons l'IA chaque jour, souvent sans nous en rendre compte (recommandations, reconnaissance faciale, traduction).
\end{itemize}
\end{aretenir}

\coursec{Avantages et inconvénients de l'IA}

L'IA rend de grands services... mais elle pose aussi de vrais problèmes. Un bon utilisateur d'IA est quelqu'un qui connaît \textbf{les deux faces de la médaille}. C'est ce que nous allons voir maintenant.

\textbf{Problématique :} quels sont les avantages et les limites de l'IA, et pourquoi ne faut-il pas lui faire une confiance aveugle ?

\courssub{Les avantages de l'IA}

\begin{itemize}
\item \textbf{Gain de temps} : l'IA accomplit en quelques secondes des tâches longues et fastidieuses (trier des milliers de photos, chercher une information, rédiger un premier brouillon).
\item \textbf{Aide à la décision} : en analysant beaucoup de données, elle aide les médecins, les agriculteurs ou les enseignants à mieux décider.
\item \textbf{Disponibilité permanente} : un chatbot répond jour et nuit, même le dimanche, même pendant les vacances.
\item \textbf{Accessibilité} : la traduction automatique et les assistants vocaux aident les personnes qui ne lisent pas l'anglais, ou les personnes handicapées (dictée vocale pour les malvoyants).
\item \textbf{Précision sur les tâches répétitives} : bien entraînée, l'IA fait moins d'erreurs qu'un humain fatigué sur des tâches monotones.
\end{itemize}

\courssub{Les inconvénients et les dangers de l'IA}

\begin{itemize}
\item \textbf{Les erreurs et les « inventions »} : une IA peut affirmer des choses \textbf{fausses} avec beaucoup d'assurance. Un chatbot peut inventer une date, un nom ou une citation qui n'existe pas. On appelle parfois cela une \emph{hallucination} de l'IA.
\item \textbf{Les biais et les injustices} : l'IA reproduit les défauts de ses données d'apprentissage. Si les exemples étaient injustes ou incomplets, ses réponses le seront aussi.
\item \textbf{La dépendance} : à force de tout demander à l'IA, on risque de ne plus réfléchir par soi-même. Un élève qui recopie les réponses d'un chatbot n'apprend rien !
\item \textbf{Les questions d'emploi} : certaines tâches autrefois faites par des humains sont maintenant automatisées, ce qui transforme les métiers.
\item \textbf{La vie privée} : pour fonctionner, beaucoup d'IA collectent des données personnelles (voix, photos, habitudes). Que deviennent ces données ?
\item \textbf{Les mauvais usages} : l'IA peut servir à créer de fausses images ou de fausses vidéos (les \emph{deepfakes}), à propager des rumeurs ou à tricher.
\end{itemize}

\begin{attention}[Ne jamais croire une IA sur parole]
Erreur fréquente : penser que « si l'IA l'a dit, c'est vrai ». \textbf{Non !} Une IA calcule la réponse la plus \emph{probable}, pas la plus \emph{vraie}. Elle peut se tromper, inventer, ou être dépassée (ses connaissances s'arrêtent à la date de son entraînement). \textbf{Vérifie toujours} les informations importantes dans un livre, sur un site officiel ou auprès d'un adulte.
\end{attention}

\begin{center}
\begin{tabularx}{\linewidth}{|X|X|}
\hline
\rowcolor{popblueL} \textbf{Avantages} & \textbf{Inconvénients} \\
\hline
Rapidité et gain de temps & Erreurs et informations inventées \\
\hline
Disponibilité 24 h/24 & Biais et réponses injustes \\
\hline
Aide à la décision & Risque de dépendance et de paresse intellectuelle \\
\hline
Accessibilité (traduction, handicap) & Menaces sur la vie privée \\
\hline
Précision sur les tâches répétitives & Mauvais usages (deepfakes, triche, rumeurs) \\
\hline
\end{tabularx}
\end{center}

\begin{exoresolu}[Avantage ou inconvénient ?]
\textbf{Énoncé.} Aminatou utilise un chatbot pour préparer son exposé d'histoire. Voici ce qui s'est passé. Pour chaque fait, dis s'il illustre un avantage ou un inconvénient de l'IA, et explique pourquoi.
\begin{enumerate}
\item Le chatbot lui a fourni un plan d'exposé en 10 secondes.
\item Le chatbot a affirmé que le Cameroun est devenu indépendant en 1972, ce qui est faux.
\item Aminatou a recopié le texte du chatbot sans le comprendre, et n'a rien retenu pour l'interrogation suivante.
\end{enumerate}
\tcblower
\textbf{Solution détaillée.}
\begin{enumerate}
\item \textbf{Avantage} : le gain de temps. L'IA produit rapidement un brouillon ou un plan, ce qui laisse plus de temps pour approfondir.
\item \textbf{Inconvénient} : l'erreur (ou « hallucination »). L'IA a affirmé une fausse information avec assurance (la bonne date est le 1\textsuperscript{er} janvier 1960 pour le Cameroun francophone). Il faut toujours vérifier les faits importants.
\item \textbf{Inconvénient} : la dépendance. Recopier sans comprendre empêche d'apprendre : l'IA doit servir à \emph{mieux comprendre}, pas à remplacer le travail personnel.
\end{enumerate}
\end{exoresolu}

\begin{exercice}[À toi de jouer]
\begin{enumerate}
\item Cite trois avantages de l'IA et illustre chacun par un exemple.
\item Cite trois inconvénients ou dangers de l'IA et illustre chacun par un exemple.
\item Pourquoi une IA peut-elle donner une réponse injuste ou fausse ? Rappelle le rôle des données d'apprentissage.
\end{enumerate}
\end{exercice}

\begin{corrige}[Exercice : À toi de jouer]
\begin{enumerate}
\item Par exemple : gain de temps (un chatbot rédige un brouillon en quelques secondes) ; disponibilité (réponses à toute heure) ; accessibilité (traduction automatique, dictée vocale pour les malvoyants).
\item Par exemple : erreurs et inventions (une fausse date donnée avec assurance) ; biais (données d'apprentissage injustes reproduites) ; atteinte à la vie privée (collecte des données personnelles) ; mauvais usages (deepfakes, triche).
\item Parce qu'elle apprend à partir d'exemples : si les données d'apprentissage sont fausses, incomplètes ou injustes, l'IA reproduit ces défauts. Elle calcule la réponse la plus probable, pas la plus vraie.
\end{enumerate}
\end{corrige}

\begin{aretenir}[Bilan de la section]
\begin{itemize}
\item L'IA apporte \textbf{rapidité, disponibilité, aide à la décision et accessibilité}.
\item Mais elle peut \textbf{se tromper, inventer, reproduire des biais}, créer une \textbf{dépendance} et menacer la \textbf{vie privée}.
\item Règle d'or : \textbf{ne jamais faire une confiance aveugle à l'IA} ; toujours vérifier les informations importantes.
\end{itemize}
\end{aretenir}

\coursec{Utilisation responsable de l'IA}

Tu connais maintenant les forces et les faiblesses de l'IA. Reste la question la plus importante pour toi, élève de 3ème : \emph{comment bien l'utiliser ?} Utiliser l'IA de manière \textbf{responsable}, c'est en profiter tout en évitant ses pièges, pour toi-même et pour les autres.

\textbf{Problématique :} quelles règles faut-il respecter pour utiliser l'IA de façon responsable, honnête et sûre ?

\courssub{Les règles d'une utilisation responsable}

\begin{methode}[Les 6 règles d'or pour utiliser l'IA de manière responsable]
\begin{enumerate}
\item \textbf{Vérifier les informations.} L'IA peut se tromper ou inventer : contrôle toujours les faits importants (dates, définitions, calculs) avec un livre, un site officiel ou un enseignant.
\item \textbf{Réfléchir d'abord soi-même.} Cherche la solution par toi-même \emph{avant} de demander à l'IA. Utilise-la pour vérifier, comprendre ou approfondir, jamais pour éviter de travailler.
\item \textbf{Ne pas tricher.} Recopier un texte d'IA et le présenter comme ton propre travail est une forme de plagiat (et de tricherie à l'examen). Cite tes sources quand tu t'en inspires.
\item \textbf{Protéger ses données personnelles.} Ne donne jamais à une IA ton mot de passe, ton adresse, tes photos privées ou celles des autres. Ce que tu partages peut être conservé et réutilisé.
\item \textbf{Respecter les autres.} N'utilise pas l'IA pour créer de fausses images ou de fausses informations sur quelqu'un, ni pour te moquer ou harceler : c'est interdit et puni par la loi.
\item \textbf{Garder l'esprit critique.} Demande-toi toujours : « cette réponse est-elle logique ? d'où vient-elle ? qui l'a produite et pourquoi ? »
\end{enumerate}
\end{methode}

\begin{attention}[L'IA et les devoirs : la frontière entre aide et triche]
\begin{itemize}
\item \textbf{Utilisation responsable} : « Explique-moi la leçon sur les fractions », « Corrige mon exercice et dis-moi où je me suis trompé », « Propose-moi des questions pour réviser ».
\item \textbf{Utilisation irresponsable} : « Rédige mon exposé à ma place », puis le recopier sans le lire. Résultat : tu n'apprends rien, tu risques une sanction, et le jour du BEPC, l'IA ne sera pas dans la salle !
\end{itemize}
\end{attention}

\begin{savaistu}[L'IA et la loi]
Au Cameroun, la loi n° 2010/012 du 21 décembre 2010 relative à la cybersécurité et à la cybercriminalité punit notamment la diffusion de fausses informations, l'atteinte à la vie privée et les fraudes numériques. Créer ou partager une fausse image d'une personne, même « pour rire », peut avoir de graves conséquences. Un bon citoyen numérique utilise l'IA dans le respect des lois et des autres.
\end{savaistu}

\courssub{Être un bon citoyen numérique}

Utiliser l'IA de manière responsable, c'est aussi adopter les bons comportements du \cle{citoyen numérique} :

\begin{itemize}
\item \textbf{Honnêteté} : je ne triche pas, je cite mes sources, je ne me fais pas passer pour l'auteur d'un texte produit par une machine.
\item \textbf{Respect} : je n'utilise pas l'IA pour nuire, insulter, tromper ou humilier.
\item \textbf{Prudence} : je protège mes données personnelles et celles des autres.
\item \textbf{Esprit critique} : je vérifie avant de croire et avant de partager.
\end{itemize}

\begin{exoresolu}[Responsable ou pas ?]
\textbf{Énoncé.} Pour chaque situation, dis si l'utilisation de l'IA est responsable ou irresponsable, et justifie.
\begin{enumerate}
\item Brice demande à un chatbot de lui expliquer le théorème de Pythagore, puis refait les exercices tout seul.
\item Sandrine demande à l'IA de rédiger sa rédaction de français et la recopie mot pour mot.
\item Ulrich envoie à un chatbot une photo de sa carte d'identité « pour tester ».
\item Vanessa vérifie sur un site officiel la date donnée par le chatbot avant de l'écrire dans son exposé.
\end{enumerate}
\tcblower
\textbf{Solution détaillée.}
\begin{enumerate}
\item \textbf{Responsable} : l'IA sert à \emph{comprendre}, puis le travail personnel reprend le dessus. C'est exactement la bonne méthode.
\item \textbf{Irresponsable} : c'est du plagiat et de la triche ; Sandrine n'apprend rien et risque une sanction.
\item \textbf{Irresponsable} : il ne faut jamais confier ses données personnelles (documents d'identité, photos privées) à une IA ; elles peuvent être conservées et mal utilisées.
\item \textbf{Responsable} : vérifier les informations avant de les utiliser est la règle d'or n° 1.
\end{enumerate}
\end{exoresolu}

\begin{exercice}[À toi de jouer]
\begin{enumerate}
\item Énonce les six règles d'or d'une utilisation responsable de l'IA.
\item Pourquoi recopier un texte produit par une IA est-il considéré comme de la triche ?
\item Donne deux exemples de données personnelles qu'il ne faut jamais confier à une IA.
\end{enumerate}
\end{exercice}

\begin{corrige}[Exercice : À toi de jouer]
\begin{enumerate}
\item Vérifier les informations ; réfléchir d'abord soi-même ; ne pas tricher (citer ses sources) ; protéger ses données personnelles ; respecter les autres ; garder l'esprit critique.
\item Parce que l'élève présente comme sien un travail qu'il n'a pas fait : c'est du plagiat. De plus, il n'apprend rien, ce qui le pénalisera le jour de l'examen.
\item Par exemple : les mots de passe, l'adresse du domicile, les photos privées, les documents d'identité, le numéro de téléphone.
\end{enumerate}
\end{corrige}

\begin{aretenir}[Bilan de la section]
\begin{itemize}
\item Utiliser l'IA de manière \textbf{responsable} : vérifier, réfléchir d'abord soi-même, ne pas tricher, protéger ses données, respecter les autres, garder l'esprit critique.
\item L'IA doit être une \textbf{aide pour apprendre}, jamais un remplacement du travail personnel.
\item Le mauvais usage de l'IA (fausses images, fraude, atteinte à la vie privée) est \textbf{puni par la loi}.
\end{itemize}
\end{aretenir}

\coursec{Apprendre à échanger avec une IA}

Dernière étape de notre découverte : savoir \emph{dialoguer} avec une IA. Car un chatbot est un peu comme un moteur de recherche très bavard : \textbf{la qualité de sa réponse dépend beaucoup de la qualité de ta question}. Apprendre à bien poser ses questions, c'est une compétence qui te servira toute ta vie.

\textbf{Problématique :} comment formuler une demande claire à une IA pour obtenir une réponse utile et fiable ?

\courssub{Qu'est-ce qu'un prompt ?}

\begin{definition}[Prompt]
Un \cle{prompt} est le \textbf{message qu'on écrit à une IA} pour lui demander quelque chose : une question, une consigne, une demande d'explication ou de rédaction. Plus le prompt est \textbf{clair et précis}, plus la réponse est \textbf{utile}.
\end{definition}

Compare ces deux prompts envoyés au même chatbot :

\begin{itemize}
\item \textbf{Prompt vague} : « Parle-moi des plantes. » $\rightarrow$ réponse longue, générale, probablement inutilisable.
\item \textbf{Prompt précis} : « Explique à un élève de 3ème, en 5 lignes simples, le rôle de la photosynthèse chez les plantes, avec un exemple. » $\rightarrow$ réponse courte, adaptée à ton niveau, directement utile.
\end{itemize}

\courssub{Les ingrédients d'un bon prompt}

\begin{methode}[Construire un bon prompt en 4 ingrédients]
\begin{enumerate}
\item \textbf{Le contexte} : dis qui tu es ou pour quoi faire. Exemple : « Je suis élève en classe de 3ème et je prépare un exposé... »
\item \textbf{La tâche} : dis clairement ce que tu attends. Exemple : « ...explique-moi... », « ...propose-moi 5 questions de révision sur... », « ...corrige ce texte... »
\item \textbf{Le format} : précise la forme de la réponse. Exemple : « en 5 lignes », « sous forme de tableau », « avec un exemple ».
\item \textbf{Le niveau} : précise le niveau de langage. Exemple : « avec des mots simples », « niveau collège ».
\end{enumerate}
Ensuite, \textbf{dialogue} : si la réponse n'est pas satisfaisante, précise ta demande (« plus court », « donne un autre exemple », « je n'ai pas compris la deuxième phrase »). Et n'oublie jamais de \textbf{vérifier} les informations importantes !
\end{methode}

\begin{exemplebox}[Prompt faible contre prompt efficace]
\begin{center}
\begin{tabularx}{\linewidth}{|X|X|}
\hline
\rowcolor{popblueL} \textbf{Prompt faible} & \textbf{Prompt efficace} \\
\hline
« Devoirs de maths » & « Je suis en 3ème. Explique-moi comment résoudre l'équation $2x + 6 = 14$, étape par étape, avec des mots simples. » \\
\hline
« Exposé » & « Propose-moi un plan en 3 parties pour un exposé de 5 minutes sur les usages de l'IA au Cameroun, niveau collège. » \\
\hline
« Traduis » & « Traduis en anglais la phrase suivante : "L'intelligence artificielle transforme notre quotidien." » \\
\hline
\end{tabularx}
\end{center}
\end{exemplebox}

\begin{experience}[TP guidé : ma première conversation avec une IA]
\textbf{Objectif :} découvrir concrètement comment dialoguer avec un chatbot et constater l'importance d'un bon prompt.

\textbf{Matériel :} un ordinateur ou un téléphone connecté à Internet, un navigateur web, un chatbot gratuit accessible (par exemple un agent conversationnel d'IA grand public).

\textbf{Manipulation :}
\begin{enumerate}
\item Ouvre le chatbot dans le navigateur.
\item Envoie un \textbf{prompt vague} : « Parle-moi du Cameroun. » Note la réponse (longueur, précision, utilité).
\item Envoie maintenant un \textbf{prompt construit} avec les 4 ingrédients : « Je suis élève en classe de 3ème au Cameroun. Explique-moi en 5 lignes simples ce qu'est l'intelligence artificielle, avec un exemple de la vie quotidienne. » Note la réponse.
\item Pose une \textbf{question de suivi} : « Donne-moi un autre exemple, cette fois dans le domaine de l'agriculture. »
\item Demande à l'IA une \textbf{information vérifiable} (par exemple une date historique), puis vérifie-la dans ton manuel ou sur un site officiel.
\end{enumerate}

\textbf{Observations :} le prompt précis donne une réponse plus courte, plus claire et mieux adaptée à ton niveau. Le chatbot garde la mémoire de la conversation : tu peux enchaîner les questions. Certaines informations données peuvent être inexactes.

\textbf{Interprétation :} la qualité de la réponse dépend de la qualité du prompt ; le dialogue permet d'améliorer les réponses ; la vérification reste indispensable.
\end{experience}

\begin{attention}[Pièges à éviter quand on discute avec une IA]
\begin{itemize}
\item \textbf{Croire que l'IA « sait » tout} : elle calcule des réponses probables ; elle peut se tromper, surtout sur l'actualité récente ou les faits très précis.
\item \textbf{Confier ses secrets} : tout ce que tu écris peut être conservé. Jamais de mots de passe, d'adresses ou de photos privées.
\item \textbf{Accepter la première réponse sans réfléchir} : relis, critique, demande des précisions, vérifie.
\item \textbf{Oublier que c'est une machine} : l'IA n'a ni émotions ni conscience, même si elle écrit « je suis content de t'aider ».
\end{itemize}
\end{attention}

\begin{exoresolu}[Améliorer un prompt]
\textbf{Énoncé.} Rodrigue écrit à un chatbot : « La pollution ». Il obtient une réponse de trois pages, trop compliquée, qu'il ne peut pas utiliser pour son exposé de SVT. Réécris son prompt en appliquant les 4 ingrédients d'un bon prompt (contexte, tâche, format, niveau).
\tcblower
\textbf{Solution détaillée.} Un bon prompt possible :
\begin{itemize}
\item \textbf{Contexte} : « Je suis élève en classe de 3ème et je prépare un exposé de SVT. »
\item \textbf{Tâche} : « Explique-moi les causes de la pollution de l'air dans les grandes villes. »
\item \textbf{Format} : « Donne ta réponse en 5 lignes, sous forme de liste. »
\item \textbf{Niveau} : « Utilise des mots simples, niveau collège. »
\end{itemize}
Prompt complet : « Je suis élève en classe de 3ème et je prépare un exposé de SVT. Explique-moi les causes de la pollution de l'air dans les grandes villes, en 5 lignes, sous forme de liste, avec des mots simples de niveau collège. » La réponse sera courte, claire et directement utilisable.
\end{exoresolu}

\begin{exercice}[À toi de jouer]
\begin{enumerate}
\item Qu'est-ce qu'un prompt ? Pourquoi sa qualité est-elle importante ?
\item Cite les quatre ingrédients d'un bon prompt.
\item Améliore ce prompt faible : « L'informatique ». (Tu dois obtenir un prompt avec contexte, tâche, format et niveau.)
\item Pourquoi faut-il vérifier les informations données par un chatbot, même quand elles semblent bien écrites ?
\end{enumerate}
\end{exercice}

\begin{corrige}[Exercice : À toi de jouer]
\begin{enumerate}
\item Un prompt est le message qu'on écrit à une IA pour lui demander quelque chose. Sa qualité est importante car plus il est clair et précis, plus la réponse est utile et adaptée.
\item Le contexte (qui je suis, pour quoi faire), la tâche (ce que j'attends), le format (forme de la réponse), le niveau (langage adapté).
\item Exemple : « Je suis élève en classe de 3ème. Explique-moi en 5 lignes simples ce qu'est un système d'exploitation, avec un exemple. »
\item Parce qu'un chatbot calcule la réponse la plus probable, pas la plus vraie : il peut se tromper ou inventer des faits avec assurance. Il faut donc vérifier les informations importantes dans des sources fiables.
\end{enumerate}
\end{corrige}

\begin{aretenir}[Bilan de la section]
\begin{itemize}
\item Un \cle{prompt} est le message envoyé à une IA ; sa qualité détermine la qualité de la réponse.
\item Un bon prompt contient 4 ingrédients : \textbf{contexte, tâche, format, niveau}.
\item On peut \textbf{dialoguer} avec l'IA pour préciser et améliorer les réponses.
\item Règle finale : \textbf{vérifier} toujours les informations importantes et ne jamais confier de données personnelles.
\end{itemize}
\end{aretenir}

\begin{aretenir}[Bilan général de la leçon]
\begin{itemize}
\item L'\cle{Intelligence Artificielle} est un ensemble de techniques informatiques permettant à une machine de réaliser des tâches qui demandent habituellement de l'intelligence humaine.
\item Contrairement au programme classique (données + règles $\rightarrow$ résultat), l'IA \textbf{apprend à partir d'exemples} (données + résultats $\rightarrow$ règles découvertes) : c'est l'\cle{apprentissage automatique}.
\item Toutes les IA actuelles sont des \textbf{IA faibles}, spécialisées ; elles savent voir, écouter, parler, écrire et prédire.
\item L'IA est présente partout (éducation, santé, agriculture, transports, commerce, sécurité) : elle est rapide et disponible, mais elle peut \textbf{se tromper, être biaisée et menacer la vie privée}.
\item Utiliser l'IA de manière \textbf{responsable} : vérifier, réfléchir d'abord soi-même, ne pas tricher, protéger ses données, respecter les autres.
\item Pour bien échanger avec une IA, on rédige des \textbf{prompts clairs} (contexte, tâche, format, niveau) et on garde toujours l'\textbf{esprit critique}.
\end{itemize}
\end{aretenir}

\coursec{Importance et usages de l'IA dans la vie courante}

Après avoir découvert ce qu'est l'intelligence artificielle — une discipline qui permet aux machines d'imiter certaines fonctions cognitives humaines comme l'apprentissage, la reconnaissance de formes ou la prise de décision — nous allons maintenant explorer concrètement où elle se cache dans notre quotidien. L'IA n'est pas confinée aux laboratoires de recherche ou aux films de science-fiction : elle est déjà présente dans votre téléphone, dans les services bancaires que vous utilisez, dans les champs des agriculteurs de votre région, et même dans les hôpitaux qui vous soignent. Comprendre ces usages, c'est mieux saisir l'impact de cette technologie sur la société camerounaise et sur votre avenir professionnel.

\courssub{IA dans la communication : assistants, traduction et correction}

\begin{definition}[Assistant vocal]
Un \cle{assistant vocal} est une application d'IA capable de reconnaître la parole (\emph{reconnaissance vocale}), de comprendre l'intention de l'utilisateur (\emph{compréhension du langage naturel}) et d'exécuter une action ou de répondre par une voix de synthèse (\emph{synthèse vocale}). Exemples : Google Assistant, Siri, Alexa.
\end{definition}

\begin{exemplebox}[Traduction automatique au service du multilinguisme camerounais]
Le Cameroun compte plus de 250 langues locales. La traduction automatique permet de franchir la barrière linguistique. Des applications comme Google Translate ou Microsoft Translator intègrent des modèles multilingues capables de traduire directement entre de nombreuses langues (français, anglais, fulfulde, bassa, etc.) sans passer systématiquement par l'anglais comme langue pivot.
\end{exemplebox}

\begin{propriete}[Correcteur de texte intelligent]
Un correcteur moderne ne se contente pas de comparer chaque mot à un dictionnaire. Il utilise un \cle{modèle de langage} entraîné sur une grande quantité de textes pour proposer la correction la plus probable en fonction du contexte (accords, temps, prépositions), et pas seulement de l'orthographe du mot isolé.
\end{propriete}

\courssub{IA dans les réseaux sociaux et le divertissement : recommandation et création}

\begin{definition}[Système de recommandation]
Un \cle{système de recommandation} prédit les éléments (vidéos, musiques, produits) qu'un utilisateur appréciera, à partir de son historique et de celui d'utilisateurs ayant des goûts similaires. Deux approches principales sont combinées :
\begin{itemize}
    \item \textbf{Filtrage collaboratif} : « Les utilisateurs qui ont aimé A ont aussi aimé B ».
    \item \textbf{Filtrage basé sur le contenu} : analyse des caractéristiques de l'élément (mots-clés, genre, thème) pour trouver des articles ressemblants.
\end{itemize}
Les plateformes (YouTube, TikTok, Spotify, Netflix) utilisent ces approches pour personnaliser votre fil d'actualité.
\end{definition}

\begin{exemplebox}[Filtres photo et réalité augmentée]
Quand vous appliquez un filtre « oreilles de chat » ou « rajeunissement » sur WhatsApp, Instagram ou Snapchat, une IA détecte en temps réel les \cle{points clés du visage} (yeux, nez, bouche, contour). Un masque numérique est ensuite déformé pour suivre vos expressions. Tout cela fonctionne directement sur votre téléphone sans envoyer la vidéo sur Internet.
\end{exemplebox}

\courssub{IA dans la santé : diagnostic, imagerie et prévention}

\begin{definition}[IA d'aide au diagnostic médical]
L'IA ne remplace pas le médecin ; elle agit comme un \cle{second lecteur} ou un \cle{outil d'aide à la décision}. Pour l'analyse d'images médicales (radiographies, IRM, scanner, fonds d'œil), l'IA aide à repérer des anomalies : nodules pulmonaires, signes de rétinopathie diabétique, contours de tumeurs. Le médecin valide toujours le diagnostic final.
\end{definition}

\begin{exemplebox}[Rappels de vaccination par SMS et chatbot]
Au Cameroun, le Programme Élargi de Vaccination (PEV) expérimente des chatbots WhatsApp et des envois SMS automatisés. Un système couplé à un modèle de prédiction de \cle{risque de perte de vue} (basé sur l'historique, la zone géographique, l'âge de l'enfant) priorise les rappels pour les familles les plus à risque de manquer une dose. L'IA optimise le ciblage et le contenu du message (langue, ton, horaire) pour maximiser le taux de réponse.
\end{exemplebox}

\courssub{IA dans l'agriculture camerounaise : météo, détection de maladies et conseil}

L'agriculture emploie plus de 60 \% de la population active au Cameroun. L'IA y apporte des solutions concrètes, souvent accessibles via un simple téléphone (USSD, SMS) ou une application légère.

\begin{exemplebox}[Détection des maladies du manioc par photo -- \emph{PlantVillage Nuru}]
Un paysan photographie une feuille de manioc suspecte. L'application \texttt{Nuru} utilise une IA entraînée sur des dizaines de milliers d'images annotées (mosaïque du manioc, pourriture brune, bactériose, plante saine). Le modèle tourne \emph{hors ligne} sur le téléphone. Résultat en quelques secondes : diagnostic + conseil de lutte (variétés résistantes, élimination des plants atteints, rotation culturale). Précision terrain : 90-95 \% sur les classes principales.
\end{exemplebox}

\begin{propriete}[Prévision météo hyperlocale pour les semis]
Les modèles météorologiques classiques ont une résolution de 10-25 km. Des startups (ex. \texttt{Ignitia}, \texttt{Tomorrow.io}) utilisent de l'IA pour produire des prévisions à 1-3 km, mises à jour toutes les heures. Le paysan reçoit par SMS : « Pluie attendue demain 14 h-18 h sur votre zone (Bafia). Semis recommandé ce matin. » Cela réduit les pertes dues aux faux départs de saison.
\end{propriete}

\courssub{IA dans la finance : lutte contre la fraude Mobile Money et chatbots bancaires}

Le Mobile Money (MTN MoMo, Orange Money) traite des millions de transactions quotidiennes au Cameroun. La fraude (arnaques, usurpation d'identité, blanchiment) est un défi majeur.

\begin{exemplebox}[Détection de fraude en temps réel sur MTN MoMo]
Quand vous initiez un transfert, une IA analyse la transaction en moins de 100 ms :
\begin{enumerate}
    \item \textbf{Analyse du contexte} : montant, heure, type d'opération, historique de l'expéditeur et du destinataire, localisation via la borne GSM, nombre de transactions récentes, type de téléphone.
    \item \textbf{Score de risque} : l'IA attribue un score entre 0 et 1 (0 = sûr, 1 = fraude probable) en comparant à des millions d'exemples passés.
    \item \textbf{Décision automatique} :
    \begin{itemize}
        \item Score faible : validation immédiate.
        \item Score moyen : demande de confirmation (code OTP, PIN, empreinte).
        \item Score élevé : blocage de la transaction et alerte de l'équipe de sécurité.
    \end{itemize}
\end{enumerate}
Le système est mis à jour très régulièrement pour s'adapter aux nouvelles techniques des fraudeurs.
\end{exemplebox}

\begin{definition}[Chatbot bancaire]
Un \cle{chatbot} bancaire (ex. \texttt{UBA Leo}, \texttt{Ecobank Ada}) combine :
\begin{itemize}
    \item un moteur de compréhension du langage pour identifier votre intention (solde, virement, recharge) ;
    \item une base de connaissances (FAQ, procédures) ;
    \item une connexion sécurisée aux services de la banque pour exécuter l'opération.
\end{itemize}
Il transfère à un agent humain s'il ne comprend pas ou si vous le demandez.
\end{definition}

\courssub{IA dans les transports et la sécurité : GPS, trafic, reconnaissance faciale}

\begin{propriete}[Estimation du trafic et ETA (Estimated Time of Arrival)]
Google Maps, Waze ou des applications locales agrègent :
\begin{itemize}
    \item \textbf{Données de sondes} : positions GPS anonymisées des téléphones (vitesse, densité).
    \item \textbf{Données officielles} : capteurs routiers, caméras, travaux signalés.
    \item \textbf{Modèle de prédiction} : l'IA modélise le réseau routier comme un graphe et propage l'information de congestion pour prédire la vitesse sur chaque tronçon pour l'heure à venir.
\end{itemize}
L'ETA affichée est la somme des temps de parcours prédits sur l'itinéraire le plus rapide.
\end{propriete}

\begin{attention}[Reconnaissance faciale et vidéosurveillance : distinction technique et éthique]
Techniquement, la reconnaissance faciale repose sur deux étapes distinctes :
\begin{enumerate}
    \item \textbf{Détection} : localiser \emph{s'il y a} un visage dans l'image (cadrer le visage).
    \item \textbf{Reconnaissance (identification/vérification)} : déterminer \emph{de qui} il s'agit en comparant les traits du visage détecté à une base de données de visages connus.
\end{enumerate}
\textbf{Piège fréquent} : confondre \emph{détection} (présence d'un visage) et \emph{reconnaissance} (identité de la personne). Éthiquement, le déploiement en espace public soulève des questions de consentement, de biais (erreurs plus fréquentes selon la couleur de peau, l'âge, le genre) et de surveillance de masse. Le Cameroun encadre partiellement ces usages via la loi n°2010/012 sur la cybercriminalité et la loi sur la protection des données personnelles.
\end{attention}

\courssub{IA dans l'éducation : tuteurs, révision et aide à la rédaction}

\begin{definition}[Tuteur intelligent (ITS -- Intelligent Tutoring System)]
Un \cle{tuteur intelligent} adapte les exercices proposés à l'élève en estimant :
\begin{itemize}
    \item ce que l'élève sait déjà (\emph{modèle de l'élève}) ;
    \item la difficulté et les prérequis de chaque exercice (\emph{modèle de l'expert}) ;
    \item la meilleure stratégie pédagogique : réviser un prérequis, proposer un indice, passer au niveau supérieur (\emph{modèle pédagogique}).
\end{itemize}
Exemples accessibles au Cameroun : \texttt{Khan Academy} (exercices adaptatifs), \texttt{Duolingo} (langues), \texttt{PhET} (simulations sciences), applications locales (\texttt{Educamer}, \texttt{Schoolap}).
\end{definition}

\begin{exemplebox}[Outil d'aide à la rédaction : au-delà de la correction]
Des assistants d'écriture (ex. \texttt{LanguageTool}, \texttt{Grammarly}, \texttt{DeepL Write}, ou intégrés dans Word/Google Docs) proposent maintenant de la \cle{régénération de texte} : reformulation plus formelle, plus concise, changement de ton, résumé, génération de plan. Ils s'appuient sur des \cle{grands modèles de langage} (LLM). \textbf{Usage responsable} : l'élève doit comprendre la suggestion, l'accepter ou la rejeter en connaissance de cause, et citer l'outil si le texte final est soumis comme production personnelle.
\end{exemplebox}

\courssub{Importance économique : nouveaux métiers et transformation de l'existant}

L'IA ne détruit pas seulement des tâches ; elle crée des métiers et en transforme d'autres. Au Cameroun, la transition est en marche.

\begin{center}
\begin{tabularx}{\linewidth}{|l|X|X|}
\hline
\rowcolor{popblueL}
\textbf{Métier émergent / transformé} & \textbf{Rôle de l'IA} & \textbf{Compétences clés (niveau entrée)} \\
\hline
Développeur IA / Data Scientist & Concevoir, entraîner, utiliser des modèles & Python, bases de données, statistiques, logique algorithmique \\
Data Analyst / Analyste de données & Extraire du sens des données pour décider & SQL, Excel/Tableur avancé, visualisation (graphiques), esprit critique \\
Data Engineer & Préparer et nettoyer les données pour l'IA & SQL, bases de données, rigueur, automatisation scripts \\
Prompt Engineer / Concepteur d'interactions IA & Rédiger les instructions optimales pour l'IA & Expression écrite claire, compréhension des limites de l'IA, tests \\
Agriculteur \emph{augmenté} & Utiliser apps météo, diagnostic maladies, prix marchés & Littératie numérique (smartphone/USSD), interprétation d'alertes \\
Agent de santé communautaire & S'appuyer sur apps diagnostic/triage, rappels vaccins & Formation apps mobiles, relationnel, éthique, escalade vers infirmier/médecin \\
Comptable / Gestionnaire & Automatisation saisie, détection anomalies, prévision trésorerie & Maîtrise tableur (Excel/LibreOffice), notions comptables, esprit critique \\
Enseignant & Différenciation pédagogique, génération ressources, détection plagiat & Pédagogie, culture numérique, esprit critique, curation ressources IA \\
\hline
\end{tabularx}
\end{center}

\begin{savaistu}[Silicon Mountain et l'écosystème IA camerounais]
Buea, surnommée \emph{Silicon Mountain}, concentre des startups IA : \texttt{Agrix Tech} (diagnostic cultures par photo + conseil SMS), \texttt{GiftedMom} / \texttt{MomConnect} (suivi grossesse/vaccination par IA), \texttt{KmerPad} (santé menstruelle), \texttt{Djaa} (logistique optimisée par IA). À Yaoundé, l'écosystème \texttt{JFN Center}, \texttt{ActivSpaces}, \texttt{Orange Digital Center} et les universités (UY1, ENSP, IUT) forment la relève. Le MINPOSTEL pilote la \emph{Stratégie Nationale Intelligence Artificielle} (SNIA 2024-2030) pour structurer la filière : formation, recherche, infrastructure, cadre éthique et réglementaire.
\end{savaistu}

\courssub{Synthèse visuelle : l'IA au quotidien}

\begin{popfigure}
\begin{tikzpicture}[
    mindmap,
    concept color=popblue,
    level 1 concept/.append style={level distance=4.5cm, sibling angle=60},
    level 2 concept/.append style={level distance=3.2cm, sibling angle=45},
    every node/.style={font=\small},
    text=popdark
]
\node[concept, scale=1.2] {IA au quotidien}
    child[concept color=poporange] { node[concept] {Communication}
        child { node[concept, scale=0.7, align=center]{Assistants vocaux\\ (Google, Siri)} }
        child { node[concept, scale=0.7, align=center]{Traduction auto.\\ (français $\leftrightarrow$ langues locales)} }
    }
    child[concept color=popgreen] { node[concept] {Santé}
        child { node[concept, scale=0.7, align=center]{Aide diagnostic\\ imagerie médicale} }
        child { node[concept, scale=0.7, align=center]{Rappels vaccin\\ chatbot/SMS} }
    }
    child[concept color=poppurple] { node[concept] {Agriculture}
        child { node[concept, scale=0.7, align=center]{Prévision météo\\ hyperlocale semis} }
        child { node[concept, scale=0.7, align=center]{Détection maladies\\ plantes par photo} }
    }
    child[concept color=poppink] { node[concept] {Finance}
        child { node[concept, scale=0.7, align=center]{Anti-fraude\\ Mobile Money} }
        child { node[concept, scale=0.7, align=center]{Chatbots\\ banques} }
    }
    child[concept color=popgold] { node[concept] {Transports}
        child { node[concept, scale=0.7, align=center]{GPS + trafic\\ temps réel} }
        child { node[concept, scale=0.7, align=center]{Reconnaissance\\ faciale / vidéo} }
    }
    child[concept color=popblue] { node[concept] {Éducation}
        child { node[concept, scale=0.7, align=center]{Tuteurs intelligents\\ révision adaptative} }
        child { node[concept, scale=0.7, align=center]{Aide rédaction\\ résumé, reformulation} }
    };
\end{tikzpicture}
\legende{Carte mentale : six domaines majeurs de l'IA au quotidien avec deux exemples concrets par domaine.}
\end{popfigure}

\courssub{Tableau récapitulatif : domaine, application, capacité d'IA mobilisée}

\begin{center}
\begin{tabularx}{\linewidth}{|l|X|X|}
\hline
\rowcolor{popblueL}
\textbf{Domaine} & \textbf{Exemple d'application} & \textbf{Capacité d'IA mobilisée} \\
\hline
Communication & Assistant vocal (Google Assistant) & \textbf{Écouter} (reconnaissance vocale), \textbf{Dialoguer} (compréhension + synthèse) \\
Communication & Traduction automatique (Google Translate) & \textbf{Lire/Écrire} (traduction multilingue) \\
Communication & Correcteur intelligent (LanguageTool) & \textbf{Lire/Écrire} (modèle de langage contextuel) \\
\hline
Réseaux sociaux & Recommandation vidéo (TikTok/YouTube) & \textbf{Prédire} (goûts utilisateur) \\
Réseaux sociaux & Filtre AR visage (Snapchat/Instagram) & \textbf{Voir} (détection points clés visage temps réel) \\
\hline
Santé & Détection rétinopathie diabétique & \textbf{Voir} (analyse image médicale) \\
Santé & Chatbot rappels vaccination & \textbf{Dialoguer} (compréhension), \textbf{Prédire} (risque oubli) \\
\hline
Agriculture & Diagnostic manioc par photo (Nuru) & \textbf{Voir} (classification image feuille) \\
Agriculture & Prévision météo hyperlocale (Ignitia) & \textbf{Prédire} (météo locale) \\
\hline
Finance & Anti-fraude MoMo (Scoring temps réel) & \textbf{Prédire} (risque fraude), \textbf{Décider} (blocage/validation) \\
Finance & Chatbot banque (UBA Leo) & \textbf{Dialoguer} (compréhension + action) \\
\hline
Transports & ETA Google Maps / Waze & \textbf{Prédire} (trafic), \textbf{Optimiser} (itinéraire) \\
Transports & Vidéosurveillance faciale & \textbf{Voir} (détection), \textbf{Identifier} (reconnaissance) \\
\hline
Éducation & Tuteur adaptatif (Khan Academy) & \textbf{Modéliser l'élève} (niveau), \textbf{Décider} (exercice suivant) \\
Éducation & Assistant rédaction (LLM) & \textbf{Lire/Écrire} (génération, résumé, reformulation) \\
\hline
\end{tabularx}
\end{center}

\begin{methode}[Repérer une IA autour de soi -- 3 questions clés]
Face à un service numérique, posez-vous ces trois questions pour savoir s'il repose sur de l'IA :
\begin{enumerate}
    \item \textbf{La machine apprend-elle de données ?} Y a-t-il un jeu d'exemples sur lequel un modèle a été \emph{entraîné} pour ajuster son fonctionnement ? (Oui = apprentissage automatique / \emph{machine learning}.)
    \item \textbf{S'adapte-t-elle à l'utilisateur ou au contexte ?} Le comportement change-t-il selon votre historique, votre localisation, l'heure, vos préférences ? (Oui = personnalisation / adaptation contextuelle.)
    \item \textbf{Prend-elle des décisions sans intervention humaine directe ?} Le système exécute-t-il une action (blocage transaction, suggestion vidéo, diagnostic, alerte) automatiquement, sans qu'un humain valide chaque cas ? (Oui = autonomie décisionnelle bornée.)
\end{enumerate}
Si vous répondez \textbf{oui aux trois}, vous avez très probablement affaire à un système d'IA. Si c'est non à la première (règles fixes codées par un humain, ex. tableur, calculatrice, jeu vidéo classique), ce n'est \textbf{pas} de l'IA au sens moderne.
\end{methode}

\begin{attention}[Erreur fréquente : tout programme n'est pas une IA]
\begin{itemize}
    \item Un \textbf{tableur} (Excel, LibreOffice Calc) exécute des formules déterministes définies par l'utilisateur. Il n'apprend pas de données.
    \item Un \textbf{jeu vidéo classique} (ex. Mario, FIFA en mode hors ligne) suit des scripts et des arbres de décision précodés. Les PNJ n'apprennent pas de vos parties.
    \item Un \textbf{script de sauvegarde automatique} copie des fichiers selon un planning fixe. Aucune adaptation, aucun apprentissage.
    \item \textbf{Critère discriminant} : la présence d'une phase d'\emph{entraînement} sur des données pour optimiser un résultat, et la capacité de \emph{généraliser} à des exemples jamais vus.
\end{itemize}
\end{attention}

\begin{exoresolu}[Identifier la capacité d'IA dans un service courant]
\textbf{Énoncé.} Pour chacun des services ci-dessous, indiquez la capacité d'IA principale mobilisée (\textbf{Voir}, \textbf{Écouter}, \textbf{Lire/Écrire}, \textbf{Prédire}, \textbf{Dialoguer}, \textbf{Décider}, \textbf{Modéliser}) et justifiez en une phrase.

\begin{enumerate}
    \item L'application \texttt{Google Lens} identifie une plante à partir d'une photo.
    \item \texttt{WhatsApp} transcrit un message vocal en texte.
    \item \texttt{Netflix} vous suggère « \emph{Parce que vous avez regardé...} ».
    \item Le chatbot de votre banque exécute un virement après votre confirmation.
    \item \texttt{Duolingo} choisit l'exercice suivant selon vos erreurs récentes.
\end{enumerate}

\textbf{Solution détaillée.}
\begin{enumerate}
    \item \textbf{Voir} -- Classification d'image (IA visuelle) sur la photo de la plante.
    \item \textbf{Écouter} -- Reconnaissance vocale convertissant l'audio en texte.
    \item \textbf{Prédire} -- Estimation de la probabilité que vous aimiez la vidéo selon votre historique.
    \item \textbf{Dialoguer} + \textbf{Décider} -- Compréhension de l'intention « virement » + exécution de l'action via l'API bancaire.
    \item \textbf{Modéliser l'élève} + \textbf{Décider} -- Estimation de votre niveau (modèle de l'élève) ; choix de l'exercice le plus utile (modèle pédagogique).
\end{enumerate}
\end{exoresolu}

\begin{exercice}[Observer l'IA dans votre environnement]
\begin{enumerate}
    \item Ouvrez votre téléphone. Identifiez \textbf{trois} applications que vous avez utilisées cette semaine et qui intègrent de l'IA. Pour chacune, nommez la capacité principale (Voir, Écouter, Lire/Écrire, Prédire, Dialoguer, Décider, Modéliser).
    \item Interrogez un commerçant, un agriculteur ou un agent de santé de votre entourage : utilise-t-il un service numérique qui « apprend » ou « s'adapte » ? Lequel ? Quelle capacité d'IA ?
    \item Consultez les offres d'emploi sur \texttt{Emploi.cm} ou \texttt{LinkedIn Cameroun} avec les mots-clés « IA », « Data », « Machine Learning ». Notez trois métiers demandés et une compétence technique requise pour chacun.
\end{enumerate}
\end{exercice}

\begin{corrige}[Exercice n°1 -- Observer l'IA dans votre environnement]
\textit{(Corrigé indicatif -- les réponses varient selon l'élève.)}
\begin{enumerate}
    \item Exemples typiques : \texttt{Google Maps} (Prédire trafic + Optimiser itinéraire), \texttt{Instagram/TikTok} (Prédire recommandation), \texttt{Clavier Google/Gboard} (Prédire mot suivant + Corriger), \texttt{Google Photos} (Voir : reconnaissance visages/objets), \texttt{Assistant Google} (Écouter + Dialoguer).
    \item Exemple agriculteur : application \texttt{Plantix} ou \texttt{Nuru} (Voir : diagnostic feuille) ; prévisions météo \texttt{Ignitia} (Prédire). Exemple agent de santé : \texttt{DHIS2} avec module prédictif (Prédire : ruptures stock, épidémies) ; chatbot PEV (Dialoguer).
    \item Métiers vus récemment : \textit{Data Analyst} (SQL, Python, Power BI), \textit{Junior ML Engineer} (Python, PyTorch, MLOps), \textit{Business Intelligence Developer} (SQL, modélisation, Tableau/Power BI), \textit{Prompt Engineer / AI Content Specialist} (LLM, évaluation, rédaction).
\end{enumerate}
\end{corrige}

\coursec{Avantages et inconvénients de l'IA}

Après avoir découvert l'importance et les usages de l'IA dans notre vie quotidienne — du téléphone intelligent qui corrige nos fautes aux systèmes de recommandation qui nous suggèrent le prochain film à regarder — il est temps de prendre du recul. Tout outil puissant a ses forces et ses limites ; l'IA ne fait pas exception. Comprendre ses avantages nous aide à l'exploiter intelligemment, tandis que connaître ses inconvénients nous protège d'une confiance aveugle. Dans cette section, nous dressons un bilan honnête et documenté, comme le ferait un ingénieur avant de choisir une technologie pour un projet réel.

\courssub{Les principaux avantages de l'IA}

L'IA excelle là où l'humain se fatigue, s'ennuie ou manque de temps. Trois atouts majeurs reviennent sans cesse dans les retours d'expérience, du cybercafé de quartier au centre hospitalier de référence.

\begin{definition}[Rapidité de traitement de grandes quantités de données]
La \cle{vitesse de calcul} des processeurs modernes (CPU, GPU, accélérateurs spécialisés) permet à un modèle d'IA d'analyser en quelques secondes des volumes de données qu'un être humain mettrait des jours, voire des mois, à parcourir. Cette rapidité ne se limite pas au calcul brut : elle s'applique à la recherche d'information, à la classification d'images, à la traduction automatique ou à la détection d'anomalies dans des flux continus (trafic réseau, capteurs agricoles, transactions bancaires).
\end{definition}

\begin{definition}[Disponibilité 24h/24 et 7j/7]
Contrairement à un opérateur humain qui a besoin de sommeil, de pauses et de congés, un service piloté par IA — chatbot d'assistance, système de surveillance vidéo, serveur de recommandation — fonctionne en continu sans interruption. Pour une administration en ligne (ex. portail \texttt{camtelecom.cm} pour le suivi de dossiers) ou une plateforme d'apprentissage accessible aux élèves du Nord comme du Sud du pays, cette permanence est un facteur d'équité territoriale.
\end{definition}

\begin{definition}[Réduction des erreurs dans les tâches répétitives]
L'humain commet des erreurs de fatigue ou d'inattention lorsqu'il répète mille fois le même geste (saisie de factures, tri de courriels, contrôle visuel de pièces sur une chaîne). L'IA, une fois correctement entraînée et validée, exécute la tâche avec une \cle{constance} quasi parfaite. Cela ne signifie pas « zéro erreur » (voir les inconvénients), mais un taux d'erreur systématiquement plus bas sur les tâches bien définies, structurées et répétitives.
\end{definition}

\begin{exemplebox}[Exemple chiffré : gain de temps sur une tâche administrative]
Imaginons le secrétariat d'un lycée qui doit classer 500 bulletins de notes par trimestre.
\begin{itemize}
    \item \textbf{Humain seul :} lecture, vérification, rangement $\approx$ 10 minutes par bulletin $\Rightarrow$ $500 \times 10\ \text{min} = 5000\ \text{min} \approx 83\ \text{heures}$ (plus de deux semaines de travail à temps plein).
    \item \textbf{Avec numérisation (OCR) + IA de classification :} numérisation par lots, extraction automatique des champs (nom, classe, moyennes), tri numérique $\approx$ 2 secondes par bulletin $\Rightarrow$ $500 \times 2\ \text{s} = 1000\ \text{s} \approx 17\ \text{minutes}$.
\end{itemize}
Le gain est d'un facteur $\approx 300$. Le secrétariat peut réinvestir ce temps dans l'accompagnement des élèves.
\tcblower
\textbf{Attention :} ce chiffre suppose une numérisation de qualité et un modèle bien entraîné sur des bulletins camerounais. Des bulletins manuscrits illisibles ou un modèle non adapté feront chuter la performance.
\end{exemplebox}

\courssub{IA et inclusion : aide aux personnes handicapées, accès à l'éducation et à l'information}

L'IA n'est pas qu'un accélérateur de productivité ; c'est aussi un levier d'accessibilité.

\begin{itemize}
    \item \textbf{Synthèse vocale (Text-to-Speech) :} un élève malvoyant peut faire lire tout manuel numérisé, toute page web, tout corrigé d'examen. Les voix neuronales actuelles sont naturelles, gèrent la ponctuation et l'intonation.
    \item \textbf{Reconnaissance vocale (Speech-to-Text) :} un élève à mobilité réduite dicte ses rédactions, ses commandes au clavier, ses recherches. Les modèles récents atteignent des taux d'erreur de mot (WER) inférieurs à \SI{5}{\percent} en français standard.
    \item \textbf{Sous-titrage et traduction en temps réel :} en classe, un élève sourd suit le cours grâce aux sous-titres générés par IA. Un élève anglophone dans un établissement francophone (et inversement) bénéficie d'une traduction simultanée.
    \item \textbf{Accès à l'éducation en zone rurale :} une application d'apprentissage adaptatif (type \texttt{Khan Academy} avec recommandations IA) fonctionne sur un smartphone bas de gamme, hors connexion après téléchargement initial. Elle comble l'absence d'enseignants spécialisés dans certaines matières.
\end{itemize}

\begin{savaistu}[L'IA au service de l'agriculture camerounaise]
Dans la région de l'Extrême-Nord, des capteurs peu coûteux (humidité du sol, température, pluviométrie) envoient leurs données vers un modèle IA hébergé sur un petit serveur local (type Raspberry Pi avec accélérateur). Le modèle prédit le moment optimal d'irrigation et alerte l'agriculteur par SMS (USSD, pas besoin de 4G). Résultats observés : économie d'eau de \SI{30}{\percent} et hausse du rendement de \SI{15}{\percent} sur le mil et le sorgho. L'IA n'est pas réservée aux grandes fermes industrielles ; elle descend jusqu'au petit exploitant.
\end{savaistu}

\courssub{Progrès en médecine, agriculture et prévention des catastrophes}

Au-delà de l'école et du bureau, l'IA sauve des vies et des récoltes.

\begin{itemize}
    \item \textbf{Médecine :} détection précoce du paludisme ou de la tuberculose sur images microscopiques ; analyse d'ECG sur montre connectée pour alerter sur une arythmie ; planification de radiothérapie ciblée.
    \item \textbf{Agriculture :} détection de maladies du cacaoyer par drone + vision par ordinateur ; prévision des rendements par imagerie satellite + séries temporelles ; optimisation des intrants (engrais, pesticides) par agriculture de précision.
    \item \textbf{Prévention des catastrophes :} modélisation des crues du fleuve Wouri ou de la Bénoué à partir de données pluviométriques, topographiques et historiques ; alertes précoces envoyées par SMS aux populations riveraines ; détection d'incendies de brousse par satellites couplée à des modèles de propagation.
\end{itemize}

\begin{popfigure}
\begin{tikzpicture}[
    scale=0.9,
    every node/.style={font=\small},
    adv/.style={fill=popgreenL, draw=popgreen, rounded corners=4pt, inner sep=4pt, minimum width=2.2cm, text centered},
    disadv/.style={fill=poppinkL, draw=poppink, rounded corners=4pt, inner sep=4pt, minimum width=2.2cm, text centered},
    pivot/.style={fill=poporange, draw=popdark, minimum width=1.2cm, minimum height=1.2cm, rounded corners=2pt},
    beam/.style={draw=popdark, thick},
    label/.style={font=\bfseries\small, text=popdark}
]
% Base
\draw[popdark, thick] (0,-2.5) rectangle (14,-2.3);
% Pivot central
\node[pivot] (pivot) at (7,0) {IA};
% Poutres avantages (gauche)
\draw[beam] (pivot.south west) -- ++(-6.5,-2.2) node[adv, anchor=south east, align=center] (adv1){Rapidité \\ traitement données};
\draw[beam] (pivot.south west) -- ++(-5.5,-2.2) node[adv, anchor=south east, align=center] (adv2){Disponibilité \\ 24h/24};
\draw[beam] (pivot.south west) -- ++(-4.5,-2.2) node[adv, anchor=south east, align=center] (adv3){Moins d'erreurs \\ tâches répétitives};
\draw[beam] (pivot.south west) -- ++(-3.5,-2.2) node[adv, anchor=south east, align=center] (adv4){Inclusion \\ handicap};
\draw[beam] (pivot.south west) -- ++(-2.5,-2.2) node[adv, anchor=south east, align=center] (adv5){Médecine, \\ agro, prévention};

% Poutres inconvénients (droite)
\draw[beam] (pivot.south east) -- ++(6.5,-2.2) node[disadv, anchor=south west, align=center] (dis1){Risque chômage \\ métiers automatisables};
\draw[beam] (pivot.south east) -- ++(5.5,-2.2) node[disadv, anchor=south west, align=center] (dis2){Coût équipements \\ et énergie};
\draw[beam] (pivot.south east) -- ++(4.5,-2.2) node[disadv, anchor=south west, align=center] (dis3){Biais et erreurs \\ données mauvaises};
\draw[beam] (pivot.south east) -- ++(3.5,-2.2) node[disadv, anchor=south west, align=center] (dis4){Vie privée \\ deepfakes};
\draw[beam] (pivot.south east) -- ++(2.5,-2.2) node[disadv, anchor=south west, align=center] (dis5){Dépendance \\ perte esprit critique};

% Légendes groupes
\node[label, anchor=south] at ($(adv3.north)+(0,0.3)$) {AVANTAGES};
\node[label, anchor=south] at ($(dis3.north)+(0,0.3)$) {INCONVÉNIENTS};
% Flèche d'équilibre
\draw[<->, popblue, thick] (adv3.north east) -- ++(0,0.5) -| (dis3.north west) node[midway, above, font=\tiny, popblue] {Équilibre à trouver};
\end{tikzpicture}
\legende{Bilan visuel : l'IA comme une balance. Chaque plateau porte des arguments documentés ; le point d'équilibre dépend du contexte d'usage, de la qualité des données et de la gouvernance humaine.}
\end{popfigure}

\courssub{Les principaux inconvénients et risques de l'IA}

Aucune technologie n'est neutre. L'IA concentre des risques techniques, économiques, éthiques et sociaux qu'un citoyen éclairé — et a fortiori un futur développeur — doit connaître.

\courssub{Risque de chômage pour certains métiers automatisables}

L'automatisation par IA ne remplace pas « le travail » en général, mais des \cle{tâches} précises : saisie, tri, traduction basique, génération de code standard, réponse aux FAQ de premier niveau. Les métiers dont la valeur ajoutée repose quasi exclusivement sur ces tâches (secrétariat d'exécution pur, téléopérateur de niveau 1, correcteur de copies à choix multiples) sont les plus exposés. En revanche, les métiers nécessitant de l'empathie, du jugement complexe, de la créativité, de la gestion d'imprévus physiques (plombier, infirmier, éducateur) le sont beaucoup moins.

\begin{attention}[Ne pas confondre « métier » et « tâche »]
Un comptable ne disparaît pas parce qu'un logiciel fait la saisie ; il évolue vers l'analyse, le conseil fiscal, l'audit. L'histoire le montre : la machine à écrire n'a pas éliminé les secrétaires, elle a transformé leur métier. L'enjeu pour toi, élève de 3ème, est d'acquérir des \textbf{compétences complémentaires à l'IA} (pensée critique, communication, résolution de problèmes nouveaux) plutôt que de craindre l'outil.
\end{attention}

\courssub{Coût des équipements et de l'énergie}

Entraîner un grand modèle de langage demande des milliers de processeurs graphiques pendant des semaines : coût estimé à plusieurs millions de dollars, empreinte carbone équivalente à plusieurs centaines de vols Paris–Yaoundé. Même l'\cle{inférence} (utilisation du modèle) consomme de l'énergie. Pour un établissement scolaire camerounais, faire tourner un modèle localement demande un ordinateur performant ou un accélérateur dédié, soit un investissement de 300 000 à 800 000 FCFA. La solution pragmatique : utiliser des API cloud (payantes à l'usage) ou des modèles légers optimisés pour tourner sur processeur classique (CPU).

\courssub{Erreurs, biais et « hallucinations » : quand les données trompent l'IA}

C'est le point technique le plus crucial. Une IA \textbf{n'a pas de compréhension du monde} ; elle apprend des corrélations statistiques dans ses données d'entraînement. Si ces données sont biaisées, l'IA reproduira, voire amplifiera, les biais.

\begin{propriete}[Biais de données et reproduction des préjugés]
\textbf{Énoncé :} Un modèle d'IA entraîné sur un corpus de données non représentatif ou porteur de préjugés historiques produira des sorties systématiquement injustes ou fausses pour les groupes sous-représentés ou stigmatisés dans ce corpus.
\textbf{Conditions :} Données d'entraînement déséquilibrées (ex. 90\% d'images de personnes à peau claire), étiquettes subjectives (ex. « profession prestigieuse » annoté par des humains biaisés), absence de contre-mesures (rééquilibrage, contraintes d'équité).
\textbf{Conséquence :} Reconnaissance faciale moins performante sur peaux foncées ; outil de recrutement qui pénalise les CV mentionnant des universités africaines ; générateur d'images qui associe « médecin » à un homme blanc et « infirmière » à une femme.
\end{propriete}

\begin{exemplebox}[Hallucination d'un chatbot : l'assurance du faux]
Tu demandes à un chatbot : « Donne-moi la référence exacte de l'article du Code du travail camerounais qui fixe la durée du congé maternité. »
Il répond avec assurance : « Article 132 du Code du travail (Loi n°92/007 du 14 août 1992) : 14 semaines. »
\textbf{Vérification :} L'article 132 traite du repos hebdomadaire. Le congé maternité est à l'article 86 (14 semaines, dont 6 avant et 8 après l'accouchement).
\tcblower
Le chatbot a \textbf{halluciné} : il a produit une réponse plausible syntaxiquement, cohérente dans le style « article de loi », mais factuellement fausse. Il ne « sait » pas ; il \textbf{prédit le token suivant probable}. C'est pourquoi \textbf{toute information factuelle fournie par une IA doit être vérifiée sur une source officielle}.
\end{exemplebox}

\courssub{Atteinte à la vie privée et désinformation (deepfakes)}

\begin{itemize}
    \item \textbf{Collecte massive de données personnelles :} assistants vocaux, applications « gratuites », objets connectés aspirent localisation, contacts, biométrie, conversations. Ces données servent à ré-entraîner les modèles (amélioration du service) ou sont revendues à des courtiers en données. Le RGPD (Europe) et la loi camerounaise sur la protection des données personnelles (loi n°2024/014) imposent consentement, finalité, droit d'accès et d'effacement — mais l'application reste faible.
    \item \textbf{Deepfakes :} vidéos ou audios synthétiques où le visage/la voix d'une personne réelle est remplacé par celui d'une autre, de façon quasi indétectable à l'œil nu. Usages malveillants : arnaque au président (virement urgent), diffamation, manipulation électorale, harcèlement (revenge porn synthétique).
\end{itemize}

\begin{savaistu}[Deepfake : le réflexe « source et cohérence »]
Une vidéo circule sur WhatsApp : le Ministre de l'Éducation annonce l'annulation du BEPC 2025.
\begin{enumerate}
    \item \textbf{Source :} La vidéo vient-elle du compte officiel vérifié (badge bleu) du Ministère ou d'un transfert anonyme ?
    \item \textbf{Cohérence :} L'annonce est-elle reprise sur le site \texttt{mineSEC.gov.cm}, à la CRTV, dans la presse écrite ?
    \item \textbf{Détails techniques :} Clignement des yeux naturel ? Synchronisation labiale parfaite ? Arrière-plan cohérent ? (Les deepfakes actuels peinent encore sur les micro-mouvements du cou, la chevelure, les reflets dans les lunettes.)
\end{enumerate}
Si un seul critère faiblit : \textbf{ne pas partager, signaler}.
\end{savaistu}

\courssub{Dépendance excessive et perte d'esprit critique}

Le risque le plus insidieux pour un élève : s'habituer à « demander à l'IA » sans réfléchir. Copier-coller une rédaction générée par un chatbot pour un devoir de français, accepter un code C proposé par un assistant sans le lire, croire une explication historique sans la croiser avec le manuel. Cela atrophie :
\begin{itemize}
    \item la \cle{mémoire de travail} (on ne retient pas ce qu'on ne restructure pas soi-même) ;
    \item la \cle{capacité de détection d'erreur} (on ne repère pas l'absurde si on n'a pas construit le raisonnement) ;
    \item la \cle{confiance en sa propre pensée} (syndrome de l'imposteur : « de toute façon l'IA fait mieux »).
\end{itemize}

\begin{methode}[Construire un avis équilibré sur l'IA : tableau avantages/inconvénients]
Pour tout sujet controversé (IA, OGM, énergie nucléaire, réseaux sociaux), suis cette démarche en 4 étapes :
\begin{enumerate}
    \item \textbf{Lister} : sur une feuille, deux colonnes « Avantages » et « Inconvénients ». Note au moins 5 points par colonne, factuels, sourcés si possible (ex. « Gain de temps 300$\times$ sur tri bulletins — source : simulation secrétariat lycée »).
    \item \textbf{Pondérer} : attribue une importance (1 = mineur, 2 = modéré, 3 = majeur) à chaque point selon \emph{ton} contexte (élève, futur développeur, citoyen camerounais).
    \item \textbf{Analyser les interactions} : un avantage peut-il atténuer un inconvénient ? (Ex. : coût équipement élevé \textit{mais} mutualisable via un serveur d'établissement + modèles légers optimisés.)
    \item \textbf{Conclure par écrit} : rédige un paragraphe de 5–8 lignes qui tranche : « Dans le contexte X, l'IA est globalement bénéfique/risquée/à encadrer car [argument 1 pondéré] et [argument 2 pondéré]. Je resterai vigilant sur [point de vigilance]. »
\end{enumerate}
Applique cette méthode à l'IA dans l'éducation au Cameroun comme exercice de fin de section.
\end{methode}

\courssub{Illustration quantitative : gain de temps sur tâches répétitives}

Le diagramme ci-dessous compare le temps moyen pour traiter 100 unités d'une tâche répétitive (classification de courriels, saisie de factures, tri de photos) par un humain seul vs un humain assisté par IA (validation seulement). Les données sont pédagogiques, basées sur des études de terrain.

\begin{popfigure}
\begin{tikzpicture}
\begin{axis}[
    ybar,
    bar width=18pt,
    width=\linewidth,
    height=9cm,
    enlargelimits=0.15,
    ylabel={Temps (minutes)},
    symbolic x coords={Classification courriels, Saisie factures, Tri photos, Rédaction réponses types},
    xtick=data,
    x tick label style={rotate=30, anchor=east, font=\small},
    ymin=0,
    ymax=130,
    legend style={at={(0.5,1.02)}, anchor=south, legend columns=2, font=\small},
    nodes near coords,
    nodes near coords align={vertical},
    every node near coord/.append style={font=\tiny, rotate=90, anchor=west, yshift=2pt},
]
\addplot[fill=popblueL, draw=popblue] coordinates {
    (Classification courriels, 90)
    (Saisie factures, 120)
    (Tri photos, 75)
    (Rédaction réponses types, 60)
};
\addplot[fill=poporangeL, draw=poporange] coordinates {
    (Classification courriels, 12)
    (Saisie factures, 18)
    (Tri photos, 10)
    (Rédaction réponses types, 8)
};
\legend{Humain seul, Humain + IA (validation)}
\end{axis}
\end{tikzpicture}
\legende{Temps moyen pour traiter 100 unités. L'IA ne remplace pas l'humain : elle réduit la partie répétitive (facteur 6 à 10), laissant à l'humain la validation, les cas ambigus, la relation client. Source : données pédagogiques inspirées d'études réelles.}
\end{popfigure}

\courssub{Bilan raisonné : l'IA est un outil, sa valeur dépend de l'usage humain}

\begin{aretenir}[Synthèse : l'IA, ni magicienne, ni maléfique]
L'intelligence artificielle est un \textbf{outil puissant, à usage dual} (bénéfique ou nuisible selon l'intention et la compétence de l'utilisateur).
\begin{itemize}
    \item \textbf{Forces :} vitesse, endurance, constance, passage à l'échelle, accessibilité (handicap, zones isolées), percées scientifiques (médecine, climat, agriculture).
    \item \textbf{Faiblesses :} pas de compréhension sémantique, hallucinations, biais des données, opacité (boîte noire), coût énergétique, risque de dépendance cognitive.
    \item \textbf{Risques sociétaux :} chômage de transition, surveillance de masse, désinformation (deepfakes), concentration du pouvoir chez quelques acteurs (Big Tech).
    \item \textbf{Notre responsabilité :} former des humains \textbf{augmentés par l'IA, pas remplacés par elle}. Cela passe par : vérification systématique des sources, apprentissage des bases (algorithmique, données, éthique), esprit critique, usage créatif et solidaire.
\end{itemize}
\end{aretenir}

\begin{exoresolu}[Analyse critique d'un usage d'IA au lycée]
\textbf{Énoncé :} Le proviseur d'un lycée de Yaoundé souhaite déployer un chatbot IA pour répondre aux questions des parents d'élèves (inscription, calendrier, frais, résultats). Il achète une solution clé en main hébergée à l'étranger, sans clause de localisation des données. Les parents interagissent en français et en langues locales (ewondo, bassa, fulfulde). Le chatbot hallucine parfois sur les montants des frais. Un parent porte plainte pour fuite de son numéro de téléphone.
\begin{enumerate}
    \item Identifie 3 avantages et 3 inconvénients concrets de ce déploiement.
    \item Propose 3 mesures techniques ou organisationnelles pour corriger les inconvénients.
    \item Rédige une conclusion argumentée (5 lignes) sur la pertinence du projet.
\end{enumerate}
\tcblower
\textbf{Solution détaillée :}
\begin{enumerate}
    \item \textbf{Avantages :} (1) Disponibilité 24h/24 pour parents travailleurs (horaires décalés). (2) Réduction de la file d'attente au secrétariat (gain de temps pour le personnel). (3) Support multilingue potentiel (si le modèle gère les langues locales).
    \textbf{Inconvénients :} (1) \textbf{Hallucination sur les frais} : risque de conflit, perte de confiance, erreur financière pour les familles. (2) \textbf{Données personnelles hors territoire} : numéro de téléphone, nom de l'enfant, niveau scolaire stockés sur serveur étranger sans garantie conforme à la loi camerounaise. (3) \textbf{Exclusion numérique} : parents sans smartphone ni data (zone rurale, personnes âgées) n'ont pas accès au service.
    \item \textbf{Mesures correctives :}
    \begin{enumerate}
        \item \textbf{Base de connaissances locale (RAG) :} brancher le chatbot sur une base de documents officiels validés (PDF des frais signés par l'intendance, calendrier officiel MINESEC) : l'IA ne génère plus, elle cite la source. Supprime l'hallucination sur les faits codifiés.
        \item \textbf{Hébergement local ou souverain :} déployer un modèle ouvert léger sur un serveur de l'établissement (ou cloud africain) avec contrat de traitement de données conforme à la loi n°2024/014.
        \item \textbf{Canal alternatif inclusif :} maintenir un standard téléphonique humain (USSD/menu vocal) et un affichage physique au secrétariat pour les non-connectés.
    \end{enumerate}
    \item \textbf{Conclusion :} Le projet est \textbf{pertinent à condition d'être encadré}. L'avantage majeur (service continu, désengorgement) est réel, mais l'hallucination sur les frais et l'exfiltration de données sont inacceptables sans base de connaissances locale et hébergement souverain. L'inclusion impose un canal de secours non numérique. Avec ces trois garde-fous, le chatbot devient un service public numérique responsable, aligné sur la stratégie « Cameroun Numérique 2030 ».
\end{enumerate}
\end{exoresolu}

\begin{exercice}[À toi de jouer : ton tableau avantages/inconvénients]
Applique la \textbf{Méthode} ci-dessus au sujet : « L'utilisation de ChatGPT (ou tout LLM) pour faire mes devoirs de maths, physique et français en classe de 3ème. »
\begin{enumerate}
    \item Remplis le tableau (5 avantages, 5 inconvénients, pondération 1–3).
    \item Écris ta conclusion argumentée (5–8 lignes) en t'adressant à ton professeur principal.
\end{enumerate}
\begin{center}
\begin{tabularx}{\linewidth}{|l|X|X|c|}
\hline
\textbf{Côté} & \textbf{Point factuel} & \textbf{Justification / Source} & \textbf{Poids (1–3)} \\
\hline
Avantage 1 & & & \\
Avantage 2 & & & \\
Avantage 3 & & & \\
Avantage 4 & & & \\
Avantage 5 & & & \\
\hline
Inconvénient 1 & & & \\
Inconvénient 2 & & & \\
Inconvénient 3 & & & \\
Inconvénient 4 & & & \\
Inconvénient 5 & & & \\
\hline
\end{tabularx}
\end{center}
\vspace{0.5cm}
\textbf{Ma conclusion :} \dotfill
\end{exercice}

\begin{corrige}[Exercice : tableau avantages/inconvénients devoirs + LLM]
\textit{Exemple de réponse d'élève rigoureux (la tienne peut différer tant que les arguments sont factuels et pondérés).}
\begin{center}
\begin{tabularx}{\linewidth}{|l|X|X|c|}
\hline
\textbf{Côté} & \textbf{Point factuel} & \textbf{Justification / Source} & \textbf{Poids} \\
\hline
Avantage 1 & Explication immédiate d'une notion bloquante & L'IA reformule le cours avec des exemples (ex. second degré : discriminant) & 3 \\
Avantage 2 & Génération d'exercices d'entraînement corrigés & « Donne-moi 5 équations du 2nd degré avec solutions détaillées » & 2 \\
Avantage 3 & Aide à la rédaction (plan, connecteurs, vocabulaire) & Pour la dissertation de français, l'IA propose un plan dialectique & 2 \\
Avantage 4 & Disponibilité le soir/week-end sans attendre le prof & Pas de créneau de permanence nécessaire & 1 \\
Avantage 5 & Traduction d'énoncés anglais/français (option LV2) & Utile pour les élèves bilingues ou en section européenne & 1 \\
\hline
Inconvénient 1 & Hallucination mathématique (faux discriminant, signe inversé) & Constaté sur modèles gratuits : erreurs sur $\Delta = b^2-4ac$ si $a,b,c$ négatifs & 3 \\
Inconvénient 2 & Perte d'apprentissage par copier-coller sans compréhension & Pédagogie : on retient ce qu'on restructure soi-même (effet de génération) & 3 \\
Inconvénient 3 & Détection par le prof (style trop lisse, connaissances hors programme) & Risque de note 0 + sanction disciplinaire (règlement intérieur) & 2 \\
Inconvénient 4 & Dépendance : angoisse face à un devoir sans connexion/IA & Perte d'autonomie, syndrome de la page blanche sans béquille & 2 \\
Inconvénient 5 & Données personnelles (devoirs, niveau) envoyées à serveur tiers & Vie privée : les conversations sont conservées par défaut par les fournisseurs & 1 \\
\hline
\end{tabularx}
\end{center}
\textbf{Conclusion :} J'utilise l'IA comme \textbf{tuteur explicatif} (je lui demande « explique-moi pourquoi $\Delta$ détermine le nombre de racines » et je reformule à ma façon), jamais comme \textbf{producteur de réponses finales}. Je vérifie chaque calcul sur ma calculatrice ou avec le prof. Ainsi, je garde l'avantage (compréhension accélérée) sans les inconvénients majeurs (hallucination, tricherie, dépendance). C'est un usage \textbf{responsable et augmentant}.
\end{corrige}

Cette section t'a donné les clés pour évaluer l'IA sans naïveté ni technophobie. La suivante — « Utilisation responsable de l'IA » — transformera ce bilan en règles d'action concrètes : charte d'usage, citation des sources, protection des données, et premiers pas pour \emph{créer} tes propres petits modèles au lieu de seulement les subir.

\coursec{Utilisation responsable de l'IA}

Dans la section précédente, nous avons vu que l'IA offre de nombreux avantages (rapidité, aide à l'apprentissage, automatisation des tâches) mais aussi des inconvénients (erreurs possibles, dépendance, risques pour la vie privée). Une question se pose alors naturellement : \emph{comment profiter des avantages de l'IA tout en évitant ses dangers ?} La réponse tient en une expression : il faut adopter une \cle{utilisation responsable} de l'IA. C'est l'objet de cette section, qui te donnera des règles simples et concrètes à appliquer chaque jour, au collège comme à la maison.

\courssub{Qu'est-ce qu'une utilisation responsable de l'IA ?}

\begin{definition}[Utilisation responsable de l'IA]
L'\cle{utilisation responsable} de l'IA consiste à s'en servir de manière \textbf{véridique} (en vérifiant les informations), \textbf{respectueuse d'autrui} (sans nuire à personne), \textbf{protectrice de ses données} (sans révéler sa vie privée) et \textbf{honnête} (sans tricher, en citant l'aide reçue).
\end{definition}

L'intuition est simple : une IA est un outil puissant, comme une voiture. Une voiture est très utile, mais il faut respecter le code de la route pour ne pas provoquer d'accident. De même, l'IA possède son « code de la route » : un ensemble de règles que tout utilisateur doit connaître et appliquer. Nous allons étudier cinq grands principes, que l'on peut appeler les \cle{cinq réflexes} de l'utilisateur responsable.

\begin{popfigure}
\begin{center}
\begin{tikzpicture}[scale=0.95]
  % Cercle central
  \node[circle, draw=popblue, fill=popblueL, thick, minimum size=2.6cm, align=center, font=\small\bfseries] (centre) at (0,0) {Utilisateur\\responsable};
  % Les 5 réflexes en cercle autour
  \node[circle, draw=poporange, fill=poporangeL, thick, minimum size=2.1cm, align=center, font=\scriptsize\bfseries] (r1) at (90:3.6) {1. Vérifier\\les informations};
  \node[circle, draw=popgreen, fill=popgreenL, thick, minimum size=2.1cm, align=center, font=\scriptsize\bfseries] (r2) at (162:3.6) {2. Protéger\\ses données};
  \node[circle, draw=poppurple, fill=poppurpleL, thick, minimum size=2.1cm, align=center, font=\scriptsize\bfseries] (r3) at (234:3.6) {3. Respecter\\autrui};
  \node[circle, draw=poppink, fill=poppinkL, thick, minimum size=2.1cm, align=center, font=\scriptsize\bfseries] (r4) at (306:3.6) {4. Être\\honnête};
  \node[circle, draw=popgold, fill=popgoldL, thick, minimum size=2.1cm, align=center, font=\scriptsize\bfseries] (r5) at (18:3.6) {5. Citer\\l'usage de l'IA};
  % Flèches vers le centre
  \draw[-{Stealth}, thick, popmuted] (r1) -- (centre);
  \draw[-{Stealth}, thick, popmuted] (r2) -- (centre);
  \draw[-{Stealth}, thick, popmuted] (r3) -- (centre);
  \draw[-{Stealth}, thick, popmuted] (r4) -- (centre);
  \draw[-{Stealth}, thick, popmuted] (r5) -- (centre);
\end{tikzpicture}
\end{center}
\legende{Les cinq réflexes de l'utilisateur responsable de l'IA : vérifier, protéger, respecter, être honnête, citer.}
\end{popfigure}

\courssub{Principe 1 : vérifier les informations fournies par une IA}

\textbf{Intuition.} Une IA génère des réponses en prédisant des mots probables : elle ne « sait » pas ce qui est vrai. Elle peut donc affirmer avec beaucoup d'assurance des choses fausses (on parle parfois d'\emph{hallucination} de l'IA). Par exemple, une IA peut inventer une date d'indépendance, un auteur de roman ou une formule de physique. Croire une IA sur parole, c'est comme croire un inconnu dans la rue sans vérifier.

\begin{propriete}[Règle de vérification]
Toute information importante fournie par une IA doit être \textbf{vérifiée auprès de sources fiables} : les livres scolaires, les sites officiels (ministères, organismes reconnus), les encyclopédies sérieuses, les enseignants. Une information non vérifiée ne doit jamais être recopiée dans un devoir.
\end{propriete}

\begin{methode}[La règle des 3 V avant de croire une réponse d'IA]
Avant d'accepter une réponse d'une IA, applique la règle des \cle{3 V} :
\begin{enumerate}
  \item \textbf{Vérifier la source} : l'IA indique-t-elle d'où vient l'information ? Cherche cette information dans un livre ou sur un site officiel.
  \item \textbf{Vérifier avec un adulte ou un document fiable} : demande confirmation à un enseignant, un parent, ou compare avec ton manuel scolaire.
  \item \textbf{Vérifier la date de l'information} : l'information est-elle récente ? Une IA peut donner des données anciennes ou dépassées (par exemple un ancien prix, un ancien classement).
\end{enumerate}
Si l'un des trois contrôles échoue, n'utilise pas l'information telle quelle.
\end{methode}

\courssub{Principe 2 : protéger ses données personnelles}

\textbf{Intuition.} Quand tu écris dans un outil d'IA en ligne, tes messages peuvent être enregistrés sur des serveurs lointains et servir à entraîner le modèle. Tout ce que tu saisis peut donc sortir de ton contrôle. Il faut traiter un chatbot comme un lieu public : on n'y crie pas ses secrets.

\begin{propriete}[Règle de protection des données]
On ne saisit \textbf{jamais} dans un outil d'IA : ses mots de passe, son numéro de pièce d'identité (CNI), ses numéros de compte mobile money ou bancaire, son adresse précise, ni de photos privées de soi-même ou d'autres personnes.
\end{propriete}

\begin{attention}[Erreur fréquente : confier sa vie privée à un chatbot]
Beaucoup d'élèves saisissent sans réfléchir leur nom complet, leur adresse, leur numéro de téléphone ou le nom de leur école dans un chatbot public. C'est dangereux : ces informations peuvent être conservées, réutilisées ou détournées (usurpation d'identité, démarchage, harcèlement). \textbf{Réflexe} : avant d'envoyer un message, demande-toi : « est-ce que j'afficherais cela sur un panneau au marché ? » Si la réponse est non, ne l'écris pas.
\end{attention}

\courssub{Principe 3 : respecter les droits d'autrui}

\textbf{Intuition.} Les IA modernes peuvent créer des images et des textes très réalistes. Cette puissance peut être détournée pour fabriquer de fausses images d'une personne (par exemple un montage humiliant d'un camarade) ou pour répandre des rumeurs. Ces actes blessent des personnes réelles et peuvent être punis par la loi.

\begin{propriete}[Règle du respect d'autrui]
Il est interdit d'utiliser une IA pour \textbf{créer ou diffuser} de fausses images, de faux enregistrements ou des rumeurs concernant une personne. On ne publie jamais un contenu qui porte atteinte à la dignité, à l'honneur ou à la réputation d'autrui.
\end{propriete}

Retiens que « ce n'est qu'une blague » n'est pas une excuse : une fausse image partagée dans un groupe WhatsApp de classe peut humilier durablement un camarade et t'exposer à des sanctions disciplinaires, voire judiciaires.

\courssub{Principe 4 : l'honnêteté scolaire}

\textbf{Intuition.} L'IA peut rédiger un exposé entier en quelques secondes. Mais si tu recopies ce texte sans le comprendre, tu n'as rien appris : au prochain contrôle, tu seras seul devant ta copie, sans IA. L'IA doit être un \emph{professeur particulier} qui t'explique, pas un \emph{nègre} qui travaille à ta place.

\begin{exemplebox}[Les deux élèves et l'exposé]
\textbf{Cas 1 — mauvaise utilisation.} Amina demande à une IA : « Rédige mon exposé sur le fleuve Sanaga », recopie le texte mot pour mot et le rend. Elle \textbf{triche} : le travail n'est pas le sien, elle ne maîtrise pas le contenu, et elle sera incapable de répondre aux questions du professeur. Elle n'a rien appris.

\textbf{Cas 2 — bonne utilisation.} Junior demande à l'IA : « Explique-moi le rôle économique du fleuve Sanaga » et « Quels sont les barrages construits sur ce fleuve ? ». Il lit les réponses, les vérifie dans son manuel de géographie, prend des notes, puis \textbf{rédige lui-même} son exposé avec ses propres phrases. Il utilise l'IA correctement : comme un outil pour comprendre et s'entraîner.
\end{exemplebox}

\begin{aretenir}[La bonne question à se poser]
Avant d'utiliser une réponse d'IA dans un travail scolaire, demande-toi : « serais-je capable de refaire ce travail seul, sans l'IA, et de l'expliquer au professeur ? » Si oui, ton usage est honnête. Si non, tu es en train de copier.
\end{aretenir}

\courssub{Principe 5 : citer l'usage de l'IA}

\textbf{Intuition.} Quand tu utilises un livre pour rédiger un devoir, tu cites le livre (titre, auteur). C'est une règle d'honnêteté intellectuelle. Il en va de même pour l'IA : si elle t'a aidé à produire un travail, tu dois le déclarer.

\begin{propriete}[Règle de citation de l'IA]
Quand une IA a contribué à un travail (recherche d'idées, explication, correction), il faut le mentionner, par exemple en fin de devoir : « \emph{Ce travail a été préparé avec l'aide d'un assistant IA pour [préciser l'usage] ; le texte final a été rédigé et vérifié par moi-même.} »
\end{propriete}

Citer l'IA n'est pas un aveu de faiblesse : c'est une preuve d'honnêteté et de maturité, de plus en plus exigée dans les établissements scolaires.

\courssub{Éthique de l'IA et cadre légal}

\begin{definition}[Éthique de l'IA]
L'\cle{éthique de l'IA} est l'ensemble des règles morales qui guident la \textbf{conception} et l'\textbf{utilisation} des machines intelligentes, afin qu'elles respectent la \textbf{dignité humaine} et la \textbf{justice} : ne pas nuire, ne pas discriminer, ne pas tromper, laisser à l'être humain le dernier mot dans les décisions importantes.
\end{definition}

L'éthique concerne donc deux groupes de personnes : ceux qui \emph{fabriquent} les IA (ils doivent éviter les programmes racistes, mensongers ou dangereux) et ceux qui les \emph{utilisent} — c'est-à-dire toi, avec les cinq principes étudiés ci-dessus.

\begin{savaistu}[Le Cameroun et la sécurité numérique]
Le Cameroun dispose d'une agence spécialisée : l'\cle{ANTIC} (Agence Nationale des Technologies de l'Information et de la Communication), chargée de la sécurité des technologies de l'information et de la lutte contre la cybercriminalité. De plus, la \textbf{loi n° 2010/012 du 21 décembre 2010} relative à la cybersécurité et à la cybercriminalité punit au Cameroun des actes comme la fraude informatique, l'usurpation d'identité ou la diffusion de contenus illicites. Utiliser une IA pour escroquer ou diffamer quelqu'un peut donc avoir de vraies conséquences judiciaires, même pour des faits commis « pour rire » depuis un téléphone.
\end{savaistu}

\courssub{Un réflexe final : l'arbre de décision avant de partager}

Avant de publier ou de transférer un contenu produit avec une IA (texte, image, information), fais-le passer par trois questions. Si une seule réponse est « non », ne partage pas.

\begin{popfigure}
\begin{center}
\begin{tikzpicture}[
  question/.style={diamond, draw=popblue, fill=popblueL, thick, align=center, inner sep=1.5pt, font=\scriptsize\bfseries, aspect=2.2},
  stop/.style={rectangle, rounded corners, draw=poppink, fill=poppinkL, thick, align=center, font=\scriptsize\bfseries, minimum height=0.8cm},
  ok/.style={rectangle, rounded corners, draw=popgreen, fill=popgreenL, thick, align=center, font=\scriptsize\bfseries, minimum height=0.8cm},
  lbl/.style={font=\scriptsize\itshape, midway, above},
  node distance=0.9cm and 1.6cm]
  \node[question, align=center] (q1){Est-ce\\vrai ?};
  \node[stop, right=of q1, align=center] (n1){Ne pas\\partager};
  \node[question, below=1.3cm of q1, align=center] (q2){Est-ce\\respectueux ?};
  \node[stop, right=of q2, align=center] (n2){Ne pas\\partager};
  \node[question, below=1.3cm of q2, align=center] (q3){Ai-je le droit\\de le diffuser ?};
  \node[stop, right=of q3, align=center] (n3){Ne pas\\partager};
  \node[ok, below=1.3cm of q3, align=center] (fin){Je peux partager\\(en citant l'IA)};
  \draw[-{Stealth}, thick] (q1) -- node[lbl]{non} (n1);
  \draw[-{Stealth}, thick] (q1) -- node[lbl, left]{oui} (q2);
  \draw[-{Stealth}, thick] (q2) -- node[lbl]{non} (n2);
  \draw[-{Stealth}, thick] (q2) -- node[lbl, left]{oui} (q3);
  \draw[-{Stealth}, thick] (q3) -- node[lbl]{non} (n3);
  \draw[-{Stealth}, thick] (q3) -- node[lbl, left]{oui} (fin);
\end{tikzpicture}
\end{center}
\legende{Arbre de décision : « Puis-je partager ce contenu généré par IA ? » Trois questions, trois conditions obligatoires.}
\end{popfigure}

\begin{exoresolu}[Le devoir d'histoire de Sandrine]
\textbf{Énoncé.} Sandrine doit rédiger un paragraphe sur l'indépendance du Cameroun. Elle demande à une IA, qui lui répond que le Cameroun français est devenu indépendant le 1\textsuperscript{er} janvier 1960. Sandrine veut recopier cette phrase dans son devoir et l'envoyer aussi dans le groupe WhatsApp de la classe pour « aider » ses camarades. Que doit-elle faire pour être une utilisatrice responsable ?
\tcblower
\textbf{Solution détaillée.}
\begin{enumerate}
  \item \textbf{Vérifier (règle des 3 V).} Sandrine ouvre son manuel d'histoire et confirme la date : le 1\textsuperscript{er} janvier 1960 est bien la date de l'indépendance du Cameroun français. Elle peut aussi demander confirmation à son professeur. L'information est vraie et à jour (c'est un fait historique).
  \item \textbf{Rédiger elle-même.} Au lieu de recopier la phrase de l'IA, elle reformule avec ses propres mots, en s'appuyant sur ses notes de cours. Ainsi, elle comprend et retient la leçon.
  \item \textbf{Citer l'IA.} À la fin de son devoir, elle écrit : « Date vérifiée avec l'aide d'un assistant IA et du manuel scolaire ; texte rédigé par moi-même. »
  \item \textbf{Avant de partager sur WhatsApp}, elle applique l'arbre de décision : l'information est vraie (oui), respectueuse (oui), et elle a le droit de la diffuser (oui, c'est un fait public). Elle peut donc partager — mais elle préfère envoyer sa propre reformulation plutôt que le texte brut de l'IA, pour rester honnête.
\end{enumerate}
Sandrine a appliqué les cinq réflexes : vérifier, protéger (elle n'a saisi aucune donnée personnelle), respecter, être honnête et citer.
\end{exoresolu}

\begin{exercice}[À toi de juger]
Pour chaque situation, dis si l'utilisation de l'IA est responsable ou non, et justifie ta réponse en citant le principe concerné :
\begin{enumerate}
  \item Kevin demande à une IA de lui expliquer les fractions, refait les exercices seul, puis vérifie ses réponses avec l'IA.
  \item Brice crée avec une IA une fausse photo de son voisin de classe et la publie dans le groupe de la classe « pour rigoler ».
  \item Larissa tape dans un chatbot : « Je m'appelle Larissa N., j'habite quartier Nkomkana à Yaoundé, mon numéro est... » pour « se présenter ».
  \item Moussa utilise un plan proposé par une IA pour son exposé, rédige lui-même le contenu et indique en bas de page que l'IA l'a aidé à structurer le plan.
\end{enumerate}
\end{exercice}

\begin{corrige}[Exercice « À toi de juger »]
\begin{enumerate}
  \item \textbf{Responsable} : Kevin utilise l'IA pour comprendre et s'entraîner (principe 4, honnêteté scolaire), pas pour copier.
  \item \textbf{Non responsable} : Brice viole le principe 3 (respect des droits d'autrui) ; créer et diffuser une fausse image de quelqu'un nuit à sa dignité et peut être puni par la loi de 2010 sur la cybersécurité et la cybercriminalité.
  \item \textbf{Non responsable} : Larissa viole le principe 2 (protection des données personnelles) ; on ne saisit jamais son identité complète, son adresse ni son numéro de téléphone dans un chatbot public.
  \item \textbf{Responsable} : Moussa rédige lui-même (principe 4) et cite l'usage de l'IA (principe 5) ; il a aussi intérêt à vérifier les informations de son plan (principe 1).
\end{enumerate}
\end{corrige}

\begin{aretenir}[L'essentiel de la section]
\begin{itemize}
  \item Utiliser l'IA de manière responsable, c'est être \textbf{véridique, respectueux, protecteur et honnête}.
  \item Les cinq réflexes : \textbf{vérifier} les informations (règle des 3 V), \textbf{protéger} ses données personnelles, \textbf{respecter} autrui, être \textbf{honnête} à l'école, \textbf{citer} l'usage de l'IA.
  \item L'\textbf{éthique de l'IA} vise des machines qui respectent la dignité humaine et la justice.
  \item Au Cameroun, la loi de 2010 sur la cybersécurité et la cybercriminalité et l'ANTIC encadrent et sanctionnent les usages abusifs du numérique.
\end{itemize}
\end{aretenir}

\coursec{Apprendre à échanger avec une IA}

Après avoir réfléchi aux \emph{enjeux éthiques} et aux \emph{bonnes pratiques} pour une utilisation responsable de l'IA (section précédente), nous devons maintenant acquérir une compétence technique concrète : \textbf{savoir lui parler}. Une IA conversationnelle n'est pas un moteur de recherche ; elle ne « trouve » pas une réponse toute faite, elle \emph{génère} une réponse en fonction de la manière dont vous lui posez la question. La qualité de votre échange détermine directement la qualité du résultat. C'est l'objet de cette dernière section : transformer un dialogue hasardeux en une collaboration efficace.

\courssub{Qu'est-ce qu'un prompt ?}

\begin{definition}[Prompt]
Un \textbf{prompt} (ou \emph{invite} en français) est le message — question, consigne, instruction ou fragment de texte — qu'un utilisateur écrit pour déclencher une réponse de la part d'une IA conversationnelle (comme ChatGPT, Gemini, Copilot, ou un modèle local). C'est l'unique interface de commande : tout ce que l'IA « sait » de votre intention passe par ce texte.
\end{definition}

Contrairement à une recherche Google où quelques mots-clés suffisent (\texttt{cacao Cameroun rendement}), une IA a besoin d'une \emph{phrase complète}, structurée et contextualisée. Imaginez que vous donniez une consigne à un stagiaire très intelligent mais qui ne connaît rien de votre classe, de votre niveau, ni de votre objectif : s'il reçoit « Parle-moi de l'agriculture », il risque d'écrire une thèse de doctorat ou un article pour enfants. Le prompt \textbf{est} le contrat de mission que vous signez avec l'IA.

\courssub{Les qualités d'un bon prompt}

Un prompt efficace repose sur quatre piliers. S'il en manque un, la réponse risque d'être hors sujet, trop vague, trop technique ou mal formatée.

\begin{propriete}[Qualités d'un bon prompt]
\begin{enumerate}
    \item \textbf{Clarté :} Le langage est simple, direct, sans ambiguïté. Pas de sous-entendus.
    \item \textbf{Précision :} Les contraintes sont explicites (longueur, ton, éléments à inclure ou exclure).
    \item \textbf{Contexte :} L'IA sait \emph{qui} elle est, \emph{à qui} elle s'adresse, et \emph{dans quelle situation} se place la demande.
    \item \textbf{Objectif explicite :} Le résultat attendu est décrit (une explication, un code, un tableau, une idée de sujet, une correction d'exercice…).
\end{enumerate}
\end{propriete}

\courssub{Structure d'un prompt efficace : la méthode RTCF}

Pour ne rien oublier, les experts utilisent une structure en quatre blocs. C'est la « méthode RTCF » (Rôle, Tâche, Contexte, Format). Appliquez-la systématiquement pour vos devoirs, vos révisions ou vos projets.

\begin{methode}[La méthode RTCF — Construire un prompt complet]
\begin{enumerate}
    \item \textbf{R — Rôle (Persona) :} Définissez qui l'IA doit incarner. \textit{Ex : « Agis comme un professeur de SVT au Cameroun. »}
    \item \textbf{T — Tâche (Instruction) :} Décrivez l'action précise à réaliser. \textit{Ex : « Explique la photosynthèse. »}
    \item \textbf{C — Contexte (Cible \& Situation) :} Précisez le public, le niveau, le pays, le programme. \textit{Ex : « Pour un élève de 3ème suivant le programme MINESEC, avec des exemples locaux (cacao, bananier). »}
    \item \textbf{F — Format (Sortie) :} Imposer la forme de la réponse. \textit{Ex : « En 5 lignes maximum, avec un schéma textuel (flèches) et une analogie simple. »}
\end{enumerate}
\textbf{Prompt final assemblé :} « Agis comme un professeur de SVT au Cameroun (R). Explique la photosynthèse (T) à un élève de 3ème programme MINESEC avec exemples locaux cacao/bananier (C) en 5 lignes max, avec schéma textuel et analogie (F). »
\end{methode}

\begin{popfigure}
\begin{tikzpicture}[
    node distance=0.8cm and 1cm,
    block/.style={draw=popblue, thick, fill=popblueL, rounded corners, minimum height=1.8cm, minimum width=5.5cm, text width=5cm, align=center, font=\small},
    label/.style={font=\bfseries\small, text=popdark, anchor=south},
    arrow/.style={-Stealth, thick, poporange}
]
% Les 4 blocs empilés
\node[block, align=center] (R){\textbf{R — RÔLE}\\\texttt{Agis comme un prof de SVT au Cameroun}};
\node[block, below=of R, align=center] (T){\textbf{T — TÂCHE}\\\texttt{Explique la photosynthèse}};
\node[block, below=of T, align=center] (C){\textbf{C — CONTEXTE}\\\texttt{Élève 3ème MINESEC, ex. cacao/bananier}};
\node[block, below=of C, align=center] (F){\textbf{F — FORMAT}\\\texttt{5 lignes, schéma textuel, analogie}};

% Flèches de liaison
\draw[arrow] (R.south) -- (T.north);
\draw[arrow] (T.south) -- (C.north);
\draw[arrow] (C.south) -- (F.north);

% Étiquette "Prompt final" à droite du groupe
\node[right=1.2cm of R, align=left, font=\small, poppurple, anchor=north west] (lbl) {\textbf{Prompt complet}\\\textit{Copiez-collez ce bloc dans l'IA}};
\draw[poppurple, thick] (R.north east) -- (lbl.north west);
\draw[poppurple, thick] (F.south east) -- (lbl.south west);
\end{tikzpicture}
\legende{Les quatre composantes empilées de la méthode RTCF. Chaque bloc affine la réponse : le Rôle donne le ton, la Tâche l'action, le Contexte la pertinence, le Format l'utilisabilité.}
\end{popfigure}

\courssub{Technique de l'amélioration progressive (Itération)}

Rarement le premier prompt est parfait. L'échange avec une IA est un \textbf{dialogue itératif} : vous observez la réponse, vous identifiez ce qui manque ou ce qui dérange, et vous \emph{reformulez} en ajoutant une contrainte.

\begin{exemplebox}[De « faible » à « efficace » : l'itération en action]
\textbf{Objectif :} Obtenir une aide pour un exposé d'histoire-géo sur l'urbanisation à Douala.

\begin{enumerate}
    \item \textbf{Prompt 1 (Faible) :} \texttt{Parle-moi de l'urbanisation à Douala.}
    \item \textit{Réponse 1 :} Texte généraliste, niveau universitaire, pas de chiffres récents, pas de structure pour un exposé.
    \item \textbf{Prompt 2 (Précision + Format) :} \texttt{Fais un plan d'exposé de 3 parties sur l'urbanisation à Douala pour un élève de 3ème.}
    \item \textit{Réponse 2 :} Meilleur plan, mais toujours pas de données locales (quartiers, population), ton trop scolaire.
    \item \textbf{Prompt 3 (RTCF complet) :} \texttt{Agis comme un prof d'Histoire-Géo au Cameroun (R). Propose un plan détaillé d'exposé oral de 5 min sur l'urbanisation à Douala (T) pour la classe de 3ème, programme MINESEC (C). Inclut : intro accroche, 3 parties avec arguments + exemples quartiers (New Bell, Bonamoussadi, PK10), conclusion ouverte. Format : puces, mots-clés en gras, temps estimé par partie (F).}
    \item \textit{Réponse 3 :} Directement utilisable pour préparer les fiches bristol.
\end{enumerate}
\end{exemplebox}

\courssub{Dialoguer par étapes : la stratégie du « suivi »}

Ne cherchez pas à tout obtenir en un seul message gigantesque. Découpez votre besoin en \emph{étapes logiques} et utilisez l'historique de la conversation.

\begin{propriete}[Dialogue par étapes (Chaînage de prompts)]
\begin{enumerate}
    \item \textbf{Étape 1 — Cadrage large :} Posez une question ouverte pour explorer le sujet. \textit{« Quelles sont les causes principales de l'exode rural vers Douala ? »}
    \item \textbf{Étape 2 — Approfondissement ciblé :} Rebondissez sur un point de la réponse. \textit{« Développe la cause "recherche d'emploi" avec des chiffres récents de l'INS si possible. »}
    \item \textbf{Étape 3 — Production finale :} Demandez la mise en forme. \textit{« Résume les 3 causes principales en un tableau Markdown avec colonnes : Cause, Mécanisme, Exemple local, Source potentielle. »}
\end{enumerate}
\end{propriete}

Cette approche évite la « surcharge cognitive » de l'IA (perte du fil sur les longs prompts) et vous permet de valider la justesse de chaque étape avant de passer à la suivante.

\courssub{Attitude critique : vérifier, croiser, conclure}

L'IA \textbf{hallucine} : elle peut inventer des dates, des noms d'auteurs, des articles de loi, des formules chimiques. Elle n'a pas de conscience de la vérité, seulement une probabilité statistique de mots qui « vont bien ensemble ».

\begin{attention}[Erreur fréquente : la confiance aveugle]
Ne copiez \textbf{jamais} une réponse d'IA dans votre devoir, votre exposé ou votre code sans l'avoir :
\begin{itemize}
    \item \textbf{Relue intégralement} par vous-même.
    \item \textbf{Comparée} avec votre cours, votre manuel ou une source fiable (site .gouv.cm, manuel homologué, encyclopédie reconnue).
    \item \textbf{Questionnée} sur ses sources : \texttt{« Quelles sont tes sources pour le chiffre de 3,8 millions d'habitants à Douala en 2023 ? Donne-moi le lien exact du rapport de l'INS. »}
\end{itemize}
Si l'IA ne peut pas citer une source précise et vérifiable, considérez l'info comme \textbf{non validée}.
\end{attention}

\begin{methode}[Rédiger une conclusion justifiée après échange avec l'IA]
Après avoir dialogué avec l'IA pour comprendre un concept ou résoudre un problème, forcez-vous à produire \textbf{votre propre synthèse écrite} (sur votre cahier ou en bas du TP). Structurez-la ainsi :
\begin{enumerate}
    \item \textbf{Ce que j'ai demandé :} Résumez le prompt final utilisé (RTCF).
    \item \textbf{Ce que l'IA a apporté :} L'idée clé, la structure, l'exemple, le bout de code, la correction d'erreur.
    \item \textbf{Ce que j'ai vérifié / corrigé :} « L'IA a dit X, mais mon cours dit Y / le site de l'INS dit Z → je retiens Y. »
    \item \textbf{Ce que je retiens EN MES MOTS :} Une phrase ou deux qui prouvent que vous avez compris, sans copier-coller.
\end{enumerate}
Cette trace écrite est votre \emph{preuve d'apprentissage} ; l'IA n'était qu'un tuteur temporaire.
\end{methode}

\courssub{Tableau comparatif : Prompt faible vs Prompt amélioré}

\begin{popfigure}
\begin{center}
\begin{tabularx}{\linewidth}{|>{\bfseries}l|X|X|X|}
\hline
\rowcolor{popblueL}
\textbf{Critère} & \textbf{Prompt faible} & \textbf{Prompt amélioré (RTCF)} & \textbf{Qualité de la réponse obtenue} \\
\hline
\textbf{Exemple 1 : SVT (Photosynthèse)} &
\texttt{Explique la photosynthèse.} &
\texttt{Agis comme un prof SVT 3ème Cameroun (R). Explique la photosynthèse (T) niveau 3ème MINESEC avec ex. cacao (C). 5 lignes, schéma fléché, analogie cuisine (F).} &
\begin{minipage}{\linewidth}\vspace{2pt}
\textbf{Faible :} Cours universitaire, équations chimiques complexes, pas d'exemple local.\newline
\textbf{Amélioré :} Niveau adapté, analogie « cuisine » (soleil = four, feuilles = cuisines, CO2 = ingrédients), schéma textuel clair, vocabulaire du programme.
\end{minipage} \\
\hline
\textbf{Exemple 2 : Maths (Fonctions)} &
\texttt{Comment tracer une fonction ?} &
\texttt{Agis comme un tuteur maths 3ème (R). Montre la méthode pas à pas pour tracer $f(x)=2x-3$ (T) pour un élève qui prépare le BEPC (C). Étapes numérotées, tableau de valeurs, repère orthonormé décrit en texte, pièges à éviter (F).} &
\begin{minipage}{\linewidth}\vspace{2pt}
\textbf{Faible :} Définition générale, pas de méthode pas à pas, pas de repère.\newline
\textbf{Amélioré :} Méthode BEPC standard (tableau 3 lignes, placer points, tracer droite), pièges (échelle, ordonnée à l'origine), prêt pour fiche de révision.
\end{minipage} \\
\hline
\textbf{Exemple 3 : Vie courante (Budget)} &
\texttt{Aide-moi à faire un budget.} &
\texttt{Agis comme un conseiller éco-gestion pour lycéen (R). Crée un modèle de budget mensuel simple (T) pour un élève de 3ème à Yaoundé qui a 15 000 FCFA/mois (transport, goûter, crédit téléphone, épargne) (C). Tableau Markdown avec colonnes Poste, Prévu, Réel, Écart + 3 conseils épargne (F).} &
\begin{minipage}{\linewidth}\vspace{2pt}
\textbf{Faible :} Modèle entreprise complexe, devises \$, catégories inadaptées.\newline
\textbf{Amélioré :} Monnaie FCFA, montants réalistes (ex: transport 3000, goûter 4000), tableau copiable dans Excel/Sheets, conseils concrets (ex: « marche 2 arrêts »).
\end{minipage} \\
\hline
\end{tabularx}
\end{center}
\legende{Comparaison concrète : le prompt RTCF transforme une réponse générique en un outil pédagogique directement utilisable par un élève de 3ème au Cameroun.}
\end{popfigure}

\courssub{Exercice résolu : Construire son prompt pour un besoin réel}

\begin{exoresolu}[Préparer une fiche de révision d'anglais : les temps du récit]
\textbf{Situation :} Vous avez un contrôle d'anglais demain sur les temps du récit (\textit{Past Simple, Past Continuous, Past Perfect}). Vous voulez une fiche de révision claire, avec des exemples camerounais, et un mini-exercice auto-corrigé.

\textbf{Question :} Rédigez le prompt RTCF complet que vous enverrez à l'IA.

\textbf{Solution détaillée :}
On applique la méthode RTCF point par point.
\begin{enumerate}
    \item \textbf{R (Rôle) :} « Agis comme un professeur d'anglais au Cameroun, correcteur au BEPC. »
    \item \textbf{T (Tâche) :} « Crée une fiche de révision complète sur les temps du récit (Past Simple, Past Continuous, Past Perfect). »
    \item \textbf{C (Contexte) :} « Pour un élève de 3ème francophone qui prépare le BEPC. Exemples ancrés au Cameroun (ex: match Lions Indomptables, marché Mokolo, trajet école). Niveau A2/B1. »
    \item \textbf{F (Format) :} « Structure : 1) Règle simple par temps (forme + emploi + mots-clés). 2) 3 phrases d'exemple par temps (contexte Cameroun). 3) Tableau récapitulatif comparatif. 4) Mini-exercice : 5 phrases à trous (verbes entre parenthèses) + corrigé détaillé en bas (caché ou à la fin). Ton encourageant. »
\end{enumerate}
\textbf{Prompt final à copier-coller :}
\begin{lstlisting}
Agis comme un prof d'anglais au Cameroun, correcteur BEPC (R).
Crée une fiche de révision sur les temps du récit (Past Simple, Past Continuous, Past Perfect) (T)
pour un élève de 3ème francophone BEPC, exemples Cameroun (match Lions, marché Mokolo) niveau A2/B1 (C).
Format : 1) Règle simple par temps (forme, emploi, mots-clés). 2) 3 ex. par temps (contexte CM). 3) Tableau récap comparatif. 4) Mini-exo 5 phrases à trous + corrigé détaillé à la fin. Ton encourageant (F).
\end{lstlisting}
\textbf{Analyse de la réponse attendue :} L'IA produira un document structuré, au bon niveau, avec des phrases comme « \textit{The Lions \textbf{were playing} when it \textbf{started} to rain} » (Past Continuous + Past Simple) au lieu de « \textit{John was reading when the phone rang} ». Le corrigé vous permettra de vous auto-évaluer.
\end{exoresolu}

\courssub{Exercices d'entraînement}

\begin{exercice}[À vous de jouer : Prompts pour le BEPC]
Rédigez le prompt RTCF complet (les 4 lignes R, T, C, F) pour chacun des besoins suivants. Testez-les ensuite sur une IA accessible (via téléphone ou salle informatique) et comparez les résultats.

\begin{enumerate}
    \item \textbf{Histoire-Géo :} Vous devez présenter oralement (3 min) « Les défis de l'industrialisation au Cameroun ». Vous voulez un plan détaillé avec arguments chiffrés (PIB, emploi) et exemples d'usines locales (Socapalm, Alucam, brasseries).
    \item \textbf{Mathématiques :} Vous bloquez sur un exercice de probabilités : « Un sac contient 4 boules rouges, 3 bleues, 2 vertes. On tire 2 boules sans remise. Calculer la probabilité d'avoir 2 boules de même couleur. » Vous voulez la \emph{méthode pas à pas} (arbre, calculs) et l'explication du « sans remise ».
    \item \textbf{Français / Expression écrite :} Sujet de dissertation : « Le téléphone portable : un outil indispensable ou un danger pour la jeunesse ? » Vous voulez une introduction type (amorce, sujet, thèse, plan), 2 arguments pour le « pour » avec exemples locaux, 2 arguments pour le « contre », et une conclusion ouverte.
    \item \textbf{Informatique (Algorithme) :} Écrire un algorithme en pseudo-code (style MINESEC : \texttt{Algorithme}, \texttt{Variables}, \texttt{Debut}, \texttt{Fin}) qui demande l'âge de 30 élèves, compte combien sont majeurs ($\ge$ 18 ans) et affiche le résultat. Vous voulez le code commenté et une table de trace pour 5 valeurs test.
\end{enumerate}
\end{exercice}

\begin{corrige}[Exercice 1 — Histoire-Géo (Industrialisation)]
\textbf{R :} Agis comme un prof d'Histoire-Géo 3ème Cameroun, spécialiste économie locale. \\
\textbf{T :} Propose un plan d'exposé oral de 3 min sur les défis de l'industrialisation au Cameroun. \\
\textbf{C :} Élève 3ème programme MINESEC. Inclut : intro (définition, enjeu), 3 défis majeurs (financement, énergie, main-d'œuvre qualifiée) avec chiffres récents (PIB part industrie, taux électrification) et exemples usines (Socapalm, Alucam, Brasseries, Cimencam). Conclusion (pôle industriel Kribi, ZIC). \\
\textbf{F :} Plan numéroté (I, A, 1, a), temps de parole estimé par partie, mots-clés en gras, sources suggérées (INS, MINMIDT, Banque Mondiale) pour vérification.
\end{corrige}

\begin{corrige}[Exercice 2 — Maths (Probabilités)]
\textbf{R :} Agis comme un tuteur de maths 3ème, expert BEPC. \\
\textbf{T :} Résous l'exercice : sac 4R, 3B, 2V. Tirage 2 boules sans remise. Proba(2 mêmes couleurs). \\
\textbf{C :} Élève 3ème qui prépare le BEPC. Montre la méthode arbre de probabilités (branches, probabilités conditionnelles). Explique pourquoi « sans remise » change les dénominateurs. \\
\textbf{F :} Étapes numérotées : 1) Arbre textuel (ASCII ou description). 2) Calculs détaillés pour RR, BB, VV. 3) Somme finale. 4) Piège fréquent (oublier majoration conditionnelle). 5) Résultat final sous forme fraction simplifiée et décimale.
\end{corrige}

\begin{corrige}[Exercice 3 — Français (Dissertation)]
\textbf{R :} Agis comme un prof de français 3ème, correcteur BEPC. \\
\textbf{T :} Rédige une introduction type et un plan détaillé dialectique pour : « Le téléphone portable : outil indispensable ou danger pour la jeunesse ? » \\
\textbf{C :} Élève 3ème Cameroun. Arguments ancrés : « Pour » (scolarité : Khan Academy, dictionnaires ; sécurité : appeler parents ; argent mobile : Mobile Money). « Contre » (addiction réseaux sociaux, cyberharcèlement, arnaques, perturbation sommeil/études). \\
\textbf{F :} 1) Introduction complète (amorce, définition, thèse nuancée, plan annoncé). 2) Plan en 2 parties (I. Outil d'émancipation / II. Risque d'aliénation) avec 2 arguments + exemple CM chacun. 3) Conclusion (synthèse, ouverture sur éducation numérique). Ton pédagogique.
\end{corrige}

\begin{corrige}[Exercice 4 — Informatique (Algorithme)]
\textbf{R :} Agis comme un prof d'informatique 3ème, programme MINESEC (pseudo-code français standard). \\
\textbf{T :} Écris l'algorithme : saisie âges 30 élèves, comptage majeurs ($\ge$18), affichage compteur. \\
\textbf{C :} Élève 3ème révision algorithmique. Variables types entiers. Boucle \texttt{Pour} ou \texttt{TantQue}. Structure conditionnelle \texttt{Si}. \\
\textbf{F :} 1) Code complet indenté (Algorithme, Variables, Debut, Fin). 2) Commentaires sur chaque ligne clé. 3) Table de trace pour 5 âges test (ex: 15, 18, 19, 16, 20) avec colonnes : Itération, AgeLu, CompteurMajeur. 4) Résultat final attendu.
\end{corrige}

\begin{savaistu}[L'ingénierie de prompts : un métier d'avenir]
Savoir formuler des prompts précis, structurés et itératifs est devenu une compétence professionnelle à part entière : on l'appelle \textbf{l'ingénierie de prompts} (\textit{Prompt Engineering}). Aux États-Unis comme au Cameroun, des entreprises recrutent des « \textit{Prompt Engineers} » pour optimiser l'usage de l'IA dans le marketing, le droit (rédaction de contrats), le code (génération/débogage), l'éducation (création de cours), la santé (résumés de dossiers). En 3ème, en maîtrisant la méthode RTCF, vous ne faites pas que « tricher » : vous apprenez à \textbf{piloter un outil puissant} — une compétence aussi fondamentale aujourd'hui que savoir utiliser un traitement de texte ou un tableur il y a 20 ans.
\end{savaistu}

\begin{aretenir}[Bilan : Bien dialoguer avec une IA]
\begin{itemize}
    \item Le \cle{prompt} est votre unique levier de contrôle : soignez-le comme une consigne officielle.
    \item Appliquez la méthode \textbf{RTCF} (Rôle, Tâche, Contexte, Format) pour toute demande importante.
    \item Procédez \textbf{par itérations} : affinez, corrigez, demandez des exemples, simplifiez.
    \item Gardez un \textbf{esprit critique} : l'IA hallucine, vérifiez les faits (chiffres, dates, code, sources) avec vos cours ou sources officielles.
    \item \textbf{Rédigez votre propre conclusion} après l'échange : c'est la trace de \emph{votre} apprentissage, pas celle de la machine.
\end{itemize}
\end{aretenir}

\newpage

\coursec{Activité d'intégration}

\courssub{Objectifs de l'activité}
Cette activité d'intégration vise à mobiliser l'ensemble des savoirs et savoir-faire de la leçon « Découverte d'un environnement d'IA » à travers une situation authentique ancrée dans le contexte camerounais. Elle permet à l'élève de :
\begin{itemize}
    \item Analyser un projet concret d'utilisation de l'IA dans un service public de proximité (le grand marché) en identifiant les \cle{avantages} attendus et les \cle{inconvénients} ou \cle{risques} potentiels (vie privée, erreurs, coût, dépendance technique).
    \item Mettre en œuvre une \cle{communication efficace avec une IA} : rédaction d'une \cle{consigne} (ou \cle{prompt}) claire et structurée pour obtenir des réponses pertinentes.
    \item Exercer son \cle{esprit critique} et son \cle{savoir-être numérique} en produisant un avis argumenté, nuancé et justifié (position « pour », « contre » ou « pour sous conditions »), puis en le défendant oralement devant la classe.
    \item Travailler la \cle{collaboration} en groupe (rôles répartis : secrétaire, porte-parole, chronomètreur, analyste) et la \cle{restitution collective} sous forme de tableau de synthèse.
\end{itemize}

\courssub{Situation}
\textbf{Contexte :} La mairie d'arrondissement de \textbf{Nkolbisson} (Yaoundé VII) lance un projet pilote « Marché Intelligent » pour moderniser la gestion du grand marché, fréquenté par plus de 5\,000 personnes par jour (vendeurs, acheteurs, transporteurs). Le maire, conseillé par une entreprise locale « \emph{CamAI Solutions} », propose d'installer trois applications d'intelligence artificielle :

\begin{enumerate}
    \item \textbf{Module « Accès Sécurisé » :} Un système de \cle{reconnaissance faciale} à l'entrée principale. Chaque vendeur titulaire d'une carte de place (environ 1\,200 vendeurs recensés) est photographié lors de son inscription. Le système compare le visage présenté à la caméra avec la base de données pour ouvrir le tourniquet. Objectif : lutter contre la revente illicite des cartes de place et contre les intrus.
    \item \textbf{Module « Assistant Client » :} Un \cle{agent conversationnel} (ou \emph{chatbot}) accessible via un code affiché sur 50 panneaux dans les allées et sur la page Facebook de la mairie. Il répond en français et en anglais aux questions fréquentes : « Où trouver le rayon ignames ? », « Quels sont les tarifs de location des étals ce mois-ci ? », « Comment déposer une plainte pour vol ? ». Il fonctionne 24h/24 sans opérateur humain de permanence.
    \item \textbf{Module « Prédi'Affluence » :} Un programme d'IA entraîné sur 3 ans de données de comptage (capteurs aux entrées), de calendrier (jours de paie, fêtes, saison des pluies) et de météo. Il prédit le nombre de visiteurs pour les 7 jours à venir avec une marge d'erreur annoncée de $\pm 12\%$. La mairie veut l'utiliser pour adapter le nombre d'agents de nettoyage (actuellement 12 agents) et de policiers municipaux (actuellement 4 agents) présents sur site.
\end{enumerate}

\textbf{Données chiffrées du projet :}
\begin{itemize}
    \item Coût d'installation (matériel + logiciel + formation) : \textbf{42 millions FCFA}.
    \item Coût de maintenance annuel estimé : \textbf{3,5 millions FCFA} (serveur, mises à jour, électricité, connexion Internet).
    \item Durée du contrat : 3 ans renouvelable.
    \item Les photos des visages et les conversations avec l'agent conversationnel sont conservées sur un serveur situé à l'étranger.
\end{itemize}

La mairie organise une \textbf{concertation citoyenne} avant le vote du conseil municipal. Votre classe représente une délégation d'élèves du collège de Nkolbisson, invitée à rendre un avis consultatif.

\courssub{Consignes}
Travaillant par groupes de 4 élèves (rôles : \emph{Secrétaire} note les idées, \emph{Porte-parole} présente, \emph{Chronomètreur} gère le temps, \emph{Analyste} vérifie la cohérence avec le cours), réalisez les 4 étapes suivantes en 50 minutes :

\begin{enumerate}
    \item \textbf{Analyse critique (15 min) :} Identifiez \textbf{deux avantages majeurs} et \textbf{deux risques majeurs} pour ce projet. Pour chaque risque, précisez sa nature (exemples : \emph{atteinte à la vie privée}, \emph{erreur de reconnaissance}, \emph{coût financier élevé}, \emph{dépendance technique}, \emph{exclusion des personnes sans téléphone}). Remplissez le tableau de l'annexe 1 (à reproduire sur votre cahier).
    \item \textbf{Rédaction d'une consigne pour l'IA (10 min) :} Rédigez \textbf{une consigne (prompt) complète et structurée} que vous enverriez à une IA conversationnelle pour obtenir une liste de \textbf{5 précautions} à imposer à la mairie avant de signer le contrat. Votre consigne doit respecter la méthode vue en classe : donner le \cle{contexte}, préciser le \cle{rôle} que l'IA doit jouer, formuler la \cle{tâche} attendue et indiquer la \cle{forme} de la réponse souhaitée. Si la salle n'a pas Internet, l'enseignant joue le rôle de l'IA et répond oralement aux consignes lues à haute voix.
    \item \textbf{Avis argumenté (15 min) :} Rédigez collectivement un \textbf{avis officiel d'environ 10 lignes} (80 à 120 mots), signé du groupe, prenant position : \textbf{« POUR »}, \textbf{« CONTRE »} ou \textbf{« POUR SOUS CONDITIONS »}. Chaque argument (au minimum 3) doit s'appuyer sur un fait de la situation ou sur une notion du cours (avantage, inconvénient, principe d'utilisation responsable : transparence, consentement, contrôle humain, respect de la vie privée).
    \item \textbf{Restitution et vote (10 min) :} Le porte-parole de chaque groupe lit l'avis (1 minute maximum). L'enseignant construit au tableau le \textbf{tableau de synthèse collectif} (colonnes : Avantages identifiés / Risques identifiés / Conditions posées / Position finale). La classe vote à main levée pour la position la plus convaincante.
\end{enumerate}

\courssub{Ressources à mobiliser}
\begin{itemize}
    \item \textbf{Notions de cours :} définition de l'IA (ensemble de techniques permettant à une machine d'imiter certaines capacités humaines : percevoir, comprendre, décider), exemples d'usages quotidiens (recommandation de vidéos, traduction automatique, reconnaissance d'images, assistants vocaux, prédiction).
    \item \textbf{Avantages de l'IA :} automatisation des tâches répétitives, rapidité de traitement de grandes quantités de données, disponibilité 24h/24 et 7j/7, aide à la décision (prévisions), accessibilité (traduction, synthèse vocale).
    \item \textbf{Inconvénients et risques :} \cle{erreurs} et \cle{biais} (données d'apprentissage non représentatives), \cle{atteinte à la vie privée} (photos, enregistrements), \cle{coût} (matériel, électricité, maintenance), \cle{dépendance} (panne, fournisseur unique), \cle{impact social} (suppression d'emplois, exclusion des personnes non équipées), \cle{consommation d'énergie}.
    \item \textbf{Principes d'une utilisation responsable de l'IA :} \cle{transparence} (savoir qu'on échange avec une machine), \cle{consentement} (accepter librement que ses données soient utilisées), \cle{contrôle humain} (un humain vérifie les décisions importantes), \cle{respect de la vie privée} (protéger les données personnelles), \cle{vérification des informations} (une IA peut se tromper).
    \item \textbf{Technique de la consigne (prompt) :} une bonne consigne donnée à une IA contient :
    \begin{itemize}
        \item le \textbf{contexte} : qui êtes-vous ? quelle est la situation ?
        \item le \textbf{rôle} : quel expert l'IA doit-elle jouer ?
        \item la \textbf{tâche} : que doit-elle faire précisément ?
        \item la \textbf{forme} de la réponse : longueur, tableau, liste, ton ?
    \end{itemize}
\end{itemize}

\courssub{Démarche et corrigé commenté}
\begin{exoresolu}[Corrigé type et guide pédagogique pour l'enseignant]
\textbf{Étape 1 -- Analyse critique (Tableau Annexe 1)}

\begin{center}
\begin{tabularx}{\linewidth}{|l|X|X|}
\hline
\rowcolor{popblueL}
\textbf{Module} & \textbf{Avantages identifiés (exemples attendus)} & \textbf{Risques identifiés (exemples attendus)} \\
\hline
Reconnaissance faciale (Accès) & 1. Sécurisation : fin de la fraude aux cartes de place (revente, prêt).\newline 2. Fluidité : entrée rapide, sans contrôle manuel lent. & 1. \textbf{Vie privée :} le visage est une donnée personnelle sensible. Risque de piratage ou d'utilisation des photos à d'autres fins. Les vendeurs ont-ils donné leur accord librement ?\newline 2. \textbf{Erreurs :} le système peut se tromper (changement de coiffure, lunettes, mauvaise lumière) et bloquer un vendeur légitime ou laisser passer un intrus. \\
\hline
Agent conversationnel (Assistant) & 1. Service disponible 24h/24, même la nuit et le week-end.\newline 2. Réponses en français et en anglais : accessible à une population variée. & 1. \textbf{Fiabilité :} l'IA peut donner des réponses fausses (tarif inventé, procédure erronée). Qui est responsable en cas d'erreur ?\newline 2. \textbf{Exclusion :} il faut un téléphone, une connexion et savoir lire. Les vendeuses âgées ou non scolarisées sont exclues du service. \\
\hline
Prédiction d'affluence & 1. Meilleure organisation : agents de nettoyage et policiers ajustés aux besoins réels, donc économies et propreté.\newline 2. Anticipation : détection des jours de forte affluence (fêtes, jours de paie). & 1. \textbf{Erreur de prédiction ($\pm 12\%$) :} trop d'agents (argent public gaspillé) ou pas assez (insalubrité, insécurité).\newline 2. \textbf{Dépendance et coût caché :} maintenance spécialisée, abonnements, formation. Que se passe-t-il si l'entreprise fait faillite ? \\
\hline
\end{tabularx}
\end{center}

\textbf{Étape 2 -- Rédaction d'une consigne (prompt) structurée}

\begin{lstlisting}
[CONTEXTE] Je suis eleve en classe de 3eme au Cameroun. Ma mairie
veut installer 3 applications d'IA au marche central (5000
visiteurs par jour) : 1/ reconnaissance faciale a l'entree
(1200 vendeurs), 2/ agent conversationnel pour informer les
clients, 3/ prediction de l'affluence pour organiser les agents.
Cout : 42 millions FCFA + 3,5 millions par an.

[ROLE] Tu es un expert en informatique et en protection des
donnees personnelles. Tu conseilles les citoyens.

[TACHE] Donne-moi 5 precautions que la mairie doit imposer
avant de signer le contrat, pour proteger la vie privee des
vendeurs, eviter les erreurs de l'IA et maitriser les couts.

[FORME] Reponds par une liste numerotee de 5 points. Pour
chaque point : un titre court, puis deux phrases d'explication
simples, comprehensibles par un eleve de 3eme.
\end{lstlisting}

\textbf{Étape 3 -- Avis argumenté (exemples de positions types)}

\paragraph{Position « POUR SOUS CONDITIONS » (recommandée pour montrer l'esprit critique) :}
\begin{quote}
\textbf{Avis du Groupe 3 -- POUR SOUS CONDITIONS}
Nous soutenons la modernisation du marché par l'IA car elle apporte plus de sécurité (fin de la fraude aux cartes), un service disponible jour et nuit (agent conversationnel) et une meilleure organisation des agents (prédiction de l'affluence). \textbf{Mais} nous posons trois conditions : (1) le \textbf{consentement} des vendeurs doit être demandé et une \textbf{alternative sans photo du visage} (badge) doit exister ; (2) un \textbf{contrôle humain} doit vérifier toute décision importante (refus d'accès, alerte) car l'IA peut se tromper ; (3) les \textbf{données personnelles} doivent être protégées et le \textbf{coût total} (installation + maintenance) garanti par écrit. Sans ces garanties, le projet menace la vie privée et engage l'argent de la commune sans protection.
\end{quote}

\paragraph{Position « CONTRE » (exemple) :}
\begin{quote}
\textbf{Avis du Groupe 1 -- CONTRE}
Nous refusons ce projet dans sa forme actuelle. Le coût total (42 millions + 3,5 millions par an) est très élevé pour une mairie qui a d'autres priorités : latrines, eau potable, toitures du marché. Les photos de 1\,200 vendeurs seront stockées sur un serveur à l'étranger, hors de tout contrôle local. L'agent conversationnel peut donner de faux tarifs. La prédiction, fausse à $12\%$ près, ne justifie pas de déplacer des agents. L'argent public doit d'abord servir aux besoins essentiels.
\end{quote}

\textbf{Étape 4 -- Restitution collective (tableau de synthèse au tableau)}

L'enseignant construit ce tableau avec les apports de tous les groupes :

\begin{center}
\begin{tabularx}{\linewidth}{|X|X|X|X|}
\hline
\rowcolor{popblueL}
\textbf{Avantages cités} & \textbf{Risques cités} & \textbf{Conditions exigées} & \textbf{Positions finales} \\
\hline
Sécurité des accès (fin de la fraude) & Atteinte à la vie privée (photos) & 1. Consentement + alternative sans visage (badge) & 6 groupes : POUR SOUS CONDITIONS \\
Service 24h/24 (agent conversationnel) & Erreurs de reconnaissance & 2. Contrôle humain des décisions importantes & 2 groupes : CONTRE \\
Meilleure organisation des agents & Réponses fausses de l'IA & 3. Vérification des informations données & 1 groupe : POUR \\
Image moderne de la commune & Coût et dépendance au fournisseur & 4. Protection et hébergement maîtrisé des données & \\
Réponses en deux langues & Exclusion des personnes sans téléphone & 5. Budget de maintenance garanti 3 ans & \\
\hline
\end{tabularx}
\end{center}

\textbf{Commentaires pédagogiques pour l'enseignant :}
\begin{itemize}
    \item \textbf{Valoriser} les groupes qui citent la protection des \cle{données personnelles} et le \cle{consentement} : une photo du visage est une donnée personnelle qu'on ne peut pas utiliser sans accord.
    \item \textbf{Corriger} la confusion fréquente : « L'IA décide » $\rightarrow$ Non, l'IA \emph{propose} ou \emph{automatise} selon des règles définies par des humains. Insister sur le \cle{contrôle humain} des décisions importantes.
    \item \textbf{Préciser} qu'un agent conversationnel ne « comprend » pas comme un humain : il produit des réponses probables à partir de ses données d'apprentissage, et peut donc se tromper. D'où la règle d'or : \emph{toujours vérifier une information importante donnée par une IA}.
    \item \textbf{Noter} l'argument de la consommation d'énergie : serveurs et caméras fonctionnent en permanence, ce qui augmente la facture d'électricité de la mairie.
\end{itemize}
\end{exoresolu}

\courssub{Bilan de l'activité}
Cette activité a permis de vérifier l'acquisition des compétences visées par la leçon dans une situation complexe et réaliste :
\begin{itemize}
    \item \textbf{Savoirs mobilisés :} les élèves ont reconnu différents usages de l'IA (reconnaissance d'images, agent conversationnel, prédiction) et ont cité des avantages (automatisation, disponibilité, rapidité) et des inconvénients (erreurs, vie privée, coût, exclusion) conformes au programme.
    \item \textbf{Savoir-faire « Rédiger et justifier » :} la rédaction d'une consigne structurée (contexte, rôle, tâche, forme) a montré la capacité à formuler une demande précise à une IA. L'avis argumenté (10 lignes, 3 arguments minimum, vocabulaire du cours : \emph{consentement, contrôle humain, vie privée, vérification}) démontre une construction logique maîtrisée.
    \item \textbf{Savoir-faire « Appliquer à la vie courante » :} l'ancrage au marché de Nkolbisson (budget en FCFA, langues, réalités locales) a sorti la réflexion de l'abstraction. Les élèves ont perçu l'IA non comme de la magie mais comme un \textbf{choix technique et citoyen} qui engage l'argent public et les droits des personnes.
    \item \textbf{Savoir-être (éthique numérique) :} le débat et le vote ont révélé une prise de conscience : \emph{l'efficacité technique ne suffit pas}. La position « POUR SOUS CONDITIONS », majoritaire, traduit une maturité numérique : accepter l'innovation \textbf{si et seulement si} les garde-fous (consentement, contrôle humain, protection des données, alternative pour les non-équipés) sont garantis.
    \item \textbf{Piste d'approfondissement (projet) :} concevoir une affiche « Vos droits face à l'IA au marché » (droit de donner ou refuser son accord, droit de vérifier une information, droit de parler à un humain) à destination des vendeurs, avec des pictogrammes simples.
\end{itemize}

\newpage

\coursec{Méthodes et exercices résolus}

\coursec{Leçon N01 : Découverte d’un environnement d’IA}

\courssub{Généralités sur l'Intelligence Artificielle}

\begin{definition}[Intelligence Artificielle (IA)]
L'\cle{Intelligence Artificielle} est un ensemble de théories et de techniques visant à créer des machines (programmes, robots) capables de \textbf{simuler} l'intelligence humaine : apprentissage, raisonnement, perception, compréhension du langage naturel, résolution de problèmes.
\end{definition}

\begin{propriete}[Distinction fondamentale : IA faible vs IA forte]
\begin{itemize}
    \item \textbf{IA faible (ou étroite) :} Spécialisée dans \textbf{une tâche précise} (reconnaissance vocale, recommandation, jeu d'échecs, traduction). Elle ne possède ni conscience, ni compréhension réelle, ni capacité de transfert vers d'autres domaines. \emph{C'est l'unique forme d'IA déployée aujourd'hui.}
    \item \textbf{IA forte (ou générale / AGI) :} Hypothétique. Capacité à \textbf{comprendre, apprendre et appliquer toute tâche intellectuelle} qu'un humain peut faire, avec conscience de soi et adaptation autonome à des situations inédites sans réentraînement. \emph{N'existe pas à ce jour.}
\end{itemize}
\end{propriete}

\begin{methode}[Définir l'Intelligence Artificielle et distinguer ses types]
\textbf{Objectif :} Savoir expliquer ce qu'est l'IA et faire la différence entre IA faible et IA forte.
\begin{enumerate}
    \item \textbf{Repérer la définition :} L'IA est un ensemble de théories et de techniques visant à créer des machines capables de simuler l'intelligence humaine (apprentissage, raisonnement, perception).
    \item \textbf{Identifier le type d'IA :}
    \begin{itemize}
        \item \textbf{IA faible (étroite) :} Spécialisée dans une tâche précise (reconnaissance vocale, recommandation, jeu d'échecs). C'est l'IA d'aujourd'hui.
        \item \textbf{IA forte (générale) :} Capacité à comprendre, apprendre et appliquer toute tâche intellectuelle qu'un humain peut faire. N'existe pas encore.
    \end{itemize}
    \item \textbf{Citer des exemples concrets :} Associer chaque exemple au bon type (ex. : assistant vocal = IA faible ; robot conscient de films = IA forte).
    \item \textbf{Justifier :} Expliquer \emph{pourquoi} un système est une IA faible (pas de conscience, pas de compréhension réelle, entraîné sur des données spécifiques).
\end{enumerate}
\end{methode}

\begin{exoresolu}[Identifier le type d'IA dans des situations camerounaises]
\textbf{Énoncé :} Pour chacun des systèmes suivants, indiquez s'il relève de l'IA faible ou de l'IA forte et justifiez votre réponse en une phrase.
\begin{enumerate}
    \item L'algorithme de recommandation de vidéos sur TikTok.
    \item Le système de détection de fraude bancaire utilisé par une banque à Douala.
    \item Un robot humanoïde capable de discuter de philosophie, de cuisiner le ndolé et de résoudre une équation du second degré sans avoir été programmé pour chacune de ces tâches.
    \item Le correcteur orthographique automatique de votre traitement de texte.
\end{enumerate}

\tcblower
\textbf{Solution détaillée :}
\begin{enumerate}
    \item \textbf{IA faible.} L'algorithme analyse vos habitudes de visionnage pour prédire ce qui vous plaira ; il ne « comprend » pas le contenu des vidéos, il applique un modèle statistique entraîné sur une tâche unique : la recommandation.
    \item \textbf{IA faible.} Le système repère des schémas anormaux dans les transactions (montants, lieux, heures) grâce à un apprentissage supervisé sur des fraudes passées. Il ne raisonne pas comme un enquêteur humain.
    \item \textbf{IA forte (hypothétique).} Ce robot manifesterait une intelligence générale : transfert de compétences, adaptation à des tâches inédites sans réentraînement, compréhension du sens. Aucun système actuel n'atteint ce niveau.
    \item \textbf{IA faible.} Le correcteur compare vos mots à un dictionnaire et applique des règles grammaticales statistiques ; il ne « sait » pas ce que vous voulez dire, il prédit la correction la plus probable.
\end{enumerate}
\textbf{Conclusion :} Tous les systèmes d'IA déployés aujourd'hui au Cameroun et dans le monde sont de l'IA faible. L'IA forte reste un sujet de recherche fondamentale.
\end{exoresolu}

\courssub{Importance et usages de l'IA dans la vie courante}

\begin{methode}[Repérer les usages de l'IA dans la vie courante et au Cameroun]
\textbf{Objectif :} Identifier des applications concrètes de l'IA dans son environnement immédiat (école, maison, administration, commerce).
\begin{enumerate}
    \item \textbf{Lister les domaines :} Santé, éducation, transport, agriculture, commerce, administration, divertissement, sécurité.
    \item \textbf{Associer une technologie IA :}
    \begin{itemize}
        \item Vision par ordinateur (reconnaissance d'images, caméras de surveillance intelligentes).
        \item Traitement automatique du langage naturel (chatbots, traduction, assistants vocaux).
        \item Apprentissage automatique / \emph{Machine Learning} (prédiction, classification, recommandation).
        \item Robotique (drones agricoles, automate industriel).
    \end{itemize}
    \item \textbf{Contextualiser au Cameroun :} Penser aux exemples locaux : diagnostic médical par téléphone (applications de télémédecine), prévisions météorologiques pour les agriculteurs, chatbots des administrations (impôts, CNPS), reconnaissance faciale aux frontières, optimisation des itinéraires pour les motos-taxis.
    \item \textbf{Expliquer le gain :} Gain de temps, réduction des coûts, accès à des services éloignés, aide à la décision.
\end{enumerate}
\end{methode}

\begin{exoresolu}[Usages de l'IA dans un lycée de Yaoundé]
\textbf{Énoncé :} Le proviseur du Lycée Bilingue d'Etoug-Ebe souhaite moderniser l'établissement. Il vous demande de proposer \textbf{trois} applications concrètes de l'IA pour améliorer la vie scolaire. Pour chacune, précisez : (a) le problème résolu, (b) la technologie IA utilisée, (c) un avantage et une limite.

\tcblower
\textbf{Solution détaillée :}

\textbf{Proposition 1 : Suivi automatisé de l'assiduité par reconnaissance faciale}
\begin{itemize}
    \item \textbf{Problème :} Perte de temps à l'appel manuel, fraudes (émargement par un camarade).
    \item \textbf{Technologie :} Vision par ordinateur (reconnaissance faciale en temps réel à l'entrée des salles).
    \item \textbf{Avantage :} Fiabilité, gain de temps pédagogique, statistiques automatiques pour l'administration.
    \item \textbf{Limite :} Coût d'installation (caméras, serveur), risque sur la vie privée des mineurs (consentement parental obligatoire), panne électrique / réseau.
\end{itemize}

\textbf{Proposition 2 : Chatbot d'orientation scolaire et administrative}
\begin{itemize}
    \item \textbf{Problème :} Élèves désorientés sur les filières, procédures d'inscription complexes, secrétariat saturé.
    \item \textbf{Technologie :} Traitement automatique du langage naturel (NLP) — agent conversationnel en français / anglais / langues locales.
    \item \textbf{Avantage :} Disponible 24/7, réponses instantanées, décharge le personnel administratif.
    \item \textbf{Limite :} Ne remplace pas le conseil humain personnalisé ; peut halluciner (donner de fausses dates d'examen) si mal entraîné ; nécessite une maintenance continue.
\end{itemize}

\textbf{Proposition 3 : Système de détection précoce du décrochage scolaire}
\begin{itemize}
    \item \textbf{Problème :} Élèves qui quittent l'école sans alerte précoce.
    \item \textbf{Technologie :} \emph{Machine Learning} (classification supervisée) sur les notes, absences, retards, comportement.
    \item \textbf{Avantage :} Permet une intervention ciblée du conseiller d'orientation / assistant social \emph{avant} la rupture.
    \item \textbf{Limite :} Risque de biais (élèves de zones rurales pénalisés par le manque de données), étiquetage stigmatisant, besoin de données historiques de qualité.
\end{itemize}

\textbf{Conclusion :} L'IA apporte des solutions concrètes mais chaque déploiement doit intégrer coûts, éthique, infrastructure et formation du personnel.
\end{exoresolu}

\courssub{Avantages et inconvénients de l'IA}

\begin{methode}[Analyser les avantages et les inconvénients d'une solution d'IA]
\textbf{Objectif :} Mener une analyse critique équilibrée (coûts / bénéfices, risques / opportunités) pour prendre une décision éclairée.
\begin{enumerate}
    \item \textbf{Cadrer le sujet :} Définir précisément l'application IA étudiée (ex. : « utilisation de ChatGPT pour rédiger mes dissertations »).
    \item \textbf{Lister les avantages (bénéfices) :} Efficacité, rapidité, accessibilité, personnalisation, coût réduit, disponibilité 24/7, appui à la créativité.
    \item \textbf{Lister les inconvénients (risques / limites) :} Hallucinations (erreurs plausibles), biais discriminatoires, dépendance technologique, perte de compétences propres, vie privée / données, coût énergétique, inégalité d'accès (fracture numérique).
    \item \textbf{Classer par dimension :} Technique, économique, sociale, éthique, environnementale, juridique.
    \item \textbf{Synthétiser :} Rédiger une conclusion nuancée : « L'IA est un levier puissant à condition de... »
\end{enumerate}
\end{methode}

\begin{exoresolu}[Analyse critique : un élève utilise une IA générative pour ses devoirs]
\textbf{Énoncé :} Marie, élève en 3ème à Bamenda, utilise une IA générative (type ChatGPT) pour faire ses devoirs de français, d'histoire et de mathématiques. Réalisez une analyse structurée (avantages / inconvénients) et donnez votre avis argumenté sur cette pratique.

\tcblower
\textbf{Solution détaillée :}

\textbf{Tableau d'analyse :}

\begin{center}
\begin{tabularx}{\linewidth}{|l|X|X|}
\hline
\rowcolor{popblueL}
\textbf{Dimension} & \textbf{Avantages (Opportunités)} & \textbf{Inconvénients (Risques / Limites)} \\
\hline
Pédagogique & Explications instantanées, reformulations claires, exemples variés, disponibilité permanente. & \textbf{Hallucinations :} réponses fausses mais plausibles (dates, formules). \textbf{Passivité :} l'élève ne construit pas son raisonnement. Perte d'entraînement à l'écriture et au calcul mental. \\
\hline
Éthique & Accès à la connaissance pour tous (réduction des inégalités si connexion dispo). & \textbf{Triche / Plagiat :} travail non personnel. Violation du règlement intérieur. \textbf{Biais :} réponses stéréotypées (genre, origine). \\
\hline
Technique & Interface simple, multilingue, gratuite (version de base). & Dépendance à la connexion Internet et aux serveurs étrangers. Modèle « boîte noire » : on ne sait pas comment la réponse est produite. \\
\hline
Économique & Gratuit ou peu cher vs cours particuliers. & Coût caché : données personnelles collectées, consommation énergétique massive des centres de données. \\
\hline
Social & Aide les élèves en difficulté ou timides. & Creuse la fracture numérique (ceux qui n'ont pas de smartphone / data sont exclus). Isolement : moins d'échanges prof-élève / élève-élève. \\
\hline
\end{tabularx}
\end{center}

\textbf{Avis argumenté :}
L'utilisation de l'IA générative par Marie est \textbf{bénéfique uniquement en mode « tuteur »} (demander une explication, un plan, une correction \emph{après} avoir essayé seule) et \textbf{nuisible en mode « substitut »} (copier-coller la réponse).
\textbf{Recommandation :} Le professeur doit encadrer l'usage : autoriser l'IA pour la recherche d'idées, la reformulation, la vérification ; interdire la production finale soumise pour notation. L'élève doit citer l'outil utilisé (transparence).
\end{exoresolu}

\courssub{Utilisation responsable de l'IA}

\begin{methode}[Appliquer les principes d'une utilisation responsable de l'IA]
\textbf{Objectif :} Respecter les règles éthiques, légales et de bon sens lors de l'usage d'un système d'IA (protection des données, vigilance aux biais, transparence, esprit critique).
\begin{enumerate}
    \item \textbf{Protéger ses données personnelles :} Ne jamais saisir d'informations sensibles (numéro CNI, mot de passe, coordonnées bancaires, notes médicales, photos d'enfants) dans une IA publique. Utiliser des comptes anonymes si possible.
    \item \textbf{Vérifier les sources (esprit critique) :} Croiser la réponse de l'IA avec une source fiable (site officiel, manuel, professeur). L'IA n'est pas une source primaire.
    \item \textbf{Repérer les biais :} Se demander : « Cette réponse favorise-t-elle un groupe ? Ignore-t-elle un contexte local (camerounais) ? » Signaler les contenus haineux, discriminatoires ou illégaux.
    \item \textbf{Citer l'IA (transparence / intégrité académique) :} Mentionner l'outil utilisé (ex. : « Ce plan a été suggéré par ChatGPT-4o le 12/10/2025 »). Ne pas faire passer une production IA pour la sienne.
    \item \textbf{Respecter le droit d'auteur :} Ne pas demander à l'IA de reproduire une œuvre protégée (texte intégral d'un roman, code source propriétaire).
    \item \textbf{Limiter l'impact environnemental :} Éviter les requêtes inutiles ; préférer des modèles légers / locaux pour des tâches simples.
\end{enumerate}
\end{methode}

\begin{exoresolu}[Cas pratique : utilisation responsable dans un cybercafé de Douala]
\textbf{Énoncé :} Vous gérez un cybercafé à Douala. Un client veut utiliser une IA en ligne pour : (1) rédiger son CV avec son vrai nom, adresse, numéro de téléphone ; (2) générer une fausse attestation de travail ; (3) traduire un document médical confidentiel. En tant que gestionnaire responsable, que lui conseillez-vous pour chaque cas ? Justifiez par un principe éthique ou légal.

\tcblower
\textbf{Solution détaillée :}

\textbf{Cas 1 : CV avec données réelles}
\begin{itemize}
    \item \textbf{Conseil :} « Utilisez l'IA pour la \textbf{structure} et la \textbf{formulation} (modèle de phrases), mais \textbf{ne collez pas} vos données personnelles réelles dans la conversation. Remplacez-les par des placeholders [NOM], [TÉLÉPHONE] et mettez les vraies infos dans votre traitement de texte \emph{après} copier-coller. »
    \item \textbf{Principe :} \textbf{Protection des données personnelles (RGPD / Loi camerounaise sur la cybercriminalité).} Les IA publiques entraînent souvent sur les données saisies ; vos infos peuvent réapparaître chez d'autres utilisateurs ou être vendues.
\end{itemize}

\textbf{Cas 2 : Fausse attestation de travail}
\begin{itemize}
    \item \textbf{Conseil :} \textbf{Refusez catégoriquement.} « Je ne peux pas vous aider à produire un faux document. C'est un délit de faux et usage de faux (Code pénal camerounais, art. 123). L'IA ne doit pas servir à la fraude. »
    \item \textbf{Principe :} \textbf{Légalité et intégrité.} L'IA amplifie la capacité à falsifier ; l'utilisateur en porte la responsabilité pénale. Le cybercafé peut être complice s'il facilite.
\end{itemize}

\textbf{Cas 3 : Document médical confidentiel}
\begin{itemize}
    \item \textbf{Conseil :} \textbf{Interdisez l'envoi vers une IA publique.} « Ce document relève du secret médical. Utilisez un traducteur \textbf{hors ligne} (ex. : LibreTranslate en local, ou dictionnaire) ou faites appel à un traducteur assermenté humain sous contrat de confidentialité. »
    \item \textbf{Principe :} \textbf{Confidentialité / Secret professionnel / Souveraineté des données de santé.} Envoyer un dossier médical sur un serveur américain ou chinois viole la déontologie et expose le patient.
\end{itemize}

\textbf{Conclusion :} Le gestionnaire de cybercafé a un rôle d'éducateur numérique : il doit afficher une charte « Usage responsable de l'IA » et former ses clients aux bons réflexes.
\end{exoresolu}

\courssub{Apprendre à échanger avec une IA}

\begin{methode}[Rédiger un prompt efficace (Ingénierie de prompt de base)]
\textbf{Objectif :} Obtenir une réponse pertinente, structurée et adaptée en formulant une instruction claire, contextualisée et contrainte.
\begin{enumerate}
    \item \textbf{Donner un RÔLE (Persona) :} « Agis comme un professeur de mathématiques de 3ème au Cameroun... »
    \item \textbf{Préciser la TÂCHE :} « Explique-moi la factorisation par mise en évidence... »
    \item \textbf{Fixer le CONTEXTE / PUBLIC :} « Pour un élève de 14 ans qui a des difficultés en algèbre, programme MINESEC 3ème. »
    \item \textbf{Imposer le FORMAT :} « Réponse en 3 étapes numérotées, avec un exemple numérique complet, vocabulaire simple, un exercice d'application à la fin. »
    \item \textbf{Ajouter des CONTRAINTES :} « N'utilise pas la formule du discriminant. Donne un exemple avec des nombres entiers positifs. Ton encourageant. »
    \item \textbf{Itérer (affiner) :} Si la réponse ne convient pas : « Reformule plus simplement », « Donne un autre exemple avec la vie courante (achat au marché) », « Corrige l'erreur à l'étape 2 ».
\end{enumerate}
\textbf{Mnémotechnique :} \cle{R-T-C-F-C-I} (Rôle, Tâche, Contexte, Format, Contraintes, Itération).
\end{methode}

\begin{exoresolu}[Rédaction d'un prompt pour réviser le BEPC]
\textbf{Énoncé :} Vous êtes en 3ème au Lycée de Nkongsamba. Vous voulez que l'IA vous génère une fiche de révision sur « Les puissances de 10 et notation scientifique » (programme Mathématiques 3ème). Rédigez un prompt complet en appliquant la méthode R-T-C-F-C-I.

\tcblower
\textbf{Solution détaillée :}

\begin{lstlisting}[language=text]
Agis comme un professeur de mathématiques de 3ème au Cameroun, expert du programme MINESEC et du BEPC.

TÂCHE : Crée une fiche de révision complète et structurée sur "Les puissances de 10 et la notation scientifique".

CONTEXTE : 
- Élève de 14 ans, série scientifique, préparation au BEPC session 2026.
- Niveau : initiation, pas de logarithmes, pas de chiffres significatifs complexes.
- Difficultés fréquentes : signe de l'exposant, passage écriture décimale <-> notation scientifique, opérations (+, -, ×, ÷).

FORMAT EXIGÉ :
1. Titre clair + rappel du programme officiel.
2. Définitions simples (puissance de 10, notation scientifique : a × 10^n avec 1 ≤ a < 10).
3. Tableau de correspondance : 10^3, 10^2, 10^1, 10^0, 10^-1, 10^-2, 10^-3 avec écriture décimale.
4. Règles de conversion (décimal → scientifique et inverse) en 3 étapes numérotées.
5. Règles des opérations (×, ÷) : on additionne/soustrait les exposants.
6. 3 exemples corrigés détaillés (un grand nombre, un petit nombre, une opération).
7. 4 exercices d'application progressifs (QCM, conversion, opération, problème concret).
8. Pièges à éviter (liste à puces).
9. Mnémotechnique pour retenir le sens du déplacement de la virgule.

CONTRAINTES :
- Vocabulaire du programme camerounais (ex. "partie entière", "partie décimale", "exposant").
- Exemples concrets camerounais : distance Yaoundé-Douala (km), diamètre d'un cheveu (mm), budget du MINEDUB (FCFA).
- Ton bienveillant, encourageant, "Tu y arrives !".
- Pas de LaTeX, texte brut lisible sur téléphone.
- Longueur max : 1 page A4 (≈ 500 mots).

ITÉRATION (si besoin) : 
- Si trop long : "Condense en supprimant les explications théoriques, garde l'essentiel BEPC."
- Si trop abstrait : "Ajoute un exemple avec le prix du kg de poisson au marché."
\end{lstlisting}

\textbf{Analyse du prompt :} Chaque composante R-T-C-F-C-I est remplie. Le contexte « Cameroun / BEPC / 3ème » garantit l'alignement programme. Les contraintes de format (numérotation, tableaux, pièges) produisent une fiche directement utilisable. L'itération prévue permet d'affiner sans recommencer.
\end{exoresolu}

\begin{methode}[Évaluer la réponse d'une IA : esprit critique et détection d'hallucinations]
\textbf{Objectif :} Ne pas croire aveuglément la sortie de l'IA ; vérifier la cohérence, la justesse factuelle et la pertinence.
\begin{enumerate}
    \item \textbf{Test de cohérence interne :} La réponse se contredit-elle ? Les étapes logiques s'enchaînent-elles ?
    \item \textbf{Test de plausibilité (bon sens) :} Le résultat est-il dans un ordre de grandeur réaliste ? (Ex. : population du Cameroun = 2,5 millions ? Non, ~28 millions).
    \item \textbf{Vérification croisée (Triangulation) :} Comparer avec \textbf{au moins une source fiable} : manuel scolaire, site officiel (.gov.cm, .edu), encyclopédie reconnue, professeur.
    \item \textbf{Repérer les signaux d'hallucination :} Citations inventées (articles de loi inexistants, auteurs fictifs), faux code (bibliothèques qui n'existent pas), dates erronées, raisonnement circulaire.
    \item \textbf{Tester les limites :} Poser une question piège dont on connaît la réponse (ex. : « Quelle est la capitale de la région du Nord-Ouest ? » → Bamenda. Si l'IA dit « Bamenda est au Nord » → erreur).
    \item \textbf{Décider :} Accepter / Corriger / Rejeter / Demander une source à l'IA (« Sur quoi te bases-tu ? Donne-moi l'URL exacte »).
\end{enumerate}
\end{methode}

\begin{exoresolu}[Détecter les erreurs d'une IA sur l'histoire du Cameroun]
\textbf{Énoncé :} Un élève demande à une IA : « Résume l'indépendance du Cameroun en 5 points clés. » L'IA répond :
\begin{enumerate}
    \item Le Cameroun a obtenu son indépendance le 1er janvier 1961 après une guerre de libération contre la France.
    \item Ahmadou Ahidjo a été le premier Premier ministre, puis président en 1962.
    \item Le pays s'appelait « République Fédérale du Cameroun » dès 1960.
    \item La réunification avec le Cameroun britannique a eu lieu en 1972 après un référendum.
    \item L'hymne national « Ô Cameroun, berceau de nos ancêtres » a été adopté en 1960 par le président Paul Biya.
\end{enumerate}
En vous basant sur vos connaissances d'Histoire-Géographie (programme 3ème), identifiez \textbf{au moins 4 erreurs} et corrigez-les.

\tcblower
\textbf{Solution détaillée :}

\textbf{Erreur 1 : Date et nature de l'indépendance.}
\begin{itemize}
    \item \textbf{IA :} « 1er janvier 1961 après une guerre contre la France. »
    \item \textbf{Réalité :} Indépendance du Cameroun français le \textbf{1er janvier 1960} (proclamée par Ahmadou Ahidjo), négociée, pas une guerre de libération type Algérie. Le 1er janvier 1961 est la date de l'indépendance du Cameroun britannique (Southern Cameroons) qui choisit de rejoindre le Cameroun (réunification).
\end{itemize}

\textbf{Erreur 2 : Titre d'Ahmadou Ahidjo.}
\begin{itemize}
    \item \textbf{IA :} « Premier Premier ministre, puis président en 1962. »
    \item \textbf{Réalité :} Ahidjo est \textbf{Premier ministre} du Cameroun français autonome (1958), puis \textbf{Président de la République} dès l'indépendance \textbf{1er janvier 1960}. Il n'y a pas de « premier ministre » après 1960 (régime présidentiel).
\end{itemize}

\textbf{Erreur 3 : Nom de l'État en 1960.}
\begin{itemize}
    \item \textbf{IA :} « République Fédérale du Cameroun dès 1960. »
    \item \textbf{Réalité :} 1960-1961 : \textbf{République du Cameroun} (État unitaire). La \textbf{République Fédérale du Cameroun} naît le \textbf{1er octobre 1961} après la réunification (deux États fédérés : Cameroun Oriental et Cameroun Occidental).
\end{itemize}

\textbf{Erreur 4 : Date de la réunification / fin du fédéralisme.}
\begin{itemize}
    \item \textbf{IA :} « Réunification en 1972 après un référendum. »
    \item \textbf{Réalité :} Réunification = \textbf{1er octobre 1961}. Le référendum du \textbf{20 mai 1972} met fin au fédéralisme → \textbf{République Unie du Cameroun} (puis République du Cameroun en 1984). L'IA confond réunification et unification.
\end{itemize}

\textbf{Erreur 5 (bonus) : Hymne national et Paul Biya.}
\begin{itemize}
    \item \textbf{IA :} « Adopté en 1960 par le président Paul Biya. »
    \item \textbf{Réalité :} L'hymne « Ô Cameroun » (paroles : René Jam Afane / musique : Samuel Minkio Bamba) est officiel dès \textbf{1960} sous \textbf{Ahmadou Ahidjo}. Paul Biya devient président en \textbf{1982}.
\end{itemize}

\textbf{Conclusion :} Cette réponse contient des hallucinations factuelles graves (mélange de dates, de régimes, de personnalités). Un élève qui l'apprendrait aurait 0/20 au BEPC. \textbf{Toujours croiser avec le manuel d'Histoire 3ème.}
\end{exoresolu}

\begin{attention}[Les 5 erreurs qui coûtent des points au BEPC]
\begin{enumerate}
    \item \textbf{Confondre IA faible et IA forte :} Dire « l'IA comprend » ou « l'IA pense » au lieu de « l'IA prédit / simule / traite ». \emph{Sanction :} Perte de points sur la définition (savoir essentiel).
    \item \textbf{Citer des exemples hors contexte camerounais sans adaptation :} « Voitures autonomes à Yaoundé » (infrastructure inexistante), « Chirurgie robotisée au district de santé » (équipement absent). \emph{Sanction :} Exemple non pertinent = 0 pt sur l'application concrète.
    \item \textbf{Lister des avantages sans inconvénients (ou inversement) :} Analyse unilatérale. Le programme exige « Avantages, inconvénients ET utilisation responsable ». \emph{Sanction :} Compétence « Analyser » non validée.
    \item \textbf{Ne pas protéger ses données dans un prompt :} Copier-coller son numéro CNI, mot de passe Facebook, ou bulletin de notes dans le chat. \emph{Sanction :} Savoir-être « Utilisation responsable » non acquis ; risque réel hors examen.
    \item \textbf{Accepter la réponse de l'IA sans vérification :} Recopier une définition fausse, une date inventée, un code bugué. \emph{Sanction :} Zéro à l'exercice pratique ; en vie réelle : échec au BEPC ou diffusion de fausses infos.
\end{enumerate}
\end{attention}

\newpage

\coursec{Exercices}

\courssub{Je vérifie mes connaissances}

\begin{exercice}[QCM : Applications relevant de l'IA]
Parmi les cinq applications suivantes, cochez celles qui relèvent de l'Intelligence Artificielle (au sens apprentissage automatique / traitement automatique du langage) et justifiez brièvement chaque choix.
\begin{enumerate}
    \item Une calculatrice scientifique qui résout une équation du second degré.
    \item Un assistant vocal (type Google Assistant ou Siri) qui répond à une question orale.
    \item Un tableur qui calcule la somme d'une colonne de nombres.
    \item Un service de traduction automatique (type Google Translate) qui traduit un texte du français vers l'anglais.
    \item Un jeu de dames classique où l'ordinateur suit des règles fixes programmées à l'avance (arbre de décision sans apprentissage).
\end{enumerate}
\end{exercice}

\begin{exercice}[Vrai ou Faux justifié]
Pour chaque affirmation, répondez par \textbf{Vrai} ou \textbf{Faux} et justifiez votre réponse en une phrase précise.
\begin{enumerate}
    \item L'IA pense et raisonne exactement comme un être humain.
    \item Une IA peut se tromper ou donner une réponse inexacte (hallucination, biais).
    \item Il est sans danger de saisir ses mots de passe personnels dans un chatbot public.
    \item L'apprentissage automatique (Machine Learning) permet à une IA d'améliorer ses performances sans être explicitement reprogrammée pour chaque nouvelle tâche.
\end{enumerate}
\end{exercice}

\begin{exercice}[Définitions et concepts clés]
Donnez une définition claire et concise (2 à 3 lignes) pour chacun des termes suivants :
\begin{enumerate}
    \item Intelligence Artificielle (IA)
    \item Apprentissage automatique (Machine Learning)
    \item Prompt (invite)
    \item Hallucination (dans le contexte des IA génératives)
    \item Biais algorithmique
\end{enumerate}
\end{exercice}

\begin{exercice}[Calcul direct : Tokenisation]
Un modèle de langage découpe le texte en \emph{tokens}. En moyenne, 1 token $\approx$ 0,75 mot en français (soit 1 mot $\approx$ 1,33 token).
\begin{enumerate}
    \item Combien de tokens environ pour un prompt de 120 mots ?
    \item Si la limite de contexte (fenêtre de contexte) est de 4096 tokens, quel est le nombre maximal de mots français que l'on peut théoriquement soumettre en une seule requête ?
\end{enumerate}
\end{exercice}

\courssub{Je m'entraîne}

\begin{exercice}[Amélioration de prompts : méthode Rôle–Tâche–Contexte–Format]
Les trois prompts ci-dessous sont maladroits. Réécrivez chacun en appliquant la structure \textbf{Rôle -- Tâche -- Contexte -- Format} pour obtenir une réponse de meilleure qualité.
\begin{enumerate}
    \item \texttt{« Fais-moi un CV. »}
    \item \texttt{« Explique-moi la photosynthèse. »}
    \item \texttt{« Écris une lettre de motivation pour un stage. »}
\end{enumerate}
Pour chaque proposition, indiquez explicitement : le \textbf{Rôle} attribué à l'IA, la \textbf{Tâche} précise, le \textbf{Contexte} apporté et le \textbf{Format} attendu.
\end{exercice}

\begin{exercice}[Tri : IA, pas IA, usage abusif]
Classez les huit situations suivantes dans le tableau ci-dessous en cochant la case correspondante. Justifiez en une phrase les cas classés « Usage abusif d'IA ».
\begin{center}
\begin{tabularx}{\linewidth}{|l|X|X|X|}
\hline
\textbf{Situation} & \textbf{IA} & \textbf{Pas IA} & \textbf{Usage abusif d'IA} \\ \hline
Reconnaissance faciale pour déverrouiller un téléphone & & & \\ \hline
Tableur calculant une moyenne & & & \\ \hline
Chatbot répondant aux questions d'un élève & & & \\ \hline
Calculatrice scientifique & & & \\ \hline
Traduction automatique d'un document administratif & & & \\ \hline
Jeu de dames classique (règles fixes, sans apprentissage) & & & \\ \hline
Assistant vocal programmant un réveil sur commande orale & & & \\ \hline
Création d'un deepfake pour usurper l'identité d'un camarade & & & \\ \hline
\end{tabularx}
\end{center}
\end{exercice}

\begin{exercice}[Étude de cas : IA et Mobile Money au Cameroun]
Lisez le texte suivant :

\begin{quote}
\og Au Cameroun, les opérateurs de Mobile Money (MTN MoMo, Orange Money) déploient des systèmes d'IA pour détecter les transactions frauduleuses en temps réel. Lorsqu'un utilisateur reçoit un SMS suspect lui demandant son code PIN, un chatbot officiel peut l'alerter instantanément. Cependant, des arnaqueurs utilisent aussi des IA génératives pour imiter la voix d'un proche (clonage vocal) et demander un transfert d'urgence. \fg
\end{quote}

À partir de ce texte, dégagez :
\begin{enumerate}
    \item Deux avantages de l'IA pour la sécurité du Mobile Money.
    \item Deux inconvénients ou risques liés à l'usage de l'IA dans ce contexte.
    \item Deux règles d'usage responsable que tout utilisateur devrait respecter pour se protéger.
\end{enumerate}
\end{exercice}

\begin{exercice}[Argumentation : L'IA pour les devoirs ?]
Rédigez un paragraphe d'environ 8 lignes répondant à la question : \og Faut-il autoriser les élèves à utiliser l'IA pour leurs devoirs ? \fg
Votre réponse doit contenir \textbf{au moins deux arguments} (un pour, un contre ou deux nuances) et une \textbf{conclusion justifiée} (votre position personnelle argumentée).
\end{exercice}

\courssub{Je me prépare au BEPC}

\begin{exercice}[Sujet type BEPC : L'IA dans l'éducation au Cameroun]
\textbf{Contexte :} Le Ministère des Enseignements Secondaires (MINESEC) envisage d'intégrer des outils d'IA générative dans les salles informatiques des lycées publics pour accompagner les élèves en algorithmique et en rédaction.

\textbf{Données :}
\begin{itemize}
    \item L'IA peut proposer des corrections de code en C ou des suggestions de structure pour une dissertation.
    \item Le risque de copier-coller sans compréhension est réel.
    \item La connexion Internet n'est pas stable dans tous les établissements.
    \item Les données des élèves (notes, identités) sont sensibles.
\end{itemize}

\textbf{Questions :}
\begin{enumerate}
    \item Citez deux avantages pédagogiques majeurs de l'intégration de l'IA en classe de 3ème.
    \item Identifiez deux risques éthiques ou techniques liés à cette intégration.
    \item Proposez trois règles concrètes d'usage responsable à inscrire dans la \emph{Charte Numérique de l'Élève}.
    \item Expliquez pourquoi la formation des enseignants à la rédaction de prompts (ingénierie de prompt) est indispensable avant le déploiement.
    \item Donnez un exemple de prompt efficace qu'un élève de 3ème pourrait utiliser pour s'entraîner à l'algorithmique (structure Rôle–Tâche–Contexte–Format).
\end{enumerate}
\end{exercice}

\begin{exercice}[Analyse critique : Détecter une hallucination et vérifier l'information]
\textbf{Contexte :} Un élève de 3ème demande à une IA générative : \og Donne-moi la biographie de l'inventeur camerounais du téléphone portable. \fg
L'IA répond : \og \textbf{Arthur Zang}, né en 1987 à Yaoundé, est l'inventeur du smartphone. Il a lancé le premier téléphone portable africain en 2010 sous la marque \textit{ZangPhone}. \fg

\textbf{Questions :}
\begin{enumerate}
    \item Identifiez \textbf{deux erreurs factuelles (hallucinations)} présentes dans cette réponse. (Rappel : Arthur Zang a inventé le \textbf{Cardiopad}, tablette médicale ; l'inventeur du téléphone portable est Martin Cooper chez Motorola en 1973).
    \item Quelle \textbf{règle d'usage responsable} l'élève doit-il appliquer systématiquement avant d'utiliser une information fournie par une IA dans un devoir ?
    \item Réécrivez un \textbf{prompt amélioré} qui demande à l'IA de citer ses sources et de ne répondre que si elle est certaine.
\end{enumerate}
\end{exercice}

\begin{exercice}[Éthique et protection des données : Charte IA au cybercafé scolaire]
\textbf{Contexte :} Le cybercafé du lycée met à disposition un accès à une IA générative en ligne (via navigateur). La directrice veut établir une \emph{Charte d'Usage de l'IA} affichée à côté de chaque poste.

\textbf{Questions :}
\begin{enumerate}
    \item Citez \textbf{trois types de données personnelles} qu'un élève ne doit \textbf{jamais} saisir dans la zone de dialogue (prompt) d'une IA publique.
    \item Pourquoi est-il déconseillé d'utiliser un \textbf{compte unique partagé} par toute la classe pour accéder à l'IA ? (Pensez à l'historique, la confidentialité et la personnalisation).
    \item Rédigez \textbf{deux règles claires} pour la \emph{Charte d'Usage de l'IA} :
    \begin{itemize}
        \item Règle 1 : Transparence (citation de l'IA).
        \item Règle 2 : Esprit critique (vérification).
    \end{itemize}
    \item L'établissement envisage d'héberger un modèle de langage \textbf{localement} (sur un serveur de l'école, sans Internet) plutôt que d'utiliser une API publique. Donnez \textbf{deux avantages} majeurs de cette solution pour la protection des données des mineurs.
\end{enumerate}
\end{exercice}

\newpage

\coursec{Corrigés des exercices}

\begin{corrige}[Exercice 1 — QCM : Applications relevant de l'IA]
\textbf{1. Calculatrice scientifique : NON.} Elle applique des formules mathématiques fixes (algorithme déterministe) : aucune donnée n'est apprise, aucun langage naturel n'est traité.

\textbf{2. Assistant vocal (Google Assistant, Siri) : OUI.} Il combine la reconnaissance vocale (transformer la voix en texte), le traitement automatique du langage naturel (comprendre la question) et la synthèse vocale : ce sont des domaines typiques de l'IA.

\textbf{3. Tableur calculant une somme : NON.} La fonction \texttt{SOMME} exécute une instruction explicite et déterministe : même entrée, même résultat, sans apprentissage.

\textbf{4. Traduction automatique : OUI.} Les services modernes (Google Translate) reposent sur des réseaux de neurones entraînés sur des millions de textes bilingues : c'est du traitement automatique du langage par apprentissage.

\textbf{5. Jeu de dames à règles fixes : NON (au sens moderne).} L'ordinateur explore un arbre de décision programmé à l'avance, sans apprendre de ses parties. C'est de l'informatique classique (parfois appelée « IA symbolique » historique), mais pas de l'IA au sens apprentissage automatique / traitement du langage demandé ici.
\end{corrige}

\begin{attention}[Point de vigilance]
Le critère à citer dans la justification est : \emph{l'application apprend-elle à partir de données ou traite-elle le langage naturel ?} Si elle ne fait qu'exécuter des règles fixes, ce n'est pas de l'IA au sens du cours.
\end{attention}

\begin{corrige}[Exercice 2 — Vrai ou Faux justifié]
\textbf{1. FAUX.} L'IA simule certaines tâches cognitives (reconnaître, prédire, générer) grâce à des calculs statistiques sur des données ; elle n'a ni conscience, ni émotions, ni compréhension humaine du monde.

\textbf{2. VRAI.} Une IA générative peut produire des informations fausses mais présentées avec assurance (hallucinations), et ses réponses peuvent refléter les biais présents dans ses données d'entraînement.

\textbf{3. FAUX.} Tout ce qui est saisi dans un chatbot public peut être stocké, analysé ou divulgué : les mots de passe, codes PIN et données personnelles ne doivent jamais y être saisis.

\textbf{4. VRAI.} C'est le principe même du \cle{Machine Learning} : le système ajuste ses paramètres à partir d'exemples (données) et améliore ses performances sans qu'un programmeur réécrive le code pour chaque cas nouveau.
\end{corrige}

\begin{corrige}[Exercice 3 — Définitions et concepts clés]
\textbf{1. \cle{Intelligence Artificielle (IA)} :} ensemble de techniques informatiques permettant à une machine d'accomplir des tâches qui demandent habituellement de l'intelligence humaine : comprendre le langage, reconnaître des images, prendre des décisions, générer du contenu.

\textbf{2. \cle{Apprentissage automatique (Machine Learning)} :} branche de l'IA où la machine améliore ses performances sur une tâche en analysant de grandes quantités de données (exemples), sans être explicitement reprogrammée pour chaque cas.

\textbf{3. \cle{Prompt (invite)} :} message texte saisi par l'utilisateur pour donner une instruction ou poser une question à une IA ; sa qualité (rôle, tâche, contexte, format) conditionne fortement la qualité de la réponse.

\textbf{4. \cle{Hallucination} :} réponse fausse, inventée ou incohérente produite par une IA générative, mais formulée de manière fluide et convaincante, ce qui la rend difficile à détecter sans vérification.

\textbf{5. \cle{Biais algorithmique} :} erreur systématique ou discrimination dans les résultats d'une IA, provenant de données d'entraînement non représentatives ou de choix de conception (ex. un système de reconnaissance faciale moins fiable sur certains types de visages).
\end{corrige}

\begin{corrige}[Exercice 4 — Calcul direct : Tokenisation]
\textbf{Données :} 1 mot $\approx$ 1,33 token ; 1 token $\approx$ 0,75 mot.

\textbf{1. Nombre de tokens pour 120 mots :}
\[
120 \text{ mots} \times 1{,}33 \text{ token/mot} = 159{,}6 \approx 160 \text{ tokens}.
\]

\textbf{2. Nombre maximal de mots pour 4096 tokens :}
\[
\frac{4096 \text{ tokens}}{1{,}33 \text{ token/mot}} \approx 3079{,}7 \approx 3080 \text{ mots}.
\]
Vérification avec l'autre conversion : $4096 \times 0{,}75 = 3072$ mots. Les deux résultats sont cohérents (l'écart vient de l'arrondi de 1,33) ; on retient \textbf{environ 3 000 à 3 080 mots}.
\end{corrige}

\begin{attention}[Point de vigilance]
Attention au sens de la conversion : pour passer des mots aux tokens on \emph{multiplie} par 1,33 ; pour passer des tokens aux mots on \emph{divise} par 1,33 (ou on multiplie par 0,75). Inverser le sens est l'erreur classique.
\end{attention}

\begin{corrige}[Exercice 5 — Amélioration de prompts : méthode Rôle–Tâche–Contexte–Format]
\textbf{1. CV :}
\begin{quote}
\og Agis comme un conseiller en orientation professionnelle. Rédige un modèle de CV pour un élève de 3ème au Cameroun, sans expérience professionnelle, qui cherche un stage de découverte dans une entreprise de Douala. Présente le CV en 4 rubriques : État civil, Formation, Compétences, Centres d'intérêt, sous forme de liste structurée. \fg
\end{quote}
Rôle : conseiller en orientation ; Tâche : rédiger un modèle de CV ; Contexte : élève de 3ème, sans expérience, stage à Douala ; Format : 4 rubriques en liste structurée.

\textbf{2. Photosynthèse :}
\begin{quote}
\og Agis comme un professeur de SVT de collège. Explique la photosynthèse à un élève de 3ème qui prépare le BEPC, en utilisant un vocabulaire simple et un exemple concret (une plante verte au soleil). Rédige environ 10 lignes avec une petite conclusion et une définition des mots clés en gras. \fg
\end{quote}
Rôle : professeur de SVT ; Tâche : expliquer la photosynthèse ; Contexte : élève de 3ème préparant le BEPC ; Format : 10 lignes, mots clés définis, conclusion.

\textbf{3. Lettre de motivation :}
\begin{quote}
\og Agis comme un rédacteur professionnel. Écris une lettre de motivation pour un stage d'observation d'une semaine dans une banque à Yaoundé, pour une élève de 3ème sérieuse et ponctuelle, intéressée par l'informatique. Respecte la forme classique d'une lettre administrative : en-tête, objet, corps en 3 paragraphes, formule de politesse. \fg
\end{quote}
Rôle : rédacteur professionnel ; Tâche : écrire une lettre de motivation ; Contexte : stage d'observation, banque à Yaoundé, profil de l'élève ; Format : lettre administrative en 3 paragraphes.
\end{corrige}

\begin{methode}[Ce que le correcteur attend]
\begin{enumerate}
    \item Identifier le \cle{Rôle} (ex. : conseiller, professeur, rédacteur).
    \item Définir la \cle{Tâche} précise (ex. : rédiger un CV, expliquer une notion, écrire une lettre).
    \item Préciser le \cle{Contexte} (niveau, lieu, profil, contraintes).
    \item Imposer le \cle{Format} de la réponse (rubriques, nombre de lignes, structure administrative).
\end{enumerate}
Un prompt incomplet (sans contexte ni format) reste un prompt faible, même formulé poliment.
\end{methode}

\begin{corrige}[Exercice 6 — Tri : IA, pas IA, usage abusif]
\begin{center}
\begin{tabularx}{\linewidth}{|X|c|c|c|}
\hline
\rowcolor{popblueL}
\textbf{Situation} & \textbf{IA} & \textbf{Pas IA} & \textbf{Usage abusif} \\ \hline
Reconnaissance faciale (déverrouillage) & X & & \\ \hline
Tableur calculant une moyenne & & X & \\ \hline
Chatbot répondant à un élève & X & & \\ \hline
Calculatrice scientifique & & X & \\ \hline
Traduction automatique d'un document & X & & \\ \hline
Jeu de dames à règles fixes & & X & \\ \hline
Assistant vocal programmant un réveil & X & & \\ \hline
Deepfake pour usurper une identité & X & & X \\ \hline
\end{tabularx}
\end{center}

\textbf{Justification du cas « Usage abusif » :} le \cle{deepfake} utilise bien une IA (génération d'image/voix par apprentissage), mais il vise à \emph{usurper l'identité} d'un camarade : c'est une atteinte à la dignité et à l'image d'autrui, punie par la loi (cybercriminalité). L'outil est légitime, l'usage est abusif et illégal.
\end{corrige}

\begin{attention}[Point de vigilance]
Une même situation peut être cochée dans deux colonnes : le deepfake est à la fois « IA » et « usage abusif ». Le tableau ne demande pas un classement exclusif pour cette ligne.
\end{attention}

\begin{corrige}[Exercice 7 — Étude de cas : IA et Mobile Money au Cameroun]
\textbf{1. Deux avantages pour la sécurité :}
\begin{itemize}
    \item \textbf{Détection en temps réel} des transactions frauduleuses : l'IA analyse des milliers d'opérations par seconde et repère les comportements anormaux, impossible à faire manuellement.
    \item \textbf{Alerte instantanée} de l'utilisateur via un chatbot officiel en cas de SMS suspect (hameçonnage du code PIN), ce qui prévient les vols avant qu'ils n'aient lieu.
\end{itemize}

\textbf{2. Deux inconvénients ou risques :}
\begin{itemize}
    \item Les fraudeurs utilisent eux aussi l'IA : le \textbf{clonage vocal} permet d'imiter la voix d'un proche pour demander un transfert d'urgence (arnaque crédible).
    \item Risque de \textbf{faux positifs} (une transaction légitime bloquée par erreur) et de confiance excessive : l'utilisateur peut croire que « l'IA le protège » et baisser sa vigilance.
\end{itemize}

\textbf{3. Deux règles d'usage responsable :}
\begin{itemize}
    \item Ne \textbf{jamais communiquer son code PIN} par SMS, appel ou chatbot, même si le message semble officiel : aucun opérateur ne le demande.
    \item \textbf{Vérifier par un autre canal} (rappeler directement le proche ou le numéro officiel de l'opérateur) toute demande urgente d'argent, surtout si la voix ou le message semble étrange.
\end{itemize}
\end{corrige}

\begin{corrige}[Exercice 8 — Argumentation : L'IA pour les devoirs ?]
\textbf{Exemple de paragraphe attendu (environ 8 lignes) :}

\og L'utilisation de l'IA pour les devoirs présente un avantage réel : elle peut expliquer une notion incomprise, proposer des exemples supplémentaires ou corriger un exercice, comme un tuteur disponible à toute heure, ce qui aide surtout les élèves n'ayant pas de soutien à la maison. Cependant, recopier une réponse d'IA sans la comprendre supprime l'effort d'apprentissage : l'élève obtient une bonne note mais ne construit aucune compétence, ce qui le pénalisera au BEPC où l'IA n'est pas autorisée. De plus, l'IA peut se tromper (hallucinations) et transmettre des erreurs. En conclusion, je pense qu'il faut autoriser l'IA comme \emph{outil d'aide et d'entraînement} — pour comprendre, vérifier, s'exercer — mais l'interdire comme \emph{machine à rédiger à la place de l'élève}, en citant toujours son utilisation. \fg

\textbf{Barème indicatif (sur 4 points) :} argument « pour » clair (1 pt) ; argument « contre » ou nuance (1 pt) ; conclusion personnelle justifiée (1 pt) ; cohérence, orthographe et longueur respectée (1 pt).
\end{corrige}

\begin{attention}[Point de vigilance]
La conclusion doit être une \emph{position argumentée}, pas une simple répétition. Une réponse qui ne donne qu'un seul point de vue ne peut pas obtenir la totalité des points.
\end{attention}

\begin{corrige}[Exercice 9 — Sujet type BEPC : L'IA dans l'éducation au Cameroun]
\textbf{1. Deux avantages pédagogiques :}
\begin{itemize}
    \item \textbf{Aide individualisée} : l'IA peut corriger un programme en C, expliquer une erreur de syntaxe ou proposer une structure de dissertation, comme un tuteur disponible pour chaque élève, même dans les classes surchargées.
    \item \textbf{Entraînement illimité} : l'élève peut générer des exercices d'algorithmique supplémentaires et s'auto-évaluer à son rythme, ce qui renforce la préparation au BEPC.
\end{itemize}

\textbf{2. Deux risques éthiques ou techniques :}
\begin{itemize}
    \item \textbf{Risque éthique} : le copier-coller sans compréhension (plagiat assisté) détruit l'apprentissage et fausse l'évaluation ; de plus, la saisie des données sensibles des élèves (notes, identités) dans une IA publique menace leur confidentialité.
    \item \textbf{Risque technique} : l'instabilité de la connexion Internet dans certains établissements rend le service inégal entre les lycées (fracture numérique), et les hallucinations de l'IA peuvent transmettre des erreurs aux élèves.
\end{itemize}

\textbf{3. Trois règles pour la Charte Numérique de l'Élève :}
\begin{enumerate}
    \item Je cite toujours l'IA lorsque je m'en sers (transparence) et je reformule avec mes propres mots.
    \item Je vérifie toute information fournie par l'IA dans mon manuel ou auprès de mon professeur avant de l'utiliser.
    \item Je ne saisis jamais de données personnelles (noms, notes, mots de passe) dans une IA en ligne.
\end{enumerate}

\textbf{4. Pourquoi former les enseignants à l'ingénierie de prompt :} la qualité de la réponse d'une IA dépend directement de la qualité du prompt. Un enseignant formé sait construire des consignes précises (Rôle--Tâche--Contexte--Format) pour obtenir des explications adaptées au niveau 3ème, détecter les réponses fausses ou hors programme, et guider les élèves vers un usage pédagogique plutôt qu'un simple copier-coller. Sans cette formation, l'outil serait mal utilisé et le déploiement inefficace, voire nuisible.

\textbf{5. Exemple de prompt efficace pour l'algorithmique :}
\begin{quote}
\og Agis comme un professeur d'informatique de collège. Propose-moi trois exercices d'algorithmique sur la structure \emph{Si...Alors...Sinon}, niveau 3ème (variables, conditions simples), sans donner les solutions. Après chacune de mes réponses, corrige-moi en expliquant mes erreurs. \fg
\end{quote}
(Rôle : professeur d'informatique ; Tâche : proposer 3 exercices puis corriger ; Contexte : niveau 3ème, chapitre des structures conditionnelles ; Format : exercices sans solutions, correction expliquée.)

\textbf{Barème indicatif (sur 10 points) :} Q1 : 2 pts ; Q2 : 2 pts ; Q3 : 2 pts ; Q4 : 2 pts ; Q5 : 2 pts (prompt avec les 4 éléments identifiables).
\end{corrige}

\begin{corrige}[Exercice 10 — Analyse critique : Détecter une hallucination]
\textbf{1. Deux erreurs factuelles :}
\begin{itemize}
    \item Arthur Zang n'est \textbf{pas l'inventeur du smartphone} : il a inventé le \textbf{Cardiopad}, une tablette médicale permettant de réaliser des examens cardiaques à distance. Le téléphone portable a été inventé par \textbf{Martin Cooper} (Motorola) en 1973.
    \item Le \og ZangPhone \fg, « premier téléphone portable africain lancé en 2010 », est \textbf{totalement inventé} : cette marque n'existe pas. C'est une hallucination typique : l'IA fabrique des détails plausibles (nom, date) pour combler ses lacunes.
\end{itemize}

\textbf{2. Règle d'usage responsable :} l'élève doit \textbf{toujours vérifier l'information} fournie par une IA dans au moins une source fiable et indépendante (manuel scolaire, encyclopédie, site officiel) avant de l'utiliser dans un devoir. L'IA est un point de départ, jamais une source finale.

\textbf{3. Prompt amélioré :}
\begin{quote}
\og Agis comme un documentaliste rigoureux. Donne-moi une courte biographie d'Arthur Zang, inventeur camerounais, en précisant son invention réelle. Cite tes sources (sites ou ouvrages) pour chaque fait important. Si tu n'es pas certain d'une information, dis-le explicitement au lieu d'inventer. \fg
\end{quote}
\end{corrige}

\begin{attention}[Point de vigilance]
Même avec un prompt demandant des sources, une IA peut \emph{inventer des références} qui semblent crédibles. La vérification humaine reste obligatoire : c'est le point que le correcteur attend dans la question 2.
\end{attention}

\begin{corrige}[Exercice 11 — Éthique et protection des données : Charte IA au cybercafé scolaire]
\textbf{1. Trois types de données personnelles à ne jamais saisir :}
\begin{itemize}
    \item Les \textbf{identifiants et secrets} : mots de passe, code PIN du téléphone ou du Mobile Money.
    \item Les \textbf{données d'identification} : nom complet associé à l'adresse du domicile, numéro de téléphone, numéro de carte nationale d'identité ou de l'acte de naissance.
    \item Les \textbf{données scolaires ou familiales sensibles} : notes et bulletins nominatifs, situation familiale, informations sur des camarades (photos, problèmes personnels).
\end{itemize}

\textbf{2. Pourquoi éviter un compte unique partagé :} avec un compte commun, l'\textbf{historique} des conversations de chaque élève est visible par tous les autres (atteinte à la confidentialité) ; les réponses de l'IA sont \textbf{personnalisées} à partir de cet historique mélangé, ce qui dégrade la qualité de l'aide ; enfin, en cas d'abus (insultes, triche), il est \textbf{impossible d'identifier le responsable}, ce qui pose un problème de discipline et de responsabilité.

\textbf{3. Deux règles pour la Charte :}
\begin{itemize}
    \item \textbf{Règle 1 — Transparence :} « Tout travail réalisé avec l'aide de l'IA doit le mentionner clairement (ex. : "réponse générée avec l'aide d'une IA, vérifiée et reformulée par mes soins"). Recopier une réponse d'IA en la présentant comme son propre travail est une tricherie. »
    \item \textbf{Règle 2 — Esprit critique :} « Toute information fournie par l'IA doit être vérifiée dans une source fiable (manuel, professeur, site officiel) avant d'être utilisée : l'IA peut se tromper. »
\end{itemize}

\textbf{4. Deux avantages d'un hébergement local :}
\begin{itemize}
    \item \textbf{Les données ne quittent pas l'établissement} : les conversations des élèves (mineurs) restent sur le serveur de l'école et ne sont ni stockées ni exploitées par une entreprise étrangère, ce qui protège leur vie privée.
    \item \textbf{Contrôle total et autonomie} : l'école maîtrise l'accès, les historiques et les filtres de contenu, et le service fonctionne même sans connexion Internet, ce qui garantit l'égalité d'accès malgré les coupures réseau.
\end{itemize}

\textbf{Barème indicatif (sur 10 points) :} Q1 : 2 pts (trois types distincts) ; Q2 : 2 pts (au moins deux arguments parmi historique, personnalisation, responsabilité) ; Q3 : 3 pts (règles claires et rédigées) ; Q4 : 3 pts (deux avantages justifiés).
\end{corrige}

\newpage

\coursec{Fiche bilan}

\begin{popfigure}
\centering
\begin{tikzpicture}[
  mindmap,
  grow cyclic,
  every node/.style={concept, execute at begin node=\hskip0pt},
  concept color=popblue,
  level 1/.append style={level distance=4.5cm, sibling angle=360/5},
  level 2/.append style={level distance=3cm, sibling angle=360/6},
  % Branch colors
  branch1/.style={concept color=poporange, text=white},
  branch2/.style={concept color=popgreen, text=white},
  branch3/.style={concept color=poppurple, text=white},
  branch4/.style={concept color=poppink, text=white},
  branch5/.style={concept color=popgold, text=white},
  % Node styles
  root/.style={concept color=popblue, text=white, font=\bfseries\large, minimum size=2.8cm},
  child/.style={font=\bfseries\small, minimum size=2.2cm, text width=2.2cm, align=center},
  grandchild/.style={font=\footnotesize, minimum size=1.5cm, text width=1.8cm, align=center, fill=white, text=popdark, concept color=white, draw=popmuted, thick}
]
\node[root] {IA en 3\`eme}
  child[branch1] {node[child] {1. Généralités}
    child[grandchild] {node {Définition}}
    child[grandchild] {node {Historique}}
    child[grandchild] {node {Types (Faible/Forte)}}
  }
  child[branch2] {node[child] {2. Usages \& Importance}
    child[grandchild] {node {Quotidien (GPS, Reco)}}
    child[grandchild] {node {Éducation \& Santé}}
    child[grandchild] {node {Agriculture \& Commerce}}
  }
  child[branch3] {node[child] {3. Avantages / Inconvénients}
    child[grandchild] {node {Rapidité, Précision}}
    child[grandchild] {node {Dispo 24/7}}
    child[grandchild] {node {Biais, Emploi, Vie privée}}
  }
  child[branch4] {node[child] {4. Usage Responsable}
    child[grandchild] {node {Éthique \& Légal}}
    child[grandchild] {node {Pensée critique}}
    child[grandchild] {node {Protection données}}
  }
  child[branch5] {node[child] {5. Échanger (Prompting)}
    child[grandchild] {node {Rôle, Contexte, Tâche}}
    child[grandchild] {node {Format, Contraintes}}
    child[grandchild] {node {Itération \& Vérif.}}
  };
\end{tikzpicture}
\legende{Carte mentale : Découverte d'un environnement d'IA (Leçon 01, 3\`eme)}
\end{popfigure}

\begin{bilan}
\textbf{Définitions clés}
\begin{itemize}
  \item \cle{Intelligence Artificielle (IA)} : Ensemble de théories et de techniques visant à créer des machines capables de simuler l'intelligence humaine (apprentissage, raisonnement, perception).
  \item \cle{IA Faible (Étroite)} : Spécialisée dans une tâche précise (reconnaissance faciale, traduction, recommandation Netflix). C'est l'IA d'aujourd'hui.
  \item \cle{IA Forte (Générale)} : Hypothétique, capable de comprendre, apprendre et appliquer toute tâche intellectuelle humaine (niveau conscience).
  \item \cle{Apprentissage Automatique (Machine Learning)} : Sous-domaine où la machine apprend à partir de données sans être explicitement programmée pour chaque règle.
  \item \cle{Prompt} : Instruction textuelle (ou multimodale) donnée à une IA générative pour obtenir une réponse.
  \item \cle{Ingénierie de prompt (Prompt Engineering)} : Art de formuler des requêtes efficaces pour guider l'IA vers le résultat souhaité.
\end{itemize}

\textbf{Principes \& Règles à retenir}
\begin{center}
\begin{tabularx}{\linewidth}{|>{\bfseries}l|X|}
\hline
\rowcolor{popblueL}
\textcolor{popdark}{Domaine} & \textcolor{popdark}{Règle / Principe essentiel} \\
\hline
\textbf{Avantages} & Automatisation tâches répétitives, rapidité de traitement masse de données, disponibilité 24/7, aide à la décision (diagnostic médical, prévision météo). \\
\hline
\textbf{Inconvénients / Risques} & Biais algorithmiques (données d'entraînement), hallucinations (réponses inventées), dépendance technologique, impact emploi, empreinte écologique (énergie), vie privée. \\
\hline
\textbf{Usage Responsable} & \textbf{1. Vérifier} les sources (esprit critique). \textbf{2. Protéger} ses données perso (pas de mot de passe, CNI, notes). \textbf{3. Citer} l'IA si usage scolaire/pro. \textbf{4. Respecter} le droit d'auteur \& RGPD/Cameroun (Loi 2024/011). \textbf{5. Humain dans la boucle} : l'IA propose, l'humain décide. \\
\hline
\textbf{Structure Prompt Efficace (R.C.T.F.C)} & \textbf{R}ôle (Tu es prof de maths), \textbf{C}ontexte (Classe de 3ème, programme Cameroun), \textbf{T}âche (Explique le PGCD), \textbf{F}ormat (Liste à puces + 2 exos), \textbf{C}ontraintes (Langage simple, pas de LaTeX). \\
\hline
\end{tabularx}
\end{center}

\textbf{Méthodes express}
\begin{itemize}
  \item \textbf{Évaluer une sortie IA :} Lire $\rightarrow$ Croiser (source fiable) $\rightarrow$ Corriger (hallucination ?) $\rightarrow$ Valider.
  \item \textbf{Améliorer un prompt :} Test $\rightarrow$ Analyse erreur $\rightarrow$ Ajout contrainte/contexte $\rightarrow$ Re-test (boucle itérative).
  \item \textbf{Protéger sa vie privée :} Anonymiser les données avant copier-coller $\rightarrow$ Utiliser mode "incognito" / "chat temporaire" si dispo $\rightarrow$ Effacer l'historique régulièrement.
\end{itemize}
\end{bilan}

\begin{aretenir}[Checklist avant le BEPC]
\begin{itemize}
  \item $\square$ Je sais définir l'IA et distinguer IA faible / IA forte.
  \item $\square$ Je cite 3 usages concrets de l'IA au Cameroun (ex: diagnostic culture, traduction langue locale, banque mobile).
  \item $\square$ Je liste 3 avantages et 3 inconvénients/risques majeurs de l'IA.
  \item $\square$ Je connais les 5 piliers d'une utilisation responsable (Vérifier, Protéger, Citer, Respecter, Décider).
  \item $\square$ Je sais structurer un prompt avec la méthode R.C.T.F.C (Rôle, Contexte, Tâche, Format, Contraintes).
  \item $\square$ Je sais détecter une hallucination et expliquer comment la vérifier.
  \item $\square$ Je ne communique jamais de mot de passe, numéro CNI, ou note secrète à une IA.
  \item $\square$ Je sais expliquer le biais algorithmique avec un exemple simple (ex: reconnaissance visage moins précis sur peaux foncées).
  \item $\square$ Je connais la loi camerounaise de référence sur la protection des données (Loi 2024/011).
  \item $\square$ Je suis capable d'améliorer un prompt mauvais en ajoutant du contexte et des contraintes.
  \item $\square$ Je comprends que l'IA est un outil d'aide, pas un remplacement de ma réflexion.
  \item $\square$ Je sais citer l'IA dans un devoir : « Texte généré avec [Nom IA], date, prompt utilisé ».
\end{itemize}
\end{aretenir}

\findecours

\end{document}
